\chapter{Indirect modification in Japanese} \label{indirect japanese}

This chapter discusses the indirect domain and focus will be switched to indirect, that is predicative, modification. Although this kind of modification is often assumed to be the only one available in Japanese \citep{sproat:1991:the, baker:2003:verbal, laenzlinger:2011:elements}, no systematic analysis has been presented for indirect modifiers in this language, either. This concerns their fine-grained structure and the integration of this structure into the DP, as well as their size and their position in the DP. This chapter thus aims at closing those research gaps and at broadening our understanding of this part of the DP, and indirect modifiers in Japanese and in general. After pointing out the characteristics of indirect modifiers in Japanese, I will discuss the issues mentioned above in turn. %\is{indirect modification}

\section{Introduction}

Recall that the indirect domain is situated inside the DP-internal domain, higher than the direct domain. It exhibits several functional projections for indirect modifiers as well as a projection that hosts numerals. As in the direct domain, these projections are not dedicated to a specific type of indirect modifier, in this case to different sizes. The structure of the indirect domain inside the DP-internal domain is circled in the structure of the DP-internal domain in \figref{fig:Structure of the indirect domain in Japanese2}.

\begin{figure}

 \begin{tikzpicture}[baseline=0pt, remember picture]
\tikzset{level distance=30pt, sibling distance=-5pt}
 \Tree [.DP\textsubscript{internal} {} [.D′\textsubscript{internal} [. \node (indirect1) {FP\textsubscript{indirect-n}}; CP [.\node (indirect0) {F′\textsubscript{indirect-n}}; [.FP\textsubscript{Num} NumP [.F′\textsubscript{Num} [.FP\textsubscript{indirect-n+1} \node (indirect3) {CP}; [. \node (indirect2) {F′\textsubscript{indirect-n+1}}; [.dP [.FP\textsubscript{direct-n} XP [.F′\textsubscript{direct-n} [.FP\textsubscript{direct-n+1} XP [.F′\textsubscript{direct-n+1} [.$\sqrt{}$P XP [.$\sqrt{}$′ \edge[roof]; N ] ] F ] ] F ] ] d ] F ] ] F ] ] F ] ] D ] ]
 \node[draw,circle,fit= (indirect0) (indirect3) ] [color=blue] {};
\end{tikzpicture}

 \caption{Structure of the indirect domain in Japanese}
 \label{fig:Structure of the indirect domain in Japanese2}
\end{figure}

\subsection{Goals and overview}

This chapter deals with the following questions relevant for indirect modifiers in Japanese:

\ex. \a. What are the complex morphological elements co-occurring with indirect modifiers, especially adjectives such as \textit{taka-katta} `was expensive', in Japanese?
\b. What is the size and merging position of indirect modifiers?
\c. What influences the position of indirect modifiers?

These questions will be addressed in turn and yield the following answers:

\ex. \a. The complex morphological elements such as \textsc{k-at-ta} are made-up of three elements: two copulas, and a tense morpheme. This analysis can be adapted to seemingly ``simpler" elements such as \textsc{da}.
\b. There is no difference in Japanese between full and reduced relative clauses;  indirect modifiers come only in one size. All relative clauses in Japanese are CPs and derived via operator movement \citep{miyamoto:2017:RC}.\footnote{In fact, nothing hinges on whether they have to be CPs or IPs (although CP-status grants a landing position for the operator). The important point is that relative clauses in Japanese only come in one size.} \c. The only factor which has an effect on the position of modifiers inside the indirect domain is syntactic and prosodic complexity. More complex modifiers move to a position above numerals.

\subsection{Characteristics of Japanese indirect modifiers} \label{nishiyama analysis}

Recall that relative clauses are the prime example of indirect modifiers (\citealt{cinque:2010:adj, cinque:2020:relcl}, among others). In Section \ref{dP}, it was mentioned that languages such as \is{English} English overtly display two types of relative clauses differing in size, structure and position, namely \textit{full} and \textit{reduced} relative clauses, but that this difference is generally hard to determine to non-existent in prenominal-RC languages. Such languages usually do not show two different types of relative clauses overtly, but instead exhibit other types of noun-modifying clauses that do not have relative clauses as direct counterparts in English and other well-known European languages \citep{keenan:1977:noun, lehmann:1984:der, vries:2002:syntax, wu:2011:syntax}.\footnote{For a coherent and contemporary typological overview of prenominal relative clauses see \citet{wu:2011:syntax}, who also points out that prenominal relative clauses ``have essentially only been studied in a few well-known and often-quoted languages such as Basque, Japanese, Mandarin, and Turkish" \citep[569]{wu:2011:syntax}.} \is{reduced relative clause} \is{Basque} \is{Mandarin} \is{Turkish}

Japanese fits the description of other prenominal-RC languages. This language only exhibits one morphosyntactic type of relative clause, but it offers several constructions that do not correspond to relative clauses in European languages \citep{comrie:1998:rethinking, matsumoto:2017:japanese}. Furthermore, it exhibits the following characteristics which further complicate the analysis.


\ex. \a. Lexemes exhibit (with the exceptions of nouns and \textit{na}-adjectives in non-past) no overt morphological difference between their role as modifier and in predicative position.
\b. Modifiers can exhibit a fairly complex morphological make-up in attributive position, again mirroring their predicative forms.
\c. Japanese exhibits a rather free word order within the DP, and this extends to indirect modifiers.
\d. Relative clauses can appear above demonstratives and can in principle be restrictive and non-restrictive in this position.

\largerpage
The observation that, with a few exceptions, the morphology of lexemes in attributive position is identical to the one displayed in predicative position further differentiates Japanese from other East Asian prenominal-RC languages. For example, looking at \is{Mandarin} Mandarin (and many other Sino-Tibetan languages) or \is{Korean} Korean, relative clauses exhibit certain elements which connect them to the noun (\textsc{de} in Mandarin, \textsc{(nu-)n} in Korean),  and which are absent in non-embedded root clauses. Japanese lacks such an element. Instead, the surface form of an indirect modifier matches that of its predicative counterpart. This is, once more, exemplified below in \ref{verbal rc, chapter 5} for a verbal relative clause.\footnote{It might be tempting to deduce the size of Japanese relative clauses from the fact alone that modifiers appear in a \textit{morphologically} \textit{finite} form. Already \citet[156]{keenan:1985:relcl} mentioned ``We may note that the prenominal RCs in Japanese do not put Vrel in a non-finite or specifically relative form, but the Japanese case appears to be the exception among prenominal RCs here". However, the notion of finiteness is not very well understood even in European languages and lacks a uniform definition. See the discussion in \citet{hudson:1973:tense, stanton:2011:rrc, belikova:2008:syntactically} and \citet{belk:2017:attributes} with regard to \is{reduced relative clause}reduced relative clauses in English, Dutch and Russian among others and \citet{nikolaeva:2007:intro} for an overview; see, for instance, \citet{wurmbrand:2005:infinitives, wurmbrand:2014:tense} and \citet{landau:2004:scale, landau:2013:control} for various phenomena in English.

The biggest obstacle is that finiteness can be approached from different perspectives. Overt tense is interpreted as morphological finiteness, especially in the Latin-Greek tradition. The classical syntactic understanding of finiteness is the assignment of nominative case \citep{chomsky:1981:lectures}, but also the morpho-syntactic marking of tense, person, politeness and other categories (see \citealt{nikolaeva:2007:intro} and articles therein). Its association with the C-domain has been put into question in the literature as well, see \citet{adger:2007:finiteness, bisang:2007:finiteness} and especially \citet{wurmbrand:2020:finiteness}, who argue that it can be encoded at different levels of the clause, ranging from C to T even to \textit{v}. From a semantic point of view, finiteness is related to the distribution of overt tense and the indication of the ``existence of a predication at the time of the utterance" \citep[201]{platzack:1995:loss}. As shown for example in \citet{ogihara:1994:adverbs} and \citet{miyagawa:2012:case}, Japanese relative clauses can, but do not need to, be dependent on the higher clause.

To summarize, finiteness is not helpful in determining the status and syntactic size of relative clauses and will be sidelined here as a criterion. Nevertheless, note that I used the notion of finiteness, in the sense that tense is projected, to argue for the possible non-occurrence of realizations of \is{MOD@\textsc{mod}}\textsc{mod}.} \is{finiteness} \is{English} \is{Dutch} \is{Russian}

\ex. \label{verbal rc, chapter 5} \ag. Kare-ga hon-wo kat-ta.\\
he-\textsc{nom} book-\textsc{acc} buy-\textsc{pst}\\
`He bought a book.’
\bg. kare-ga kat-ta hon\\
he-\textsc{nom} buy-\textsc{pst} book\\
`the book (which) he bought'

This extends to all non-verbal modifiers inflected for past: \textit{na}-adjectives and \textit{no}-modifiers can co-occur with the canonical copula forms \textsc{de-ar-u} and \textsc{d(e)-at-ta} both in attributive and predicative position, and \textit{i}-adjectives can co-occur with \textsc{k-at-ta} in order to show that they are inflected for past, as shown below.

\ex. \ag. Ano doresu-wa taka-k-at-ta.\\
that dress-\textsc{top} expensive-\textsc{pred-cop-pst}\\
`That dress was expensive.'
\bg. ano taka-k-at-ta doresu\\
that expensive-\textsc{pred-cop-pst} dress\\
`that dress which was expensive' \null \hfill \textit{i}-adjective, past

I will start with a discussion of these elements before discussing the size and position of indirect modifiers.

\section{The elements PRED, COP and T} \label{copulaelements}

In this section, I present a precise analysis of the complex morphological elements Japanese modifiers can co-occur with. Following Nishiyama (\citealt{nishiyama:1999:adj}, and see also \citealt{miyama:2011:da, watanabe:2013:non, yamada:2023:looking}), I will step by step outline that they are composed of a predicative copula (\textsc{pred}) in the sense of \citet{bowers:1993:predication}, a dummy copula and a tense marker, starting with the final element, and translate the data into the cartographic structure adopted here. This analysis will then be applied to the shorter version of the copula, \textsc{i} and \textsc{da}, and I will explain why an adaptation to the monomoraic elements occurring in attributive position is not justified.

\subsection{Nishiyama's analysis}

\subsubsection{Tense marker}

The final element \textsc{ta} contained in \textsc{k-at-ta} and \textsc{d(e)-at-ta} and also occurring with verbs straightforwardly receives an analysis as a tense marker for past.\footnote{\textsc{ta} is likely derived from the premodern auxiliary verb -\textit{tari}, which was added to the nominalization form of inflecting categories and was itself inflecting.} \textsc{ta} is subject to different phonological processes, namely gemination, and also voicing, yielding the allomorph /da/. The first generalization to be drawn is thus as follows.

\ex. \textsc{ta} = \textsc{I} (past)

The analogous element for non-past is \textsc{u}.\footnote{When the vowel /u/ is added to the stem of vowel-stem verbs such as \textit{tabe} `to eat’, an epenthetic alveolar /r/ is inserted. Note that there are also consonant-stem verbs whose stem ends in /r/ and which share their paradigm with the \is{dummy copula}dummy copula (see next section). In this case, the alveolar consonant is part of the stem. Before the past ending and the suffix -\textsc{te}, /r/ is deleted and /t/ is geminated. These groups of verbs should thus not be mixed up. They can be further differentiated based on their inflectional paradigms. For discussion see \citet[364–365]{ito:2015:word}. In any case, the morpheme signaling non-past remains the same.}

\ex. \textsc{u} = \textsc{I} (non-past)

\subsubsection{Dummy copula}

When added to adjectives and \textit{no}-modifiers, \textsc{ta} is obligatorily preceded by the element \textsc{ar}, which Nishiyama analyzes as the \is{dummy copula}dummy copula. In this sense, the copula is \textit{semantically vacuous}, a ``dummy element whose function is merely to support tense" \citep[187]{nishiyama:1999:adj}. It thus performs a function of \textit{be}-support, similar to English \textit{do}-support. \is{be-support}

This view goes back to \citet{bach:1967:have} and \citet{doron:1983:verbless}, among others, and was formalized in \citet{rapoport:1987:copular}. In \is{English} English, the copula \textit{be}, although it can serve the semantically important function of identification (\textit{That woman is Tali.}; cf. \citealt[150]{rapoport:1987:copular}), is typically devoid of meaning, yet necessary in predicative and equative sentences involving non-verbal elements (examples based on \citealt[156]{rapoport:1987:copular}).

\ex. \label{nut} \a. Xeli *(is) a nut.
\b. Tali *(is) that woman over there.

As a contrast, it is not necessary in small clause predicational environments involving a lexical verb.

\ex. I consider Xeli a nut. \label{consider} \null\hfill \citep[152]{rapoport:1987:copular}

\citet{rapoport:1987:copular} first showed that there is no semantic requirement for the copula, that it is semantically empty, hence must be a syntactic element that satisfies a syntactic requirement. This requirement is that every sentence has to contain a verb, since every sentence needs to express tense \citep[157]{rapoport:1987:copular}. This can be achieved by a lexical verb \ref{consider}, or, alternatively, by the copula \ref{nut}. In the absence of a verb, the copula supports the tense (and agreement) features of \textsc{I} (see \citealt[157]{rapoport:1987:copular}).\footnote{It is not clear whether \textit{consider} is indeed a lexical verb in this analysis. See \citet{bowers:1993:predication} for an analysis of \textit{consider} and similar verbs in English.}

In many languages, the \is{dummy copula}semantically vacuous copula can be omitted in the default tense non-past/present, but is needed in other contexts, for example to express past where it bears tense inflection. This underlines its supporting character for the I node and lack of semantic contribution. Observe the following scenario in \is{Russian}Russian (cf. \citealt{babby:2010:syntactic}; cf. \citealt{doron:1983:verbless, rapoport:1987:copular} for \is{Hebrew}Hebrew).

\ex. \label{pagoda} \ag. Pagoda xorošaja. (Russian)\\
weather.\textsc{nom.sg.f} good.\textsc{nom.sg.f}\\
`The weather is good.' \null \hfill present
\bg. Pagoda *(byla) xorošaja.\\
weather.\textsc{nom.sg.f} was.\textsc{sg.f} good.\textsc{nom.sg.f}\\
`The weather was good.' \null \hfill past

The semantically vacuous copula in Japanese is identified in \citet{nishiyama:1999:adj} and \citet{urushibara:1993:syntactic} as -\textsc{ar}- and given the syntactic label \textit{V}. To avoid confusion, I will use the label \textsc{cop}.\footnote{\citet{watanabe:2013:non} assigns the dummy copula the light verb status \textit{v}. Recall that before the past tense morpheme \textsc{ta}, obligatory gemination happens.}

\ex. \textsc{ar} = \textsc{cop}

\textsc{ar} can also be used as an independent verb, which expresses existence, identical to English \textit{be}.

\exg. Kuruma-wa niwa-ni ar-u.\\
car-\textsc{top} garden-\textsc{loc} \textsc{cop-prs}\\
`The car is in the garden.’

It also appears in other scenarios \textit{be} appears in, including equation \ref{equation} and identification \ref{identification}. In both cases, another element, namely \textsc{de} (the \is{predicative copula}predicative copula, see below) is required.

\exg. Hanako-wa sensei-de-ar-u.\\
Hanako-\textsc{top} teacher-\textsc{pred-cop-prs}\\
`Hanako is a teacher.' \label{equation}

\exg. Ano ko-wa Hanako-de-ar-u.\\
that girl-\textsc{top} Hanako-\textsc{pred-cop-prs}\\
`That girl over there is Hanako.' \label{identification}

It needs to be made clear that the sole, but very strong, piece of evidence that \textsc{ar} is a \is{dummy copula}(dummy) copula is its requirement to support tense, in other words to host the tense node. Non-verbal elements can never express overt tense, with an overt non-past or past morpheme, without the presence of \textsc{ar}, as given below.

\exg. *Ano doresu-wa shikku-(de)-u / (de)-ta.\\
that dress-\textsc{top} chic-\textsc{pred-prs} / \textsc{pred-pst}\\
(`That dress is/was chic.’)

This means the dummy copula is a different element from \textsc{de}, which is not sufficient to support tense and will be reviewed below.

\subsubsection{Predicative copula} \label{chapter pred}

However, a tense marker and the dummy copula are not enough to license an adjective or a \textit{no}-modifier in \is{predicative copula}predicative position either. Additionally, the element \textsc{de}, in the case of \textit{na}-adjectives and \textit{no}-modifiers, or \textsc{ku}, in the case of \textit{i}-adjectives, is required. Following \citet{nishiyama:1999:adj}, I will analyze this element as another, more meaningful kind of copula, namely as the syntactic head for predication \textsc{pred} (\citealt{bowers:1993:predication}, there \textsc{pr}; cf. also \citealt{svenonius:1996:pred, bowers:2002:pred, dikken:2006:relators, heycock:2013:syntax}). This element is the real bearer of predication.\footnote{A similar view can also be found in the work of one of the first Western linguists specialized in Japanese: Bernard Bloch. Bloch claimed in the 1940s that the Japanese copula is an inflectional element and that every sentence needs such an inflectional element as its nucleus, making the copula an essential ingredient of the sentence (\citealt[35]{bloch:1970:bloch}, see \citealt[187–188]{nishiyama:1999:adj}).}

Evidence for this analysis and for the differences between the two copulas are visible in the following data points. Precisely, whenever a predication relation between a non-verbal \textit{na}-adjective or \textit{no}-modifier and a verb is established, the predicative copula obligatorily appears, but the dummy copula is absent. This is visible in \ref{overtpred}, where \textit{no}-modifiers are used as depictives of a verb. See especially \citet{koizumi:1995:secondary}, to whom most of these data points and observations can be attributed.

\ex. \label{overtpred} \ag. Jon-ga sakana-wo hadaka-*(de) tabe-ta.\\
John-\textsc{nom} fish-\textsc{acc} naked-\textsc{pred} eat-\textsc{pst}\\
`John ate the fish naked.' \null\hfill \citep[188]{nishiyama:1999:adj}
\bg. Jon-ga katsuo-wo nama-*(de) tabe-ta.\\
John-\textsc{nom} bonito-\textsc{acc} raw-\textsc{pred} eat-\textsc{pst}\\
`John ate the bonito raw.' \null \hfill \citep[27]{koizumi:1995:secondary}

In these examples, the predication relation is established via \textsc{pred}, which is spelled-out as /de/. While \is{predicative copula}predication needs to be established, the tense node already has a host, namely, in both sentences the verb \textit{tabe} `to eat’. Therefore, the dummy copula is not needed.

This means the two copulas differ in the \is{dummy copula}following ways:

\ex. \a. The dummy copula is semantically meaningless. It is simply there to support tense. The predicative copula is the bearer of predication.
\b. The dummy copula cannot occur independently and is not required when a verb is present. The predicative copula can independently co-occur with modifiers in tenseless environments, in which a verb establishes tense.

Following \citet{bowers:1993:predication} and its application to Japanese in \citet{nishiyama:1999:adj}, I define \is{predicative copula}\textsc{pred} as a functional category with a syntactic as well as a semantic dimension. Semantically, it fulfills the function \textit{Pred(ication)} in the sense of \citet{chierchia:1988:semantics}. Syntactically, it is a functional head which projects the phrase PredP and selects the maximal projection of any lexical category (AP, NP, PP, etc., given as XP below) as its external argument in its specifier position. From there, a predication relation with the predicate argument in its complement position is established \citep[595–596]{bowers:1993:predication}. This means \textsc{pred} is the predicative counterpart of \is{MOD@\textsc{mod}}\textsc{mod} (\citealt{rubin:1994:mod}; Section \ref{internalstructure}, see especially Section \ref{semantics mod}). The general structure is given in \figref{structure PredP}.

\begin{figure}
\centering
\begin{tikzpicture}[baseline=-0pt, remember picture]
\tikzset{level distance=25pt, sibling distance=5pt}
\Tree [.PredP {(subject) NP} [.Pred′ Pred {XP (predicate)} ]]]]
\end{tikzpicture}
\caption{Structure of PredP}
\label{structure PredP}
\end{figure}

There are many syntactic points of justification and empirical arguments why the category \textsc{pred} is necessary, and why, besides its semantic dimension, it should be kept apart from VP. These include the following (see especially \citealt[598–599]{bowers:1993:predication}):\footnote{Similar to the proposal of the syntactic head \textsc{mod} (see Section \ref{internalstructure}), the adoption of a universal head \textsc{pred} is not without controversy. For discussion see \citet{bowers:1993:predication, bowers:2002:pred, svenonius:1996:pred, dikken:2006:relators, heycock:2013:syntax} as well as \citet{matushansky:2019:against} and \citet{bruening:2020:category}, who have spoken decisively against this hypothesis. Besides theory-internal problems, the authors argue that the relevant head is redundant, as many problems it was meant to solve in the original theory can now be solved via other mechanisms, and this is indeed the case. For instance, \textit{unlike} (= dissimilar) categories can also be \is{coordination}coordinated via a Boolian phrase.

\par

Furthermore, \citet{matushansky:2019:against} tries to show that the alleged manifestations of \textsc{pred} in several languages do not hold, because each language shows exceptions, where a copular element does not appear even though one would expect it to. In Edo, the alleged \textsc{pred} elements do not appear in resultatives, in Welsh they do not appear in PPs, and in the Berber language East Riffian they are restricted to adjectives of Arabic origin \citep[§ 4]{matushansky:2019:against}.

\par

While this is problematic for the universality of the \textsc{pred}-hypothesis, Matushansky's analysis does not prove that no \textsc{pred}-like elements exist. What holds for \textsc{mod} also holds for \textsc{pred}: The relevant element is attested in Japanese. I argue that it is present in this language and, sticking with the cartographic idea of language-uniformity, cross-linguistically.}

\ex. Having a functional head \textsc{pred} allows for the coordination of unlike lexical categories.

\citet{bowers:1993:predication} explained cases such as \textit{I regard him} [[\textsubscript{\textsc{dp}} \textit{a fool}] \textit{and} [\textsubscript{\textsc{ap}} \textit{unable to do his job}]] without having to resort to non-binary structures. Both \is{coordination}conjuncts are being mediated by the predication head. \is{predicative copula}

\ex. It enables movement of lexical categories, with the possibility that the subject can remain in Spec, PredP.

In a sentence such as \textit{How stupid do} [\textsubscript{\textsc{predp}} \textit{I} [\textsubscript{\textsc{pred}} \textit{consider him}]]\textit{?}, the subject remains in situ, while phrasal material moves above it. This was a significant contribution at the time (cf. \citealt[68–69]{matushansky:2019:against} for discussion).

\ex. PredP provides a syntactic position for the subject in expletive-construc\-tion and various other small clauses, fulfilling the EPP even in the absence of thematic relations.

In tenseless small clauses such as \textit{I consider} [\textit{him stupid}], no obvious subject position is available. This subject position is, however, offered by \textsc{pred} in addition to licensing predication.

\ex. Subject and object position are kept apart, the subject being in Spec, PredP, the object in Spec, VP.

This relates to the previous point in the sense that non-thematic subjects, for example those in \textit{tough} constructions, also receive a subject position, although this position is not related to thematic role assignment, i.e. that of the agent.

\ex. Many languages have overt realizations of \textsc{pred}.

Probably the strongest piece of evidence for \textsc{pred} is offered by its lexical and syntactic manifestations in various languages. This includes \textit{as} and \textit{for}, which appear in small clauses with matrix verbs such as \textit{regard} in the absence of a copula in \is{English}English \ref{predenglish} \citep[596]{bowers:1993:predication}. Similar analyses have been proposed for the cross-linguistic counterparts of these elements, for example \is{Russian}Russian \textit{kak} `as' \ref{predrussian} (or \textit{za} `as' or \textit{v} `in’; cf. \citet{bailyn:2002:overt}.\footnote{Further possible manifestations of \textsc{pred} include instrumental case in \is{Russian}Russian \citep{bailyn:2001:syntax}, and copular elements in a variety of African languages including \is{Edo}Edo and \is{Chichewa}Chichewa \citep{baker:2003:cat}, as well as \is{Welsh}Welsh \citep{bowers:2002:pred} and \is{Modern Greek}Greek \citep{campos:2004:ds}.}

\exg. I regard [\textsubscript{SC}John as crazy].\\
\textsc{sub} \textsc{v} \textsc{obj} \textsc{pred} \textsc{a} \\ \null \hfill \citep[596]{bowers:1993:predication} \label{predenglish}

\exg. On vygladit [kak durak]. (Russian)\\
he look \textsc{pred} fool\\
`He looks like a fool.’ \null\ \hfill \citep[39]{bailyn:2002:overt} \label{predrussian}

This also includes the overt manifestations of \is{predicative copula}\textsc{pred} in Japanese. Besides depictives, presented in \ref{overtpred}, another instance where the predicative copula is manifested in the absence of the \is{dummy copula}dummy copula is coordination structures. As argued in Section \ref{coordination copula}, \textsc{de} must either autonomously or together with the sentential conjunction marker \textit{katsu} establish coordination of two nominals. Such examples were analyzed as clausal coordination in which two instances of the copula must appear, one at the end of the sentence as the main clause predicate, one after the first conjunct. An example is repeated below (cf. \citealt[364]{hiraiwa:2012:mechanism}, from which examples are adapted). \is{coordination}

\exg. Kono hito-wa Ken-no sensei *(katsu) / *(de) Naomi-no sensei-da.\\
this person-\textsc{top} Ken-\textsc{mod} teacher and / \textsc{pred} Naomi-\textsc{mod} teacher-\textsc{cop}\\
`This person is Ken's teacher and Naomi's teacher.’

Recall also that direct modifiers, such as \textit{Chūgoku-no} `Chinese’, on the other hand, can never appear in such coordination structures, not even as the second modifier in the DP. As they can never co-occur with the copula in predicative position, it is expected that they can therefore never co-occur with \textsc{de} either, if this element has copular status as well.

\ex. \ag. *Chūgoku-de kirei-na kabin\\
China-\textsc{pred} pretty-\textsc{mod} vase\\
\bg. *kirei-de Chūgoku-no 	kabin\\
beautiful-\textsc{pred} China-\textsc{mod} vase\\
(`a Chinese and beautiful vase’)


Finally, there are some historical pieces of evidence for the copular status of \textsc{de}. This form in fact goes back to an amalgamation of the premodern copula \textsc{ni} -- the short form of the premodern copula inflected for nominalization form \textit{nari} -- and the \is{coordination}coordinative suffix \textsc{te}. These two elements began to phonologically assimilate in Classical Japanese and have completed merging into one element, \textsc{de}, in Modern Japanese \citep{frellesvig:2010:history, miyama:2011:da}. \is{predicative copula} \is{Classical Japanese}

This means that /de/ in \textsc{de-ar-u} and in \textsc{de-at-ta} is the predicative copula. The representation of all three elements thus looks as follows.

\ex. \ag. Yoru-ga shizuka-de-ar-u.\\
night-\textsc{nom} quiet-\textsc{pred-cop-prs}\\
`The night is quiet.'
\bg. Yoru-ga shizuka-de-at-ta.\\
night-\textsc{nom} quiet-\textsc{pred-cop-pst}\\
`The night was quiet.' \null \hfill \citep[189, 194]{nishiyama:1999:adj}

What about /d/ in \textsc{d-at-ta}, in which the vowel /e/ is not present? The vowel /e/ is deleted here \citep[203]{nishiyama:1999:adj}, probably to avoid the clashing of two vowels, a phenomenon not uncommon in Japanese, both diachronically and synchronically.

Coming now to \textit{i}-adjectives, these exhibit, in certain contexts, the same complex tripartite structure of morphemes parallel to the complex form \textsc{de-ar-u}. This complex form is \textsc{ku-ar-u}. It seems straightforward – and this is the analysis adopted in \citet{nishiyama:1999:adj} – that \textsc{ku} corresponds to \textsc{de}. All things being equal, these two complex forms only differ in one aspect, namely the realization of the predicative copula. Japanese thus provides evidence for a second possible realization of \textsc{pred}.

One thing to consider is that \textsc{ku-ar-u} is restricted in its distribution and rather marked in canonical root clauses. 

\exg. ??Ano doresu-wa taka-ku-ar-u.\\
that dress-\textsc{top} expensive-\textsc{pred-cop-prs}\\
`That dress is expensive.'

\newpage
A prominent context in which this complex form is allowed, however, is emphatic expressions containing the \is{focus}focus-sensitive particles -\textit{mo} `even', -\textit{sae} `even' or -\textit{bakari} `only’. Interestingly, these intervene between \textsc{ku} and \textsc{ar-u}.\footnote{This was noted as early as \citet{miyata:1948:nihongo}, see \citet[185]{nishiyama:1999:adj} and \citet{kubo:1992:japanese}. It is unclear where exactly those focus markers emerge and are situated.}

\exg. Ano doresu-wa taka-ku mo ar-u.\\
that dress-\textsc{top} expensive-\textsc{pred} \textsc{emp} \textsc{cop-prs}\\
`That dress is even expensive.'

There are further instances where \textit{i}-adjectives must co-occur with these morphological complex forms: Before the modals \textit{beki} `should’ and \textit{mai} `will not be’ and in subordinate clauses with the complementizer \textit{koto} (\citealt[199]{nishiyama:1999:adj} and \citealt[51]{miyama:2011:da}). In these cases, these elements form a direct string. \is{complementizer}

\ex. \ag. Ano doresu-wa taka-ku-ar-u beki-da / mai.\\
that dress-\textsc{top} expensive-\textsc{pred-cop-prs} should-\textsc{cop} / will.not.be\\
`That dress should be / will not be expensive.'
\bg. Ano doresu-wa taka-ku-ar-u koto-wo wasure-ta.\\
that dress-\textsc{top} expensive-\textsc{pred-cop-prs} \textsc{comp}-\textsc{acc} forget-\textsc{pst}\\
`I forgot that that dress is expensive.'

In the same constructions, \textit{na}-adjectives (and \textit{no}-modifiers) must co-occur with the corresponding form of the canonical copula \textsc{de-ar-u}  (or past \textsc{de-at-ta}). In fact, \textsc{da} is not possible \citep[186]{nishiyama:1999:adj}. As shown in \ref{shikkudemo}, their use in emphatic contexts further highlights the parallelism between \textsc{k(u)} and \textsc{de}. \is{predicative copula}

\exg. Ano doresu-wa shikku-de-ar-u beki-da / *shikku-da beki-da.\\
that dress-\textsc{top} chic-\textsc{pred-cop-prs} should-\textsc{cop} / chic-\textsc{cop} should-\textsc{cop}\\
`That dress should be chic.'

\exg. Ano doresu-wa shikku-de mo ar-u.\\
that dress-\textsc{top} chic-\textsc{pred} \textsc{emp} \textsc{cop-prs}\\
`That dress is even chic.' \label{shikkudemo}

\largerpage
Moreover, \textsc{k(u)} also appears in other contexts in which \textsc{de} appears, for instance in coordinative structures involving \textit{i}-adjectives. Recall from Section \ref{coordination copula} that it is optionally followed by \textsc{te}, yielding the more canonical option. Now, we can safely identify \textsc{ku} in these contexts as the predicative copula as well. \is{coordination} \is{coordinative copula}

\exg. aka-ku-(te) utsukushi-i doresu\\
red-\textsc{pred-ger} beautiful-\textsc{mod} dress\\
`a red and beautiful dress'

The \is{predicative copula}predicative copula is used in tenseless constructions as well, for example in resultative, inchoative and causative constructions. As expected, it is not followed by the dummy copula in these contexts, because a verb is present to establish tense. This can be put forth as a further piece of evidence that Japanese has a requirement for \textsc{pred} (different from, for example, \is{English} English, visible in the translations); for similar examples and discussion also see \citet[8–10]{yamada:2023:looking}. Finally, the adverbial inflectional form of \textit{i}-adjectives is also /ku/ and receives the same analysis as \textsc{pred} \citep{nishiyama:2005:baker}, but there is some controversy \citep{namai:2002:word}.

\ex. \ag. Jon-ga kabe-wo aka-ku nut-ta.\\
John-\textsc{nom} wall-\textsc{acc} red-\textsc{pred} paint-\textsc{pst}\\
`John painted the wall red.’ \\ \null\hfill resultative \citep[153]{nishiyama:2005:baker}
\bg. utsukushi-ku nar-u\\
beautiful-\textsc{pred} become-\textsc{prs}\\
`become beautiful' \null \hfill inchoative
\cg. utsukushi-ku sur-u\\
beautiful-\textsc{pred} do-\textsc{prs}\\
`make beautiful' \null \hfill causative
\dg. haya-ku arui-ta\\
fast-\textsc{pred} walk-\textsc{pst}\\
`walked fast' \null \hfill adverbial

Next we come to past contexts. Consider first the following example in which an \textit{i}-adjective is inflected for past and the morphological complex form is \textsc{ku-at-ta}. A \is{focus}focus-sensitive particle occurs between the predicative and the dummy copula, as in the examples above.

\exg. Hanako-wa kanashi-ku sae at-ta.\\
Hanako-\textsc{top} sad-\textsc{pred} \textsc{emp} \textsc{cop-pst}\\
`Hanako was even sad.' \null\hfill \citep[109]{kubo:1992:japanese}

In past contexts, when the three elements occur uninterrupted, \textit{i}-adjectives co-occur with the complex form \textsc{k-at-ta}, not \textsc{ku-at-ta}, as given below.

\exg. Ano doresu-wa taka-k-at-ta.\\
that dress-\textsc{top} expensive-\textsc{pred-cop-pst}\\
`That dress was expensive.'

\largerpage
We can conclude that \textsc{ku-at-ta}, or \textsc{k-at-ta}, contain the same elements as \textsc{de-at-ta}, or \textsc{d-at-ta}. However, there is no free alternation and \textsc{k-at-ta} has a default character. Instead of formalizing a parallel rule of contraction and vowel reduction, \citet[192–193]{nishiyama:1999:adj} treats [u] as an epenthetic vowel and simply analyzes the consonant part /k/ as \textsc{pred}. This analysis receives support from phonological observations. As Japanese favors open syllables, an epenthetic [u] is often inserted in word-final position as well \citep{kubozono:2005:rendaku, ito:2015:word}. The analysis also becomes relevant for the treatment of the simple element \textsc{i} (see Section \ref{kitoi}) and will be adopted in the spell-out rule below.

Before closing this subsection, I point out two more scenarios in which the \is{predicative copula}predicative copulas \textsc{de} and \textsc{k(u)} appear. One concerns negation. Given what we know about predication and its syntactic realization in Japanese, it is expected that the predicative copula has to appear in contexts of negation in order to license predication.

Negation is expressed via the negative element \textsc{na}, which can probably be regarded as the negative form of the \is{dummy copula}dummy copula.\footnote{Since it inflects exactly like an \textit{i}-adjective, it is also often classified as such. Note that \textsc{i}, as it appears in predicative position, will be analyzed as an amalgamation of copulas and tense markers following \textit{i}-adjectives below, and is therefore not glossed for now.} In the case of \textit{na}-adjectives, the topic marker -\textit{wa} additionally intervenes between this element and the predicative copula  \ref{negation na long}.\footnote{Note that -\textit{wa} could very well also be analyzed as a focus marker here, as done in \citet[6]{yamada:2023:looking}, but nothing hinges on this.} In colloquial speech, /de/ and /wa/ assimilate to /ja/ \ref{negation na short}.\footnote{The predicative copula is also obligatory for the polite form of negation, in which the existential verb, i.e. the dummy copula \textsc{ar}, is followed by the polite auxiliary element -\textit{mas}-. 

\par

\exg. Ano doresu-wa shikku-de-wa-ari-mase-n.\\
that dress-\textsc{top} chic-\textsc{pred-top-cop-aux.pol-neg}\\
`That dress is not chic.'

\par

}

\exg. Ano doresu-wa utsukushi-ku-nai.\\
that dress-\textsc{top} beautiful-\textsc{pred-neg}\\
`That dress is not beautiful.' \null \hfill negation, \textit{i}-adjective \label{negation i}

\ex. \ag.  Ano doresu-wa shikku-de-wa-nai.\\
that dress-\textsc{top} chic-\textsc{pred-top-neg} \label{negation na long}\\
\bg. Ano doresu-wa shikku-ja-nai.\\
that dress-\textsc{top} chic-\textsc{pred.top-neg}\\
`That dress is not chic.' \null \hfill negation, \textit{na}-adjective \label{negation na short}

On a final note, \textsc{pred} has another realization, namely /ni/, as argued in \citet{nishiyama:1999:adj}. This element co-occurs with \textit{na}-adjectives in tenseless environments, such as resultative, inchoative and causative constructions, but not in depictives. Under such an analysis, it can be given the same status as \textsc{k(u)}, meaning that \textit{na}-adjectives show an alternation in the realization of \textsc{pred} that \textit{i}-adjectives do not show (\citealt[205–208]{nishiyama:1999:adj}).\footnote{Again, I present Nishiyama’s analysis of adverbially used \textit{na}-adjectives as co-occurring with \textsc{pred} at face value without comment, knowing there is yet to be found empirical proof for these claims. An alternative could be to treat /ni/ as the (homophonous) dative marker, as suggested to me by Jim Kotopoulis.}

\ex. \ag. Jon-ga kabe-wo makka-ni nut-ta.\\
John-\textsc{nom} wall-\textsc{acc} crimson.red-\textsc{pred} paint-\textsc{pst}\\
`John painted the wall crimson red.’ \\ \null\hfill resultative \citep[207]{nishiyama:1999:adj}
\bg. shizuka-ni nar-u\\
quiet-\textsc{pred} become-\textsc{prs}\\
`become quiet' \null \hfill inchoative
\cg. shizuka-ni sur-u\\
quiet-\textsc{pred} do-\textsc{prs}\\
`make quiet' \null \hfill causative
\dg. shizuka-ni arui-ta\\
quiet-\textsc{pred} walk-\textsc{pst}\\
`walked quietly' \null\hfill adverbial \citep[150]{nishiyama:2005:baker}

To summarize, Japanese provides lexical evidence for the head \textsc{pred}, which is conditioned in its realization by word class, syntactic environment and phonology. I propose the following spell-out rules adopted from \citet[207]{nishiyama:1999:adj}, where [+V] represents resultative, inchoative and causative use of \textit{na}-adjectives respectively. These rules mean that \textsc{pred} is always realized as /k/ when preceded by an \textit{i}-adjective, as /ni/ when preceded by a \textit{na}-adjective in a tensed environment and otherwise as /de/.

\ex. Spell-out rule for \textsc{pred} in Japanese\\
$[Pred]$\phantom{; +V} $\leftrightarrow$ /k/ / iA \_ \\
$[Pred; +V]$ $\leftrightarrow$ /ni/ / naA \_ \\
\phantom{$[Pred; +V]$ $\leftrightarrow$} /de/

\subsubsection{Structural representation of predicative use}

Having established the ingredients of the \is{predicative copula}morphologically complex forms \textsc{de-at-ta}/\textsc{d-at-ta}, \textsc{de-ar-u} and \textsc{k-at-ta}, I will now derive the relevant structure, which I adopt from \citet[189–194]{nishiyama:1999:adj}. I will start with the predicative structures for past. See examples \ref{na-adjective past} and \ref{i-adjective past} and the corresponding \figref{na-adjective past tree} and \figref{i-adjective past tree}.

\exg. Yoru-ga shizuka-d(e)-at-ta.\\
night-\textsc{nom} quiet-\textsc{pred-cop-pst}\\
`The night was quiet.' \label{na-adjective past}

\exg. Doresu-ga utsukushi-k-at-ta.\\
dress-\textsc{nom} beautiful-\textsc{pred-cop-pst}\\
`The dress was beautiful.' \label{i-adjective past}

\begin{figure}
\centering
\begin{tikzpicture}[baseline=-0pt]
\tikzset{level distance=25pt, sibling distance=5pt}
\Tree [.TP [.NP \edge[roof]; {yoru} ] [.T′ [.CopP [.PredP [.AP \edge[roof]; {shizuka} ] [.Pred {de} ] ] [.Cop {ar} ]] [.T {ta} ]]]
\end{tikzpicture}
\caption{Predicative \textit{na}-adjectives in past (\citealt[189, 194]{nishiyama:1999:adj}; = \ref{na-adjective past})}
\label{na-adjective past tree}
\end{figure}

\begin{figure}
\centering
\begin{tikzpicture}[baseline=-0pt]
\tikzset{level distance=25pt, sibling distance=5pt}
\Tree [.TP [.NP \edge[roof]; {doresu} ]  [.T′ [.CopP [.PredP [.AP \edge[roof]; {utsukushi} ] [.Pred {ku} ] ] [.Cop {ar} ]]  [.T {ta} ]]]
\end{tikzpicture}
\caption{Predicative \textit{i}-adjectives in past (\citealt[190, 192]{nishiyama:1999:adj}; = \ref{i-adjective past})}
\label{i-adjective past tree}
\end{figure}

\newpage
These structures are exactly parallel, the only difference being the different spell-out of \textsc{pred}. The structure for non-past \textsc{de-ar-u} is also exactly analogous (mind the different tense marker for non-past), as shown in \ref{na-adjective non-past} and \figref{na-adjective non-past tree}.

\exg. Yoru-ga shizuka-de-ar-u.\\
night-\textsc{nom} quiet-\textsc{pred-cop-prs}\\
`The night is quiet.' \label{na-adjective non-past}

\begin{figure}
\centering
\begin{tikzpicture}[baseline=-0pt]
\tikzset{level distance=25pt, sibling distance=5pt}
\Tree [.TP [.NP \edge[roof]; {yoru} ] [.T′ [.CopP [.PredP [.AP \edge[roof]; {shizuka} ] [.Pred {de} ] ] [.Cop {ar} ]] [.T {u} ]]]
\end{tikzpicture}
\caption{Predicative \textit{na}-adjectives in non-past (\citealt[189, 194]{nishiyama:1999:adj}; = \ref{na-adjective non-past})}
\label{na-adjective non-past tree}
\end{figure}

\subsection{Application to attributive use}

As the past forms \textsc{k-at-ta} and \textsc{d(e)-at-ta} as well as the non-past form \textsc{de-ar-u} can be used attributively without morphological difference to their predicative counterparts, the predicative structure can simply be translated into an attributive structure. According to Nishiyama, attributively used adjectives project an additional CP layer whose head is empty and contains a relative clause feature (\citealt[197]{nishiyama:1999:adj}; cf. \citealt{kaplan:1995:the}). In the tree in \figref{attributive na-adjective past tree}, based on \ref{attributive na-adjective past}, I depict the structure for a DP containing a \textit{na}-adjective inflected for past. The structures for \textsc{k-at-ta} and \textsc{de-ar-u} are analogous apart from the different realization forms of \textsc{pred} and I.

\exg. shizuka-d(e)-at-ta yoru\\
quiet-\textsc{pred-cop-prs} night\\
`the night which was quiet' \label{attributive na-adjective past}

\begin{figure}
\centering
\begin{tikzpicture}[baseline=-0pt]
\tikzset{level distance=25pt, sibling distance=5pt}
\Tree [.NP [.CP [.TP [.CopP [.PredP [.AP \edge[roof]; {shizuka} ] [.Pred {de} ] ] [.Cop {ar} ] ] [.T {ta} ] ] [.C {[RC]} ] ] [.NP \edge[roof]; {yoru} ] ]
\end{tikzpicture}
\caption{Attributive \textit{na}-adjectives in past (\citealt[197]{nishiyama:1999:adj}; = \ref{attributive na-adjective past})}
\label{attributive na-adjective past tree}
\end{figure}

\subsection{Application to predicative DA and I} \label{kitoi}

Next, the analysis can be extended to the predicative elements \textsc{da} and \textsc{i}.

\subsubsection{Predicative DA}
\largerpage[1.5]
It is a standard assumption that \textsc{da} is a variant of the canonical copula \textsc{de-ar-u} \citep{urushibara:1993:syntactic} and these elements have been treated similarly in this book. There are several pieces of evidence for this claim (cf. \citealt{miyama:2011:da}). First, those elements are almost equally distributed, and they are mutually exclusive. However, while \textsc{de-ar-u} can appear in almost all environments (except before sentence-final particles where it is marked, \citealt[45–46]{miyama:2011:da}), \textsc{da} cannot. Instances where it is not possible include attributive structures, and, as mentioned above, before certain grammatical elements such as \textit{beki} `should’, as well as in emphatic constructions. Second, that \textsc{da} is likely historically derived from \textsc{de-ar-u} is suggested by the in-between stages \textsc{de-ar} and \textsc{dea} \citep{nakahara:2002:japanese}.

From a morphosyntactic perspective, a corresponding analysis of \textsc{da} to \textsc{de-ar-u} seems to be pretty straightforward as well -- at least to a certain point. Morphologically, the initial consonant /d/ can receive an analysis as \is{predicative copula}\textsc{pred}, exactly parallel to /d/ in \textsc{d-at-ta}. The vowel /a/ is a realization of the \is{dummy copula}dummy copula \textsc{cop}.

It is crucial to note that \textsc{da} cannot be an allomorph of \textsc{de}, since the two forms are not paraphrasable, not even in emphatic constructions. Compare \ref{danotde} to \ref{shikkudemo}.

\exg. *shizuka-da mo ar-u\\
quiet-\textsc{da} \textsc{emp} \textsc{cop-prs}\\
(`it is even quiet’) \null \hfill \citep[186]{nishiyama:1999:adj} \label{danotde}
\clearpage

What is striking is the alleged absence of a tense marker. Based on the fact that \textsc{da} and \textsc{de-ar-u} are, in most environments, interchangeable, Nishiyama simply treats the two as allomorphs and assumes that all three elements -- predicative copula, dummy copula and tense marker for non-past -- are combined in /da/ as shown below \citep[196–197]{nishiyama:1999:adj}.

\exg. Yoru-ga shizuka-da.\\
night-\textsc{nom} quiet-\textsc{pred.cop.prs}\\
`The night is quiet.' \null \hfill (Nishiyama's analysis)

Nothing immediately hinges on this, but I will assume here that the predicative copula is the consonant /d/ and that the \is{dummy copula}dummy copula and the tense marker for present are merged in /a/ as shown below.

\ex. \textsc{d-a} = \textsc{pred-cop.prs}

This means that \textsc{da} is not derived from \textsc{de-ar-u}, but simply is a shorter form, in which all relevant elements, including the tense marker, are still present.

\subsubsection{Predicative I} \label{predicativei}

A similar structure can be derived for \textit{i}-adjectives, although predicative \textsc{i}, at first glance, does not seem to bear any resemblance to \textsc{ku-ar-u}. First, note that non-past is the only environment in which the velar consonant /k/ is absent in the inflectional paradigm of \textit{i}-adjectives. In \is{Classical Japanese} Classical Japanese, however, the attributive form of the inflectional paradigm of \textit{i}-adjectives was -\textsc{ki}, thus included this consonant (cf. \citealt[189–190]{nishiyama:1999:adj}). \citet{nishiyama:1999:adj} accounts for these facts as follows: Modern \textsc{i} is a direct successor of \textsc{ki} derived via a historical \is{onbin@\textit{onbin}}devoicing process (cf. \citealt{frellesvig:2010:history}). The consonant /k/ was an instance of \textsc{pred} and continued this role in Modern Japanese, although it is not pronounced. He formalizes this through an operation which he calls \is{velar filter}\textit{Velar Filter} \citep[191]{nishiyama:1999:adj}.

In other words, the head of \is{predicative copula}PredP is structurally present, but phonologically reduced,\footnote{Such phonological processes are not uncommon in the history of Japanese, where the consonant /k/ shows systematic reducion in front of the vowel /i/ in diachronic perspective \citep[340]{frellesvig:2010:history}. Furthermore, they are also attested synchronically, especially in certain verb forms. For example, the past form of \textit{kak} `to write' is \textit{kai-ta}, not *\textit{kaki-ta} \citep[191]{nishiyama:1999:adj}.} so \textit{i}-adjectives co-occur with the underlying elements \textsc{\st{k}-i}, which constitute a combination of the \is{predicative copula}predicative copula and a tense marker and are realized as /i/ (cf. \citealt[199–201]{nishiyama:1999:adj}; cf. also \citealt[2]{baker:2003:verbal}).

\ex. Nishiyama's analysis for \textit{i}-adjectives in non-past
\ag. Ano doresu-ga taka-\st{k}-i.\\
that dress-\textsc{nom} expensive-\textsc{pred-prs}\\
`That dress is expensive.'

But under this analysis no dummy copula is present. Recall that this is also the case in other languages in default non-past contexts, for instance \is{Hebrew} Hebrew \citep{doron:1983:verbless, rapoport:1987:copular} and \is{Russian}Russian \citep{babby:2010:syntactic}, as in the example in \ref{pagoda} repeated below. This can thus be applied to Japanese (\citealt[191–192]{nishiyama:1999:adj}, and see \citealt{urushibara:1993:syntactic}).

\exg. Pagoda xorošaja. (Russian)\\
weather.\textsc{nom.sg.f} good.\textsc{nom.sg.f}\\
`The weather is good.'

An obvious challenge to this account is the question how tense is supported, if no \is{dummy copula}dummy copula is present. As explained above, the dummy copula is needed to support the I node and establishes some kind of \textit{be}-support. And while no overt tense is present in the Russian example, Nishiyama argues that non-past information is included in \textsc{i}. Furthermore, how can the parametric difference between \textit{i}- and \textit{na}-adjectives, that necessarily arises, be accounted for? For lack of a better alternative, I assume for predicative \textsc{i} almost the same structure as for \textsc{da}: The predicative copula is represented via the (unpronounced) consonant /k/, and the dummy copula and the tense morpheme are \is{velar filter}merged in /i/. \is{Classical Japanese}

\ex. \textsc{(k)-i} = \textsc{pred-cop.prs}

\subsection{Why the application to attributive I and NA fails} \label{why it fails}

This section refutes Nishiyama's analysis of the attributive forms for non-past, \textsc{i} and \textsc{na}, as copular elements. It is not surprising that \citet{nishiyama:1999:adj} puts forth the same analysis for \textsc{i} in attributive contexts as for \textsc{i} in predicative contexts, thus treats those as the same elements, see \ref{example attributive i-adjective Nishiyama}.

\exg. taka-\st{k}-i doresu\\
expensive-\textsc{pred-prs} dress\\
`dress which is expensive' \label{example attributive i-adjective Nishiyama}

The only structural difference is that the adjective projects an additional CP layer as in \figref{attributive i-adjective non-past tree wrong} \citep[199–202]{nishiyama:1999:adj}.

\begin{figure}
\centering
 \begin{tikzpicture}[baseline=-0pt]
\tikzset{level distance=25pt, sibling distance=5pt}
\Tree [.NP [.CP [.TP [.PredP [.AP \edge[roof]; {taka} ] [.Pred {\st{k}} ] ] [.T {i} ] ] [.C {[RC]} ] ] [.NP \edge[roof]; {doresu} ] ]
\end{tikzpicture}
\caption{Attributive \textit{i}-adjectives in non-past according to Nishiyama (not correct), based on \ref{example attributive i-adjective Nishiyama}}
\label{attributive i-adjective non-past tree wrong}
\end{figure}

For \textit{na}-adjectives, \citet[196–199]{nishiyama:1999:adj} analyzes \textsc{na} as the very same element as \textsc{da}, thus as a combination of \is{predicative copula}predicative copula, dummy copula and tense marker for non-past. The only difference in consonant realization is conditioned by the presence of a CP. This means that the labor between predicative and attributive use is strictly divided in \textit{na}-adjectives via \textsc{da} and \textsc{na}, which are allomorphs and not interchangeable \citep[197]{nishiyama:1999:adj}. He gives the following phonological transformation rule in the spirit of Distributed Morphology \citep{halle:1993:distributed, halle:1994:some}. \is{Distributed Morphology}

\ex. [pred.cop, dum.cop, -past, rel.cl] $\leftrightarrow$ /na/ / NA \_ \\ \null \hfill \citep[197]{nishiyama:1999:adj}

\citet{nishiyama:1999:adj} treats all attributive adjectives in Japanese as embedded in a CP layer and therefore as indirect modifiers.\footnote{Nishiyama himself does consider the possibility of direct modification. He states ``[...] we know that natural language allows adjectives to act as direct modifiers (as in English). Thus, unless there is evidence to the contrary, the null hypothesis is that Japanese adjectives also act as direct modifiers [...]” \citep[201]{nishiyama:1999:adj} and mentions \is{MOD@\textsc{mod}}\textsc{mod} presented in Section \ref{internalstructure} as an alternative.} \citet{baker:2003:verbal} and \citet{hoshi:2002:kaynean} adopt his analysis and explicitly deny the possibility of direct modification in Japanese.

Nishiyama not only encourages the long-existing analysis of \textsc{na} as the attributive form of the copula \citep{teramura:1982:nihongo, miyagawa:1987:lexical, murasugi:1991:noun, kubo:1992:japanese, urushibara:1993:syntactic}, but also manages to unify \textit{i}- and \textit{na}-adjectives by analyzing \textsc{i} in a parallel fashion. This is a significant contribution. However, this analysis overgeneralizes as presented in abundance throughout this book. Although sometimes up for discussion, many direct \textit{i}- and \textit{na}-adjectives were presented in the preceding chapter (see especially Section \ref{section direct modifiers}), among them the following (cf. especially \citealt{yamakido:2005:nature} or \citealt{ayano:2010:revisiting}, from where these examples are repeated). %\is{direct modifiers}

\ex. \ag. Omoigakena-i kyaku-ga ki-ta.\\
unexpected-\textsc{mod} guest-\textsc{nom} come-\textsc{pst}\\
`An unexpected guest came.' \null \hfill \citep[67]{yamakido:2005:nature}
\bg. *Ano kyaku-ga omoigakena-i.\\
that guest-\textsc{nom} unexpected-\textsc{pred.cop.prs}\\
(`That guest is unexpected.’)

\ex. \ag. Max-ga kanzen-na baka-d-a.\\
Max-\textsc{nom} complete-\textsc{mod} idiot-\textsc{pred-cop.prs}\\
`Max is a complete idiot.' \null \hfill \citep[68]{yamakido:2005:nature}
\bg. *Ano baka-ga kanzen-d-a.\\
that idiot-\textsc{nom} complete-\textsc{pred-cop.prs}\\
(`That idiot is complete.’)

This means \textsc{i} and \textsc{na} in attributive position cannot be universally analyzed as copular elements, which is why I reject this analysis in favor of \is{MOD@\textsc{mod}}\textsc{mod} (Section \ref{internalstructure}). In turn, this means there are two different instances of \textsc{i}: a copular element in predicative position and the syntactic linker \textsc{mod} in attributive position.

\subsection{Application to \textit{no}-modifiers}

The analysis presented here can be straightforwardly extended to \textit{no}-modifiers, which share their entire inflectional paradigm with \textit{na}-adjectives with the exception of \textsc{no}. All the elements receive the same analysis, illustrated first with \textsc{da} and \textsc{de-ar-u} in combination with a clear noun \ref{clear noun}, and a more adjectival \textit{no}-modifier, i.e. a \textit{no}-adjective. \ref{noadjective}. See also \citet[195]{nishiyama:1999:adj}. %\is{no-modifiers}

\ex. \label{clear noun} \ag. Hanako-wa sensei-d-a.\\
Hanako-\textsc{top} teacher-\textsc{pred-cop.prs}\\
\bg. Hanako-wa sensei-de-ar-u.\\
Hanako-\textsc{top} teacher-\textsc{pred-cop-prs}\\
`Hanako is a teacher.'

\ex. \label{noadjective} \ag. Haishīzun-ni-wa hayame-no yoyaku-ga hisshu-d-a.\\
high.season-\textsc{loc-top} early-\textsc{mod} reservation-\textsc{nom} necessary-\textsc{pred-cop.prs}\\
\bg. Haishīzun-ni-wa hayame-no yoyaku-ga hisshu-de-ar-u.\\
high.season-\textsc{loc-top} early-\textsc{mod} reservation-\textsc{nom} necessary-\textsc{pred-cop-prs}\\
`During high season, early reservation is necessary.’ \\ adapted from: Ishimori Hiromi (2002): Tabi, chotto senchimentaru. Tokyo: Tōyō Shuppan, via BCCWJ, 2023/08/05.

The analysis also carries over to \textsc{de-at-ta}, which can-occur with a \textit{no}-modifier as well.

\exg. Hanako-wa sensei-de-at-ta.\\
Hanako-\textsc{top} teacher-\textsc{pred-cop-pst}\\
`Hanako was a teacher.' \null \hfill noun, past

Naturally, this extends to the attributive use of \textsc{de-ar-u} and \textsc{de-at-ta}. Just as \textsc{da} and \textsc{de-ar-u} are mutually exclusive in predicative position, \textsc{no} and \textsc{de-ar-u} are mutually exclusive in attributive position.

\exg. sensei-de-ar-u Hanako\\
teacher-\textsc{pred-cop-prs} Hanako\\
`Hanako, who is a teacher'

Although Nishiyama does not focus on \textit{no}-modifiers, he proposes a spell-out rule for \textsc{no} as well, which looks as follows.

\ex. [pred.cop, dum.cop, -past, rel.cl] $\leftrightarrow$ /no/ \_ \\ \null \hfill (Nishiyama’s analysis of \textsc{no}; \citealt[197]{nishiyama:1999:adj})

This spell-out rule says that \textsc{no} is a manifestation of the \is{predicative copula}predicative copula, dummy copula and non-past tense marker and appears in relative clause contexts. It is realized in the elsewhere-case, meaning when the modifier is not a clear \textit{na}-adjective.

As in the case of the other modifiers, this cannot be the right analysis given the abundance of direct modifiers among the group of \textit{no}-modifiers. Instead, I have put forth an analysis in Section \ref{internalstructure} of \textsc{no} as an instance of \is{MOD@\textsc{mod}}\textsc{mod} which is blocked when \textsc{de-ar-u} or other elements with an overt I-head appear.

However, it is worth emphasizing that the existence of such cases, those in which a \textit{no}-modifier has a clear predicative function, show that not all \textit{no}-modifiers are direct. This, in turn, refutes a claim made in \citet{nagano:2015:RA}, who argue, ``Japanese adjectives differ from English adjectives in formally disambiguating direct and indirect modification'' \citep[124]{nagano:2015:RA}. According to the authors, it is only \textit{no}-modifiers which are responsible for direct modification, all other modifiers in Japanese being indirect. However, this is simply not true, as shown in many examples throughout this book and in the data set. %\is{no-modifiers}

\subsection{The integration of indirect modifiers into the DP}

The complex morphological elements deduced in this section can be neatly integrated along with the modifier into a cartographic structure of the DP. This is exemplified in example \ref{example integration} with a \textit{na}-adjective inflected for past, but the derivation is exactly parallel for the other elements for past (\textsc{d-at-ta} for \textit{no}-modifiers, \textsc{k-at-ta} for \textit{i}-adjectives) as well as \textsc{de-ar-u} for non-past.

\exg. ano shizuka-de-at-ta yoru\\
that quiet-\textsc{pred-cop-pst} night\\
`that night which was quiet' \label{example integration}

Structurally, such indirect modifiers are merged in a functional projection in the indirect domain and are embedded in a complex set of functional layers whose heads spell-out the predicative copula, the canonical copula and the tense morpheme respectively. This is displayed in \figref{cartographic tree indirect modifiers}.\footnote{In this chapter, I will not include the universal syntactic head \textsc{mod} in the visualizations of relative clauses in cases where it is not realized. The CP status of relative clauses will be defended in the following.}

\begin{figure}[b]
\centering
\begin{tikzpicture}[baseline=-0pt]
\tikzset{level distance=25pt, sibling distance=5pt}
\Tree [.DP\textsubscript{internal} ano [.D′\textsubscript{internal} [.FP\textsubscript{indirect} [.CP {} [.C′ [.IP {} [.I′ [.CopP {} [.Cop′ [.PredP {} [.Pred′ [.AP \edge[roof]; shizuka ] [.Pred de ] ] ] [.Cop ar ] ] ] [.I ta ] ] ] C ]] [.F′\textsubscript{indirect} [.... [.NP \edge[roof]; yoru ] {} ] F ] ] D ]]
\end{tikzpicture}
\caption{Integration of indirect modifiers and their complex structure in cartographic DP (cf. \ref{example integration})}
\label{cartographic tree indirect modifiers}
\end{figure}

\subsection{A note on the polite elements} \label{desu}

At this point, the question remains why the \is{polite copula}polite copula \textsc{desu}, glossed as \textsc{cop.pol} up to now, and its counterpart for past \textsc{deshita} are mutually exclusive with the copular forms \textsc{de-ar-u} and \textsc{de-at-ta}, but follow \textsc{i} and also \textsc{k-at-ta}.\footnote{Assuming there are three stages of politeness in Japanese \citep{minami.1977:gendai}, \textsc{desu} expresses the medial level and is used when talking to strangers and superiors. In the verbal domain, the corresponding politeness is achieved via the polite auxiliary -\textit{mas}-, in past \textit{mashi-ta}. Also note that some other grammatical elements, for example the epistemic suffix \textit{mitai} `seem’, occur after \textsc{i} but replace \textsc{da} \citep{kubo:1992:japanese, urushibara:1993:syntactic, nishiyama:1999:adj}.} Observe some relevant examples.

\ex. \ag. Ano doresu-wa shikku (*d-a) desu.\\
that dress-\textsc{top} chic \textsc{pred-cop.prs} \textsc{desu}\\
`That dress is chic.' \null \hfill polite non-past, \textit{na}-adjective
 \bg. Ano doresu-wa shikku (*d-at-ta) deshi-ta.\\
that dress-\textsc{top} chic \textsc{pred-cop-pst} \textsc{desu}-\textsc{pst}\\
`That dress was chic.' \null \hfill polite past, \textit{na}-adjective


\ex. \ag. Ano doresu-wa taka-*(i) desu.\\
that dress-\textsc{top} expensive-\textsc{pred.cop.prs} \textsc{desu}\\
`That dress is expensive.' \null \hfill polite non-past, \textit{i}-adjective
\bg. Ano doresu-wa taka-k-at-ta desu.\\
that dress-\textsc{top} expensive-\textsc{pred-cop-pst} \textsc{desu}\\
`That dress was expensive.' \null \hfill polite past, \textit{i}-adjective \label{katta desu not deshita}

\ref{katta desu not deshita} is the only way to express \is{polite copula}politeness and past tense on an \textit{i}-adjective at the same time. The polite past form \textsc{deshita} never co-occurs with \textit{i}-adjectives.

\exg. *Ano doresu-wa taka-i deshi-ta.\\
that dress-\textsc{top} expensive-\textsc{pred.cop.prs} \textsc{desu}-\textsc{pst}\\
(`That dress was expensive.’) (polite)

In a recent analysis, \citet{yamada:2023:looking}, and see also \citet{miyagawa:2022:syntax}, argues that Modern Japanese has two separate instances of \textsc{desu}, which correspond to two different grammars. This explains the different distribution. In a nutshell, the element \textsc{desu} following \textit{na}-adjectives (and nouns/\textit{no}-modifiers) consists of the familiar tripartite of elements. The element \textsc{desu} following \textit{i}-adjectives, on the other hand, is just a single marker of politeness and crucially follows the copula \citep[2–4]{yamada:2023:looking}.

In the first case, \textsc{desu/deshita} consist of the identical elements as \textsc{de-ar-u/de-at-ta}, but significantly, politeness is marked. Translated into the analysis proposed here, I treat -\textsc{u} and -\textsc{ta} as tense morphemes, \textsc{de} as the \is{predicative copula}predicative copula, and the medial elements -\textsc{s}- and -\textsc{shi}- as the polite version of the \is{dummy copula}dummy copula.\footnote{Note that the difference in consonant quality between the consonants `s' and `sh' is simply phonologically governed (see \citealt{ito:2015:word}). The epenthetic vowel [i] has to be inserted after the consonant [s], preserving this consonant as well as the CV mora structure favored in Japanese, and becomes the coda of the new mora [\textctc i] (/shi/). The consonant shift of the preceding consonant from [s] to [\textctc] is systematic in the sense that the fricatives [s] and [z] always become alveolo-palatal before [i]. The mora [si] simply is never realized in Japanese \citep{lewin:1989:sprache, vance:2008:sounds}.} This means that the dummy copula is not simply semantically vacuous, but has a semantic and pragmatic dimension. Nevertheless, it needs to support tense. Therefore, I maintain the gloss \textsc{cop.pol} for the medial part, highlighting that it is a polite version of the dummy copula, as shown below.

\ex. \ag. Yoru-ga shizuka-de-s-u.\\
night-\textsc{nom} quiet-\textsc{pred-cop.pol-prs}\\
`The night is quiet.' \null \hfill polite non-past, \textit{na}-adjective
\bg. Yoru-ga shizuka-de-shi-ta.\\
night-\textsc{nom} quiet-\textsc{pred-cop.pol-pst}\\
`The night was quiet.' \null \hfill polite past, \textit{na}-adjective

\largerpage
\textsc{desu} following \textit{i}-adjectives and their copular forms is a sort-of grammaticalized modal element and by itself projects the syntactic category \textsc{pol} (cf. \citealt[10–11]{yamada:2023:looking}). There is no need to express an additional dummy copula. I suggest the following structure.\footnote{I thank Ken Hiraiwa for suggesting and discussing these patterns with me. The analysis in \citet{yamada:2023:looking} is similar in that it proposes the emergence of a new grammar with a re-analysis of \textsc{desu} as a marker for addressee honorificity. Note that this politeness marker also follows the negative element \textsc{na}, which inflects like an \textit{i}-adjective.}

\exg. Ano doresu-wa taka-k-at-ta desu.\\
that dress-\textsc{top} expensive-\textsc{pred-cop-pst} \textsc{pol}\\
`That dress was expensive.' \null \hfill polite past, \textit{i}-adjective \label{katta desu}

There are several pieces of evidence for this claim. First, \textsc{desu} in \ref{katta desu} is syntactically not necessary, in contrast to its combination with \textit{na}-adjectives. This means it has a syntactic dummy status, albeit a more semantic contribution. Furthermore, \textsc{desu} in such instances is never paraphrasable with \textsc{de-ar-u} as shown below (and see \citealt[10–11]{yamada:2023:looking}).

\ex. \ag. *Ano doresu-wa taka-k-at-ta-de-ar-u.\\
that dress-\textsc{top} expensive-\textsc{pred-cop-pst-pred-cop-prs}\\
\bg. *sensei-de-wa-na-i-de-ar-u\\
teacher-\textsc{pred-top-neg-pred.cop.prs-pred-cop-prs}\\

Furthermore, use of the \is{polite copula}grammaticalized version of \textsc{desu} has in recent times increased. It is also attested after \textsc{d-at-ta}, where the copular elements are already present as well \citep{yamada:2023:looking}.

\exg. I-i mise d-at-ta desu.\\
good-\textsc{mod} restaurant \textsc{pred-cop-pst} \textsc{pol}\\
`It was a good restaurant.’ \null \hfill \citep[13]{yamada:2023:looking} \label{desudatta}

\citet{yamada:2023:looking} also gives examples where \textsc{desu} appears multiple times inside one sentence, each time followed by a sentence-final particle \ref{desufinal}.

\exg. Watashi-wa desu ne totemo desu ne tsukare-ta desu yo.\\
I-\textsc{top} \textsc{pol} \textsc{sfp} very \textsc{pol} \textsc{sfp} become.tired-\textsc{pst} \textsc{pol} \textsc{sfp}\\
`I (you know) got very (you know) tired (you know).’ \\ \null \hfill (\citealt[Footnote 34]{yamada:2023:looking}) \label{desufinal}

\largerpage[2]
These examples show that this instance of \textsc{desu} has been grammaticalized into an element that simply marks politeness and co-occurs with the copula. To summarize:\footnote{Note that the question is left open why \textsc{desu} can never appear in attributive position.

\par

\exg. *ano shikku desu doresu\\
that chic \textsc{desu} dress\\
(`that chic dress’) (polite)

\par

Tentatively, I assume this really is a pragmatic ban, restricting politeness to the matrix clause. Some authors \citep{horie:2017:attributive} have shown that verbs in a polite speech style are in some instances possible in attributive position, especially when the root verb is in the very polite speech style. Apparently, such constructions are common in the service domain.

\par

\exg. [O-mōshikomi-ni nari-mashi-ta] kādo-wo hakkō-sase-te-itadaki-mas-u.\\
\textsc{hon}-apply-\textsc{dat} become-\textsc{aux.pol-pst} card-\textsc{acc} issue-cause.\textsc{pol}-\textsc{ger}-humbly.receive-\textsc{aux.pol-prs}\\
`We shall respectfully issue a card [for which you had sent your application form].’ \\ \null \hfill \citep[51]{horie:2017:attributive}

\par

It is not necessary, however, to use the polite form of the verb in the embedded clause, and potentially this is a new trend to be as polite as possible that just has not reached the adjectival domain yet (Ken Hiraiwa p.c.).}

%

\ex.\a. \textsc{de-s-u} = \textsc{pred-cop.pol-prs} (\textit{na}-adjectives and nouns)
\b. \textsc{d(e)-shi-ta} = \textsc{pred-cop.pol-pst} (\textit{na}-adjectives and nouns)
\c. \textsc{desu} = \textsc{pol} (\textit{i}-adjectives and elsewhere)

The element \textsc{de-s-u} is another form of the copula, and, as shown in the previous chapter (Section \ref{test predication}), equally incompatible with direct modifiers. It contrasts with another element \textsc{desu} used to mark \is{polite copula}politeness.\footnote{The analysis of the element \textsc{desu} as monostructural proposed here is, in principle, consistent with the analysis given in \citet[64–67]{miyagawa:2022:syntax}. He analyzes this element as \textit{modal} and the head of a Mod(al)P, which, as is evident from the structural precedence, occurs higher than TP.

\par

Note that the analysis given here can be carried over to the \is{epistemic copula}epistemic element \textsc{d(e)arō}, which is, like \textsc{desu}, mutually exclusive with the copula \textsc{d-a}, but has to follow \textsc{i} and any tense morpheme when co-occurring with \textit{i}-adjectives and verbs. Tentatively, I propose the following segmentation.

\par


\ex. \a. \textsc{d(e)-ar-ō} = \textsc{pred-cop-epis} (\textit{na}-adjectives and nouns)
\b. \textsc{d(e)arō} = \textsc{epis} (\textit{i}-adjectives and elsewhere)

\par

}

\subsection{Summary}

This section has adopted the landmark analysis of Nishiyama for the complex elements Japanese (non-verbal) modifiers co-occur with, showing they are combinations of a tense marker and two different copulas. This analysis can readily be incorporated into the cartographic framework assumed here. Next, I will discuss the size of indirect modifiers in Japanese.

\section{The size of indirect modifiers in Japanese} \label{size,indirect}

As shown in the introduction to this chapter, the size of Japanese relative clauses, whether they are IPs or CPs, is difficult to determine. Interest in this matter has increased since the end of the 1980s. At that time, the null hypothesis was that cross-linguistically all relative clauses are CPs. In her doctoral thesis and subsequent work, \citet{murasugi:1991:noun, murasugi:2000:anti} explicitly advocated for the IP-hypothesis for Japanese. This was picked up by Kayne (\citeyear{kayne:1994:anti}, and see \citealt{kayne:2005:japanese}), who argued that relative clauses in prenominal-RC languages are derived via movement of the remnant IP to the prenominal position. On the other hand, the CP-hypothesis is explicitly defended for example in \citet{kaplan:1995:the} and \citet{miyamoto:2017:RC}.\footnote{See the overview and discussion in \citet{hoshi:2004:para} and \citet{miyamoto:2017:RC}.} Furthermore, some authors \citep{krause:2001:rrc, miyagawa:2011:genitive, miyagawa:2012:case, miyagawa:2017:agreement} have argued for a non-unified status, some \citep{sells:2007:finiteness} proposed abandoning the difference in structural size in Japanese altogether, and others \citep{matsumoto:2017:japanese} take the position that Japanese relative clauses should be analyzed entirely differently from relative clauses in European languages.

In this section, I show that relative clauses in Japanese are only of one size, and propose a CP status. I first demonstrate they must be at least IPs, then dismiss accounts which argue that they cannot be bigger as well as those which argue they come in different sizes, before eventually showing why they are CPs. The most important point to take away, however, is that Japanese relative clauses come only in one size, meaning there is no dichotomy between CP-IP or full-reduced. This has serious consequences for Cartography. Instead of distinct merging positions for reduced and full relative clauses, a unifying analysis of relative clauses predicts that they are unordered, crucially also with regard to numerals. This aspect will then be discussed in Section \ref{complexity}.

\subsection{Japanese relative clauses are at least IPs} \label{cinque adj}

\subsubsection{High adverbials}

A popular and reliable test used to determine the size of (especially \is{reduced relative clause}reduced) relative clauses in several European languages is their interaction with different \is{CP adverb}adverbs. It is well-known that adverbs attach at different levels of the clause, the CP, IP or VP-domain (cf. for instance \citealt{ramchand:2014:deriving}) and take scope over different elements. Below, I give a reconstruction of the hierarchy of adverbs proposed by \citet[109]{cinque:1999:adverbs} where everything above Asp\textsubscript{perfect}, marked in boldface, is the IP-domain.\footnote{For other hierarchies see \citet{laenzlinger:2011:elements} and \citet{morzycki:2016:mod}.}

% {
% \tiny
% \begin{forest}
% [Mood\textsubscript{speech act} \textit{frankly}  [{}] [Mood\textsubscript{evaluative} \textit{fortunately}  [{}] [Mood\textsubscript{evidential} \textit{allegedly}  [{}] [Mod\textsubscript{epistemic} \textit{probably}  [{}] [T\textsubscript{past} \textit{once}  [{}] [T\textsubscript{future} \textit{then}  [{}] [Mod\textsubscript{irrealis} \textit{perhaps}  [{}] [Mod\textsubscript{necessity} \textit{necessarily}  [{}] [Asp\textsubscript{habitual} \textit{usually}  [{}] [Asp\textsubscript{repetitive (I)} \textit{again}  [{}] [Asp\textsubscript{frequentative (I)} \textit{often}  [{}] [Mod\textsubscript{volitional} \textit{intentionally}  [{}] [Asp\textsubscript{celerative (I)} \textit{quickly}  [{}] [T\textsubscript{anterior} \textit{already}  [{}] [Asp\textsubscript{terminative} \textit{no longer}  [{}] [ Asp\textsubscript{continuative} \textit{still}  [{}] [\textbf{Asp\textsubscript{perfect}} \textit{always}  [{}] [Asp\textsubscript{retrospective} \textit{just}  [{}] [Asp\textsubscript{proximative} \textit{soon}  [{}] [Asp\textsubscript{durative} \textit{briefly}  [{}] [Asp\textsubscript{generic/progressive} \textit{characteristically}  [{}] [Asp\textsubscript{prospective} \textit{almost}  [{}] [Asp\textsubscript{sg.completive (I)} \textit{completely}  [{}] [Asp\textsubscript{pl.completive} \textit{all}  [{}] [Voice \textit{well}  [{}] [Asp\textsubscript{celerative (II)} \textit{fast/early}  [{}] [Asp\textsubscript{repetitive (II)} \textit{again}  [{}] [Asp\textsubscript{frequentative (II)} \textit{often}  [{}] [Asp\textsubscript{sg.completive (II)} \textit{completely} ]
% ]
% ]
% ]
% ]
% ]
% ]
% ]
% ]
% ]
% ]]]
% ]]]]]]]]]]
% ]
% ]]
% ]
% ]
% ]
% \end{forest}
% }

\ex.
\raggedright
[Mood\textsubscript{speech act} \textit{frankly} [Mood\textsubscript{evaluative} \textit{fortunately} [Mood\textsubscript{evidential} \textit{allegedly} [Mod\textsubscript{epistemic} \textit{probably} [T\textsubscript{past} \textit{once} [T\textsubscript{future} \textit{then} [Mod\textsubscript{irrealis} \textit{perhaps} [Mod\textsubscript{necessity} \textit{necessarily} [Asp\textsubscript{habitual} \textit{usually} [Asp\textsubscript{repetitive (I)} \textit{again} [Asp\textsubscript{frequentative (I)} \textit{often} [Mod\textsubscript{volitional} \textit{intentionally} [Asp\textsubscript{celerative (I)} \textit{quickly} [T\textsubscript{anterior} \textit{already} [Asp\textsubscript{terminative} \textit{no longer} [ Asp\textsubscript{continuative} \textit{still} [\textbf{Asp\textsubscript{perfect}} \textit{always} [Asp\textsubscript{retrospective} \textit{just} [Asp\textsubscript{proximative} \textit{soon} [Asp\textsubscript{durative} \textit{briefly} [Asp\textsubscript{generic/progressive} \textit{characteristically} [Asp\textsubscript{prospective} \textit{almost} [Asp\textsubscript{sg.completive (I)} \textit{completely} [Asp\textsubscript{pl.completive} \textit{all} [Voice \textit{well} [Asp\textsubscript{celerative (II)} \textit{fast/early} [Asp\textsubscript{repetitive (II)} \textit{again} [Asp\textsubscript{frequentative (II)} \textit{often} [Asp\textsubscript{sg.completive (II)} \textit{completely} ]

\citet{douglas:2016:syntactic} and \citet{harwood:2018:rrc} independently predict that \is{reduced relative clause}reduced relative clauses in \is{English} English and \is{French} French respectively should be compatible with all \is{CP adverb}adverbs at least below Adv\textsubscript{generic/progressive} and possibly even as high as Asp\textsubscript{perfect}, if they indeed have a status as big as IP. That is because these projections are associated with tense and also host auxiliary verbs, elements of the IP-domain. However, this should also predict that ``any adverb higher than [...] Asp\textsubscript{perfect} will be impossible in RRCs" (\citealt[223]{douglas:2016:syntactic}) given that this is the CP domain. As it turns out, reduced relative clauses in the languages they test are not only compatible with the predicted kind of adverbs. In fact, they are even compatible with very high speaker-oriented, evaluative and evidential, thus clausal, adverbs. This is shown by \citet[226–230; 239–245]{douglas:2016:syntactic} for English and French, and the same findings are observable in \is{Italian} Italian (cf. \citealt[214]{cinque:2020:relcl}).

\ex. The cakes [unfortunately\textsubscript{evaluative} / probably\textsubscript{evidential} (being) eaten by the guests] will cause stomach upsets. (English) \null \hfill \citep[226]{douglas:2016:syntactic}

\exg. L’étudiant [malheureusement\textsubscript{evaluative} / probablement\textsubscript{evidential} arrêté par la police] est un étranger. (French)\\
\textsc{det}.student unfortunately / probably arrested by \textsc{det} police be.\textsc{3.sg} a foreigner\\
`The student [probably/unfortunately arrested by the police] is a foreigner.’ \\ \null \hfill \citep[241]{douglas:2016:syntactic}

\exg. I film [sfortunatamente\textsubscript{evaluative} / presumibilmente\textsubscript{evidential} persi durante l’ultima guerra] sono davvero molti. (Italian)\\
\textsc{det} films unfortunately / presumably lost during \textsc{det}.last war are really many\\
`The films [unfortunately/presumably lost during the last] war are really many.’ \null \hfill \citep[214]{cinque:2020:relcl}

\newpage
That means that despite being able to prove IP (over for example vP) status, it cannot be determined whether \is{reduced relative clause}reduced relative clauses in those languages are just IPs or potentially bigger, that is CPs.\footnote{Furthermore, different languages might allow reduced relative clauses of different sizes and therefore combinations with different kinds of adverbs, which is also hinted at in \citet[76–77]{krause:2001:rrc}, but not further elaborated.}

This test can be applied in Japanese. First note that all types of adverbs do exist in Japanese: VP (\textit{hayaku} `fast'), IP/TP (\textit{kinō} `yesterday') and \is{CP adverb}CP adverbs (\textit{zannennagara} `unfortunately’). Important differences are first that the canonical position for CP adverbs is very high in the clause, although they can appear in different positions. More importantly, these adverbs can never be the target of negation, which is very low, as shown in \citet{koizumi:2017:subject}. Compare \ref{wakekino}, in which the TP adverb \textit{kinō} is in the scope (c-command domain) of negation, to \ref{wakeosoraku}, in which the CP adverb \textit{osoraku} `possibly’ is not in the scope of negation.

\exg. [Kinō [[hashit-ta] wake-de-wa-na-i]].\\
yesterday run-\textsc{pst} case-\textsc{pred-top-neg-pred.cop.prs}\\
`It is not the case that (I) ran yesterday.’ \label{wakekino}

\exg. Osoraku [[[hashit-ta] wake-de-wa-na-i] darō].\\
possibly run-\textsc{pst} \textsc{pred-top-neg-pred.cop.prs} \textsc{epis}\\
`Probably, it is not the case that (he or she) ran.’\\
$\neq$ `It is not the case that (he or she) probably ran.’ \\ \null \hfill \citep[556]{koizumi:2017:subject} \label{wakeosoraku}

Since different types of adverbs exist in Japanese, it should be possible to deduce a hierarchy as well. This is attempted in \citet{narrog:2009:mod}, who argues for an infinite number of modal adverbs in Japanese and, based on ordering and \is{scope}scope principles, which he investigates in a corpus study (see \citealt[71–76]{narrog:2009:mod}), gives the following hierarchy, which also contains the categories tense, aspect and negation. As already mentioned in Section \ref{narrog}, this hierarchy turns out not to be a perfect match with the hierarchy proposed in \citet{cinque:1999:adverbs} and Narrog identifies several problems.

\ex. \emph {Narrog's hierarchy of Japanese modal categories (\citealt[235]{narrog:2009:mod})}\\
MoodP speech act I (imperative) $>$ MoodP speech act II (hortative) $>$ ModP epistemic I (speculative) $>$ TP past I $>$ ModP evidential I (reportive) $>$ TP past II $>$ MoDP epistemic II (epistemic possibility) $>$ TP past III $>$ NegP INtNeg I $>$ ModP evidential II inferential evidentiality $>$ TP past IV $>$ NegP IntNeg II $>$ ModP epistemic III epistemic necessity $>$ TP past V $>$ NegP IntNeg III $>$ ModP deontic I weak deontic necessity $>$ TP past VI $>$ NegP IntNeg IV $>$ AspP result/compl I $>$ ModP deontic II strong deontic necessity $>$ TP past VII $>$ NegP IntNeg V $>$ AspP result/compl III $>$ ModP volition $>$ TP past VIII $>$ NegP IntNeg VI $>$ AspP result/compl IV $>$ ModP evidential III future-oriented $>$ NegP InstNeg VII $>$ AspP result/progress/compl II $>$ ModP root possibility I NegP IntNeg VIII $>$ AspP result/progress/compl V

Nevertheless, high \is{CP adverb}CP adverbs are attested in Japanese and can be tested in relative clauses in a similar vein. This yields the same picture as in \is{English} English or \is{French} French. Relative clauses combine rather freely with high subject- and speaker-oriented adverbs. This is shown below first for evaluative adverbs \ref{evaluativejap}, in a direct translation of the example in \citet[214]{cinque:2020:relcl}, and then for evidential adverbs \ref{evidentialjap}. As the brackets indicate, the adverbs are understood as part of the relative clause.\footnote{Finding evidential adverbs can prove difficult, since Japanese speakers would rely on modal suffixes such as \textit{rashi-i} `seem' to express the speaker's opinion. I thank C. Takeuchi for providing valuable judgments and examples of relevant adverbs.}

\exg. [\textbf{Zannennagara/Oshikumo}\textsubscript{evaluative} sensō-de ushinaw-are-ta] eiga-wa ō-i.\\
unfortunately/regrettably war-\textsc{loc} lose-\textsc{pass-pst} movie-\textsc{top} many-\textsc{pred.cop.prs}\\
`The films [unfortunately/regrettably lost during the war] are many.’ \label{evaluativejap}

\exg. [\textbf{Osoraku/Dōyara/Tabun}\textsubscript{evidential} sensō-de ushinaw-are-ta] eiga-wa ō-i.\\
possibly/apparently/probably war-\textsc{loc} lose-\textsc{pass-pst} movie-\textsc{top} many-\textsc{pred.cop.prs}\\
`The films [possibly/apparently/probably lost during the war] are many.’ \label{evidentialjap}

Such constructions can also be constructed for relative clauses with an adjectival nucleus, as shown with the evaluative adverb \textit{saiwaini-mo} `fortunately’ in \ref{shinpai}.

\exg. [Shinpai-shi-ta kedo, \textbf{saiwaini-mo}\textsubscript{evaluative} genki-d-at-ta] josei-wa Hanako d-at-ta.\\
worry-do-\textsc{pst} but fortunately-\textsc{emp} healthy-\textsc{pred-cop-pst} woman-\textsc{top} Hanako \textsc{pred-cop-pst}\\
`The woman [I was worried about but who fortunately was fine] was Hanako.’ \label{shinpai}

These \is{CP adverb}data points in comparison with the ones illustrated above for \is{English} English, \is{French} French and \is{Italian} Italian lead to the same conclusion: Relative clauses in Japanese are also at least IPs, but on the other hand, it is not clear whether they can be bigger or not.

\subsubsection{Exponence of an overt subject}

Another indication of IP status is the availability of a subject position, which is standardly assumed to be Spec, IP (\citealt{chomsky:1986:barriers, koopman:1991:the}, see \citealt{mccloskey:1997:subject}). It was already mentioned that prenominal-RC languages generally exhibit an overt subject inside relative clauses and many relevant examples were shown in Japanese in Section \ref{internalstructure}. This is thus another indication of IP-status, although, in principle, nothing is said about a potential higher CP layer (although see \citealt{fukui:1986:theory} and subsequent work for an alternative analysis for Japanese). I give a verbal relative clause with filled Spec, IP in example \ref{kare-ga kat-ta hon} and its corresponding tree in \figref{RC kare-ga kat-ta hon tree} for illustration.

\exg. kare-ga kat-ta hon\\
he-\textsc{nom} buy-\textsc{pst} book\\
`the book (which) he bought' \label{kare-ga kat-ta hon}

\begin{figure}
\centering
\begin{tikzpicture}[baseline=0pt, remember picture]
\tikzset{level distance=25pt, sibling distance=5pt}
 \Tree [.DP\textsubscript{internal} {} [.D′\textsubscript{internal} [.FP\textsubscript{indirect} [.CP \node (op1) {}; [.C′ [.IP [.DP \edge[roof]; kare-ga ] [.I′ [.VP {} [.V′ \node (op2) {}; [.V kat ] ] ] [.I ta ] ] ] C ] ] [.F′\textsubscript{indirect} [.... [.NP \edge[roof]; hon ] {} ] F ] ] D ] ]
 \end{tikzpicture}
\caption{Relative clause with overt subject \textit{kare-ga kat-ta hon} `the book (which) he bought' (= \ref{kare-ga kat-ta hon})}
\label{RC kare-ga kat-ta hon tree}
\end{figure}

\subsection{Why Japanese relative clauses are not only IPs}

The question is whether there is a way to prove that relative clauses in Japanese are \textit{only} IPs. Below, I evaluate arguments from the literature and eventually falsify them. These include the analysis in \citet{cinque:2020:relcl} for an unequivocal IP status of relative clauses in \is{Mandarin} Mandarin Chinese and arguments from \citet{murasugi:1991:noun}.

\subsubsection{Individual-level vs. stage-level}

\citet[§ 3.6]{cinque:2020:relcl} argues that relative clauses in \is{Mandarin} Mandarin Chinese, which also only exhibit one surface type, are IPs. He bases his analysis on the alleged adherence to ordering restrictions between \is{stage/individual-level}stage-level (transitory) and individual-level (non-transitory) relative clauses presented in \citet{larson:2004:order} (and \citealt{gobbo:2005:chinese}) for various prenominal RC-languages, among them Mandarin and Japanese. The authors argue that when these types of relative clauses co-occur, the individual-level relative clause has to occur closer to the head noun than the stage-level relative clause \ref{s-i}. However, when two relative clauses with the same interpretation co-occur, \ref{i-i} and \ref{s-s}, the order is free. First, some of the relevant examples in Mandarin are \is{stage/individual-level}given (\citealt[104]{larson:2004:order}, taken from \citealt[298]{gobbo:2005:chinese}).\footnote{Recall that verbs and nouns in this language are always connected to the head noun via the element \textsc{de}. As presented in Section \ref{cartography Cinque}, this element is an indication of indirect modification. See also Footnote \ref{demandarin} in that chapter for various analyses, and Section \ref{internalstructure} for why it could be an instance of \textsc{mod}.}

\largerpage[2]
\ex. \label{ex1} \emph{stage-level $\succ$ individual-level} \label{s-i}
\ag. [Wo zuotian kanjian de]\textsubscript{s-level} [xihuan qu yinyuehui de]\textsubscript{i-level} ren shi Lisi. (Mandarin)\\
I yesterday meet \textsc{de} like go concert \textsc{de} person \textsc{cop} Lisi\\
\bg. * [Xihuan qu yinyuehui de]\textsubscript{i-level} [wo zuotian kanjian]\textsubscript{s-level} de ren shi Lisi.\\
like go concert \textsc{de} I yesterday meet \textsc{de} person \textsc{cop} Lisi\\
`The person who I met yesterday who likes to go to concerts is Lisi.'

\ex. \label{ex2} \emph{individual-level + individual-level} \label{i-i}
\ag. [Hui shuo Yidaliyu de]\textsubscript{i-level} [xihuan qu yinyuehui de]\textsubscript{i-level} ren shi Lisi. (Mandarin)\\
can speak Italian \textsc{de} like go concert \textsc{de} person \textsc{cop} Lisi\\
\bg. [Xihuan qu yinyuehui de]\textsubscript{i-level} [Hui shuo Yidaliyu de]\textsubscript{i-level} ren shi Lisi.\\
like go concert \textsc{de} can speak Italian \textsc{de} person \textsc{cop} Lisi\\
`The person who can speak Italian who likes to go to concerts is Lisi.'

\ex. \label{ex3} \emph{stage-level + stage-level} \label{s-s}
\ag. [Cong Yidali huilai de]\textsubscript{s-level} [wo zuotian kanjian de]\textsubscript{s-level} ren shi Lisi. (Mandarin)\\
from Italy return \textsc{de} l yesterday meet \textsc{de} person \textsc{cop} Lisi\\
\bg. [Wo zuotian kanjian de]\textsubscript{s-level} [cong Yidali huilai de]\textsubscript{s-level} ren shi Lisi.\\
l yesterday meet \textsc{de} from Italy return \textsc{de} person \textsc{cop} Lisi\\
`The person who returned from Italy who I met yesterday is Lisi.'

Crucially, \is{Mandarin} Mandarin relative clauses share this \is{stage/individual-level}restriction with full, but not \is{reduced relative clause}reduced, relative clauses in \is{English} English. First, observe the examples of postnominal relative clauses of different kinds of transitoriness in English below. Here, no ordering restrictions apply (see also \citealt{baker:2003:verbal}).

\ex. \a. the person [who I met]\textsubscript{s-level} [who smokes]\textsubscript{i-level}
\b. the person [who smokes]\textsubscript{i-level} [who I met]\textsubscript{s-level} \\ \null \hfill \citep[239]{cinque:2020:relcl}

Reduced relative clauses, on the other hand, do uphold such ordering preferences, as seen in the following example. \is{English}

\ex. the Wednesday Thursday lecture\\
1. stage-level $\succ$ individual-level\\
2. *individual-level $\succ$ stage-level \null \hfill \citep[112]{larson:2004:order} \label{wednesday-thursday}

In \ref{wednesday-thursday}, the first modifier must have a deictic, \is{stage/individual-level}stage-level, interpretation and refer to a lecture on a particular Wednesday, whereas the second modifier must have a generic, individual-level, interpretation and denote the usual lecture on Thursdays \citep[239]{cinque:2020:relcl}. The resulting DP denotes a lesson which usually takes place on Thursdays, but happens to take place on a Wednesday as an exception. From these facts, \citet{cinque:2020:relcl} deduces that relative clauses in Mandarin cannot be of the same type as full relative clauses, but pattern with \is{reduced relative clause}reduced relative clauses in English. Therefore, they are presumably of size IP.

Apart from the fact that it is not clear whether these observations can be extended to all relative clauses in Mandarin, things are more complicated. Direct counterarguments come from \citet{kim:2019:syntax}, who acknowledges the empirical adequacy of the \is{stage/individual-level}examples in \citet{larson:2004:order} for \is{Korean} Korean, but puts forth examples in which the order \textit{individual-level} $\succ$ \textit{stage-level} is actually preferred over the order \textit{stage-level} $\succ$ \textit{individual-level}, \ref{i-skorean}, and examples in which ordering restrictions are observed among the same type of \is{non-restrictive relative clause}relative clauses, as in \ref{s-skorean}.\footnote{These observations are explained in \citet{kim:2019:syntax} via a constraint-based account, which predicts obligatory movement of relative clauses with more familiar topics to the discourse participants (the speaker singing a song or remembering something) to a sentence-initial position, presumably at PF. The surface order is arguably the result of an output filter \citep{prince:1993:optimality} which ranks the best order based on various semantic and pragmatic constraints \citep[§ 5]{kim:2019:syntax}, those being \textit{episodic event description}, \textit{earlier event time} and \textit{higher familiarity to interlocutors} \citep[202]{kim:2019:syntax}. She establishes the hierarchy \textit{supplementary} $\succ$ \textit{non-restrictive} $\succ$ \textit{restrictive} \citep[120]{kim:2019:syntax}. Importantly, such constraints bear no relation to the interpretation as stage or individual-level relative clause \citep[239–240]{kim:2019:syntax} or to any difference between full-reduced or CP-IP. See also \citet{kotowski:2016:adj}, who did not find a significant tendency in a questionnaire-study on the order between stage- and individual-level adjectives in German.} \is{German}

\ex. \emph{individual-level $\succ$ stage-level as preferred order in Korean \citep[85]{kim:2019:syntax}} \label{i-skorean}
\ag. [Mina-ka phyengso culkye-pwuru-n-un,]\textsubscript{i-level} [chilsipnyentay-ey ywuhaynghay-ss-te-n]\textsubscript{s-level} noray-nun leylipti-i-Ø-ta.\\
Mina-\textsc{nom} usually enjoy-singing-\textsc{ipfv-rel} 1970s-\textsc{loc} be.popular-\textsc{ant-rtro-rel} song-\textsc{top} let.it.be-\textsc{cop-prs-decl}\\
\bg. ?[Chilsipnyentay-ey ywuhaynghay-ss-te-n,]\textsubscript{s-level} [Mina-ka phyengso culkye-pwuru-n-un]\textsubscript{i-level} noray-nun leylipti-i-Ø-ta.\\
 1970s-\textsc{loc} be.popular-\textsc{ant-rtro-rel} Mina-\textsc{nom} usually enjoy-singing-\textsc{ipfv-rel} song-\textsc{top} let.it.be-\textsc{cop-prs-decl}\\
 `The song that Mina usually enjoys singing that was popular in the 1970s is ``Let it be".’

\newpage
\ex. \emph{Two stage-level relative clauses, no free ordering in Korean \citep[86]{kim:2019:syntax}} \label{s-skorean}
\ag. [Mina-ka kacang cal kiekha-n-un,]\textsubscript{s-level} [kwusipnyentay-chopan-ey iss-ess-te-n]\textsubscript{s-level} il-un emma-hako-uy yehayang-i-Ø-ta.\\
Mina-\textsc{nom} most well remember-\textsc{ipfv-rel} 1990s-early-\textsc{loc} exist-\textsc{ant-rtro-rel} event-\textsc{top} Mom-with-\textsc{gen} trip-\textsc{cop-prs-decl}\\
\bg. *[Kwusipnyentay-chopan-ey iss-ess-te-n,]\textsubscript{s-level} [Mina-ka kacang cal kiekha-n-un]\textsubscript{s-level} il-un emma-hako-uy yehayang-i-Ø-ta.\\
1990s-early-\textsc{loc} exist-\textsc{ant-rtro-rel} Mina-\textsc{nom} most well remember-\textsc{ipfv-rel} event-\textsc{top} Mom-with-\textsc{gen} trip-\textsc{cop-prs-decl}\\
`The event that Mina remembers the best which happened in the early 1990s is the trip with her Mom.'

These \is{stage/individual-level}data points are problematic for \citet{larson:2004:order} and by extension also for \citet{cinque:2020:relcl} but, as \citet{kim:2019:syntax} admits, it might be far-fetched to draw any immediate fundamental conclusions from them.

Coming now to Japanese, \citet[102]{larson:2004:order} report identical judgments regarding the order of stage- and individual-level relative clauses as given below.

\ex. \label{sijapa} \emph{stage-level $\succ$ individual-level in Japanese after \citet[102]{larson:2004:order}}
\ag. [Watashi-ga kinō at-ta]\textsubscript{s-level} [tabako-wo su-u]\textsubscript{i-level} hito-wa Tanaka-san-de-s-u.\\
I-\textsc{nom} yesterday meet-\textsc{pst} tobacco-\textsc{acc} inhale-\textsc{prs} person-\textsc{top} Tanaka-Mr.-\textsc{pred-cop.pol-prs}\\
\bg. *[Tabako-wo su-u]\textsubscript{i-level} [watashi-ga kinō at-ta]\textsubscript{s-level} hito-wa Tanaka-san-de-s-u.\\
tobacco-\textsc{acc} inhale-\textsc{prs} I-\textsc{nom} yesterday meet-\textsc{pst} person-\textsc{top} Tanaka-Mr.-\textsc{pred-cop.pol-prs}\\
`The person who I met yesterday who smokes is Mr./Mrs. Tanaka.'

\ex. \emph{Two individual-level relative clauses in Japanese according to \citet[102]{larson:2004:order}}
\ag. [Tabako-wo su-u]\textsubscript{i-level} [sake-wo no-mu]\textsubscript{i-level} hito-wa Tanaka-san-de-s-u.\\
tobacco-\textsc{acc} inhale-\textsc{prs} alcohol-\textsc{acc} drink-\textsc{prs} person-\textsc{top} Tanaka-Mr.-\textsc{pred-cop.pol-prs}\\
\bg. [Sake-wo no-mu]\textsubscript{i-level} [tabako-wo su-u]\textsubscript{i-level} hito-wa Tanaka-san-de-s-u.\\
alcohol-\textsc{acc} drink-\textsc{prs} tobacco-\textsc{acc} inhale-\textsc{prs} person-\textsc{top} Tanaka-Mr.-\textsc{pred-cop.pol-prs}\\
`The person who smokes who drinks alcohol is Mr./Mrs. Tanaka.'

\ex. \emph{Two stage-level relative clauses in Japanese according to \citet[102]{larson:2004:order}}
\ag. [Watashi-ga kinō at-ta]\textsubscript{s-level} [sake-wo non-de-i-ta]\textsubscript{s-level} hito-wa Tanaka-san-de-s-u.\\
I-\textsc{nom} yesterday meet-\textsc{pst} alcohol-\textsc{acc} drink-\textsc{ger}-be-\textsc{pst} person-\textsc{top} Tanaka-Mr.-\textsc{pred-cop.pol-prs}\\
\bg. [Sake-wo non-de-i-ta]\textsubscript{s-level} [watashi-ga kinō at-ta]\textsubscript{s-level} hito-wa Tanaka-san-de-s-u.\\
alcohol-\textsc{acc} drink-\textsc{ger}-be-\textsc{pst} I-\textsc{nom} yesterday meet-\textsc{pst} person-\textsc{top} Tanaka-Mr.-\textsc{pred-cop.pol-prs}\\
`The person who I met yesterday who was drinking alcohol is Mr./Mrs. Tanaka.'

Speakers I asked judged the example in \ref{sijapa} as more natural if the \is{stage/individual-level}stage-level relative clause preceded the individual-level relative clause. However, the reversed order \textit{smokes} $\succ$ \textit{likes to go to concerts} seems not to be ungrammatical, contrary to what is argued in \citet{larson:2004:order}. Regarding the counterexamples reported in \citet{kim:2019:syntax}, I have attempted as best as possible to give Japanese translations of her examples and discuss them with consultants as out-of-the-blue examples in order to see if her observations also hold in this language. As it turns out, my consultants did not detect any ordering restrictions and judged both orders as equally natural, exemplified in \ref{free order s- and l- in Japanese}.

\ex. \label{free order s- and l- in Japanese} \emph{Free ordering between stage-level and individual-level relative clause in Japanese translated from Korean from \citet[85]{kim:2019:syntax}}
\ag. [Mīna-ga itsumo uta-u,]\textsubscript{i-level} [1970-nen-ni ninki-dat-t-ta]\textsubscript{s-level} kyoku-wa Letitbe-d-a.\\
Mina-\textsc{nom} always sing-\textsc{prs} 1970s-year-\textsc{loc} be.popular-\textsc{pred-cop-pst} song-\textsc{top} let.it.be-\textsc{pred-cop.prs}\\
\bg. [1970-nen-ni ninki-dat-t-ta,]\textsubscript{s-level} [Mīna-ga itsumo uta-u]\textsubscript{i-level} kyoku-wa Letitbe d-a.\\
1970s-year-\textsc{loc} be.popular-\textsc{pred-cop-pst} Mina-\textsc{nom} always sing-\textsc{prs} song-\textsc{top} let.it.be \textsc{pred-cop.prs}\\
`The song that Mina always sings that was popular in the 1970's is ``Let it be".’

This discussion has shown first of all that a deduction of IP status on the grounds of ordering restrictions among \is{stage/individual-level}stage- and individual-level relative clauses is hard to establish. It also emphasizes free order of modifiers of the same type in Japanese and corroborates the findings from previous chapters. A final disclaimer: I am aware that a deeper analysis of the order of relative clauses in Japanese needs to be conducted than can be granted here. Speakers generally had a hard time determining which order among the two relative clauses I presented was best. There might be factors ameliorating one order over the other that can be brought about by context.\footnote{This was also pointed out to me by Min-Joo Kim, to whom I am grateful for discussing the relevant examples.} Aside from syntactic and phonological/prosodic complexity, which will be addressed in the next section, I could not find determining factors to ameliorate one order over the other, so this is an endeavor I leave for future research.

\subsubsection{Reason and manner clauses} \label{reason and manner}

An advocate for the IP status of relative clauses in Japanese is \citet{murasugi:1991:noun}. One argument she puts forth is the lack of long distance dependencies in \textit{reason} and \textit{manner} relative clauses. See the following examples. \is{reason and manner clauses}

\exg. [Marī-ga [Jon-ga kaet-ta to] omot-te-i-ru] riyū \label{reason}\\
Mary-\textsc{nom} John-\textsc{nom} return-\textsc{pst} \textsc{comp} think-\textsc{ger}-be-\textsc{prs} reason\\
1. $\approx$ `the reason Mary thinks that John returned’\\
2. $\neq$ `the reason why John returned according to what Mary thinks' \\ \null \hfill \citep[185]{murasugi:1991:noun}


\exg. [Jon-ga [Marī-ga mondai-wo toi-ta to] it-ta] hōhō \label{manner}\\
John-\textsc{nom} Mary-\textsc{nom} problem-\textsc{acc} solve-\textsc{pst} \textsc{comp} say-\textsc{pst} way\\
1. $\approx$ `the way Mary said that John solved the problem’\\
2. $\neq$ `the way John solved the problem according to what Mary said’ \\ \null \hfill \citep[133]{murasugi:1991:noun}

In \ref{reason}, the speaker can only refer to the reason why Mary thinks that John returned, not to the actual reason of his return, according to what Mary thinks. Similarly, in \ref{manner}, the speaker can only refer to the way Mary said that John solved the problem, not to the way he actually solved it. This contrasts with relative clauses in which locative and temporal nouns are relativized. These relative clauses in fact allow two readings, a higher reading and a lower reading. This shows that they do allow long-distance relativization \citep[133–134]{murasugi:1991:noun}. \is{reason and manner clauses} \is{long distance relativization}

\exg. [Marī-ga [Jon-ga sentaku-shi-ta to] it-ta] hi \label{hi}\\
Mary-\textsc{nom} John-\textsc{nom} washing-do-\textsc{pst} \textsc{comp} say-\textsc{pst} day\\
1. $\approx$ `the day on which Mary said that John did laundry’\\
2. $\approx$ `the day on which John did the laundry according to what Mary said’ \\ \null \hfill \citep[134]{murasugi:1991:noun}

\exg. [Jon-ga [Marī-ga mondai-wo toi-ta to] it-ta] tokoro\\
John-\textsc{nom} Mary-\textsc{nom} problem-\textsc{acc} solve-\textsc{pst} \textsc{comp} say-\textsc{pst} place\\
1. $\approx$ `the place where Mary said that John solved the problem’\\
2. $\approx$ `the place where John solved the problem according to what Mary said’ \\ \null \hfill \citep[134]{murasugi:1991:noun}

Following the framework of \citet{lasnik:1992:move}, \citet[§ 3.4]{murasugi:1991:noun} assumes a \textit{pro}-account for relative clauses and that relativization is only possible for arguments, never adjuncts, to which \textit{reason}- and \textit{manner}-nouns belong, according to her, and which must be clause-bound. This contrasts them with \textit{time}- and \textit{place}-nouns, which are arguments. In favor of \textit{pro}, she dismisses operator-movement-based accounts. As she assumes Japanese relative clauses to be IPs, no projection (Spec, CP) is available to host an operator in the first place.

Abstracting away from the fact that \textit{pro}-based accounts have been dismissed on several grounds \citep[36–37]{salzmann:2017:reconstruction}, rather than having to assume an IP status for such relative clauses, the problem seems to be the peculiar behavior of \textit{reason} and \textit{manner}-relative clauses, as also argued in \citet{miyamoto:2017:RC}. In this light, note another contrast between \textit{reason}/\textit{manner}-relative clauses, on the one hand, and locative/temporal relative clauses on the other. The latter can appear as major subjects in a \is{major subject construction}major subject construction (\citealt{kuroda:1992:japanese} and already \citealt{kuno:1973:structure}; cf. Section \ref{majorsubject}). \is{operator}

\exg. Sono hi-ga [Hanako-ga [Tarō-ga hashit-ta to] it-ta].\\
that day-\textsc{nom} Hanako-\textsc{nom} Taro-\textsc{nom} run-\textsc{pst} \textsc{comp} say-\textsc{pst}\\
`The day is such that Hanako said that Taro ran.' \\ \null \hfill \citep[628]{miyamoto:2017:RC} \label{sono hi}

\newpage

As in \ref{hi}, this sentence can refer to the time that Hanako said that Taro ran, or the time that he ran according to Hanako. The corresponding (major subject) construction with the noun \textit{riyū} `reason' is not only unavailable under the desired reading (`the reason Taro ran according to Hanako’), it is entirely ungrammatical.

\exg. *Sono riyū-ga [Hanako-ga [Tarō-ga hashit-ta to] it-ta].\\
that reason-\textsc{nom} Hanako-\textsc{nom} Taro-\textsc{nom} run-\textsc{pst} \textsc{comp} say-\textsc{pst}\\
(`The reason is such that Hanako said that Taro ran.’) \\ \null \hfill \citep[628]{miyamoto:2017:RC}

\citet{miyamoto:2017:RC} draws the conclusion that relative clauses as in \ref{hi} are derived via \is{operator}operator-movement out of \is{major subject construction}major subject constructions. Since the underlying structure is not available for manner nouns in the first place, the relative clause in \ref{reason} does not yield the desired long-distance reading. The conclusion to draw from these facts is that the unavailability of long-distance dependencies in \textit{reason}- and \textit{manner}-relative clauses in Japanese is not related to the size of the relative clause, but to the nature of the relativized noun and the possible derivation of Japanese relative clauses in Japanese out of a major subject position (see Section \ref{majorsubject}).

\subsection{Relative clauses are not of different sizes} \label{NGC}

In this subsection, accounts will be reviewed according to which relative clauses in Japanese come in different sizes. These are mostly based on the phenomenon \is{nominative-genitive-conversion}\textit{Nominative Genitive Conversion} (NGC, also \textit{ga-no-conversion}). I will review some of the arguments from the literature and then report on how my consultants judged the data as well as some additional examples. It will be argued that not enough evidence can be put forward to justify an analysis along such lines.

\subsubsection{Nominative-Genitive-Conversion}

Japanese can mark subjects within a relative clause with the nominative case particle, but alternatively also with -\textit{no}, the genitive case particle.\footnote{Genitive case marking of the subject in a relative clause is a phenomenon not unique to Japanese, but also exhibited in other languages, for example \is{Turkish}Turkish, \is{Mongolian}Mongolian and \is{Dagur}Dagur \citep{krause:2001:rrc}.} As exemplified in \ref{no-ga}, both options -- nominative and genitive -- are available and this choice does not affect the meaning. Note that I adopt the gloss \textsc{gen} for this instance of -\textit{no} here as it cannot be an instance of \is{MOD@\textsc{mod}}\textsc{mod} in accordance with the refinement rule given in Section \ref{internalstructure}, i.e. since it does not follow a modifier of the noun \textit{per se}, but rather the subject of the relative clause.  

\exg. kare-ga / kare-no kat-ta hon\\
he-\textsc{nom} / he-\textsc{gen} buy-\textsc{pst} book\\
`the book (which) he bought' \label{no-ga}

In the 1970s, \citet{harada:1971:ga} argued that the different markers are purely a result of idiolectal differences. He sparked a discussion on the linguistic dimension of this phenomenon, which was picked up in the generative literature from the 90s onwards. In fact, although the two respective markers -- -\textit{ga} for nominative and -\textit{no} for genitive -- seem interchangeable in most scenarios, various restrictions are observable with genitive-marked subject relative clauses. The following have been pointed out in the \is{nominative-genitive-conversion}literature \citep{watanabe:1996:nominative, hiraiwa:2001:on, miyagawa:2011:genitive, miyagawa:2012:case, miyagawa:2017:agreement, ochi:2004:ga, ochi:2017:ga, ochi:2021:feature, akaso:2011:on, nambu:2014:an, nambu:2023:experimental}.

\ex. \a. A restriction on CP adverbs
\b. Intervention effects: a restriction on intervening adjuncts
\c. Focus is impossible with genitive marked subjects.
\d. Genitive marked subjects give rise to potential Principle B violations.
\e. Genitive marked subjects can give rise to scope ambiguities.
\f. Genitive marked subjects are not compatible with accusative marked objects. (Transitivity restriction)

There are two main approaches to this phenomenon, which both relate the licensing of case to the size of the relative clause: the `C-licensing approach’, which assumes a uniform analysis and that a CP is always present, which is where case is licensed (\citealt{watanabe:1996:nominative} and \citealt{hiraiwa:2001:multiple, hiraiwa:2001:on}), and the `D-licensing approach' \citep{miyagawa:2011:genitive, miyagawa:2012:case, miyagawa:2017:agreement}, which postulates different analyses of nominative- and genitive-marked subject relative clauses, namely as CPs and TPs. I will concentrate mainly on the latter account, since it correlates with the two different sizes of relative clauses. After outlining the main idea, focusing on the first three points mentioned above, I will point out two things. First, these arguments are not strong enough to convincingly argue for a different status of Japanese relative clauses. Second, they can also be explained in different ways, without proposing two different analyses. For a complete overview of all the patterns, the reader is referred to the relevant literature and especially \citet{ochi:2017:ga}.

\subsubsection{The idea}

Concretely, \citet{miyagawa:2011:genitive, miyagawa:2012:case, miyagawa:2017:agreement} treats genitive marking of subjects inside a relative clause as a version of Exceptional Case Marking (ECM, see \citealt{chomsky:1986:barriers, chomsky:2001:derivation}). The reason why the subject is marked by genitive and not by nominative is because T is defective \citep[1271]{miyagawa:2011:genitive}. It is defective because it does not have a higher node, namely C, from which it would receive case. Therefore, it cannot be a case-assigner, and since the subject cannot receive case from T, it must look for another case-assigning category in a higher node. It finds this category in D, which is name-bearing for the approach \citep[1271]{miyagawa:2011:genitive}. The fact that the relative clause is TP is thus a prerequisite for genitive case-marking. This is illustrated in \figref{Rc with genitive miyagawa}.

\begin{figure}
\centering
\begin{tikzpicture}[baseline=-0pt]
\tikzset{level distance=25pt, sibling distance=5pt}
\Tree [.DP {} [.D′ [.TP [.vP \textsc{sub-gen} [.v′ VP v ] ] [.T ] ] [.D ]]]
\end{tikzpicture}
\caption{Structure for relative clauses with genitive subject \citep[147]{miyagawa:2017:agreement}}
\label{Rc with genitive miyagawa}
\end{figure}

Nominative subjects, on the other hand, do receive their case from \is{nominative-genitive-conversion}T. As mentioned above, T needs to receive case from somewhere in the first place to later assign it to the subject. This means it must be selected by a higher CP. This T, which has received case from C, is not defective, but properly assigns nominative case to the subject \citep[1271]{miyagawa:2011:genitive}.

\citet[1271–1272]{miyagawa:2011:genitive} assumes that C is a phase and, intervening between T and D, stops D from assigning case to the subject, which is why the nominative subject cannot get its case from D in this scenario. This means that in Miyagawa's model, a CP layer has to be present when a subject appears in nominative case. This alternative option is given in \figref{Rc with nominative miyagawa}.

\begin{figure}
\centering
\begin{tikzpicture}[baseline=-0pt]
\tikzset{level distance=25pt, sibling distance=5pt}
\Tree [.DP {} [.D′ [.CP \node (IP1) {}; [.C′ [.TP \textsc{sub-nom} [.T′ vP T ] ] [.C ] ] ] [.D ]]]
\end{tikzpicture}
\caption{Structure for relative clauses with nominative subject \citep[147]{miyagawa:2017:agreement}}
\label{Rc with nominative miyagawa}
\end{figure}

\subsubsection{CP-elements and adjacency}

A central argument for Miyagawa's claim that genitive-marked relative clauses lack a CP layer comes from the alleged incompatibility between genitive-marked relative clauses and \is{CP adverb}CP-elements. Recall from Section \ref{cinque adj} that high speaker- and subject-oriented CP adverbs are freely compatible with Japanese relative clauses. The relevant example is repeated for convenience below (= \ref{evaluativejap}).

\exg. [\textbf{Zannennagara/Oshikumo}\textsubscript{evaluative} sensō-de ushinaw-are-ta] eiga-wa ō-i.\\
unfortunately/regrettably war-\textsc{loc} lose-\textsc{pass-pst} movie-\textsc{top} many-\textsc{pred.cop.prs}\\
`The films unfortunately/regrettably lost during the war are many.’

\citet[1273]{miyagawa:2011:genitive} argues that modal TP adverbs such as \textit{kitto} `certainly' are possible in relative clauses where the \is{nominative-genitive-conversion}subject is marked by nominative, but also in those, where it is marked by genitive, as shown in \ref{kitto}. On the other hand, he argues that \is{CP adverb}CP adverbs are only possible with nominative-marked, never with genitive-marked, subject relative clauses, which is because no CP layer is present in the latter case. He shows this using the evaluative adverb \textit{saiwaini} `fortunately' \ref{saiwai}. 

\exg. [kitto Tarō-ga / -no yon-da] hon\\
certainly Taro-\textsc{nom} / \textsc{gen} read-\textsc{pst} book\\
`the book [that Taro certainly read]’ \null \hfill \citep[1273]{miyagawa:2011:genitive} \label{kitto}

\exg. [saiwaini Tarō-ga / *-no yon-da] hon\\
fortunately Taro-\textsc{nom} / \textsc{gen} read-\textsc{pst} book\\
`the book [that Taro fortunately read]’ \null \hfill \citep[1273]{miyagawa:2011:genitive} \label{saiwai}

Recall that I have argued that these tests merely show that relative clauses are IP/TP, but that they cannot readily be used as indication whether they are IPs or CPs. A further problem with the example \ref{saiwai} is that the \is{CP adverb}adverb \textit{saiwaini} `fortunately' was not accepted by any of my consultants in the first place and arguably should be replaced by the more natural \textit{saiwaini-mo} as given below.

\exg. [saiwaini-mo Tarō-ga / Tarō-no yon-da] hon\\
fortunately-\textsc{emp} Taro-\textsc{nom} / Taro-\textsc{gen} read-\textsc{pst} book\\
`the book [that Taro fortunately read’] \label{saiwainimo}

Comparing \ref{saiwainimo} with \ref{kitto}, the speakers I consulted either judged both DPs as acceptable, regardless of whether the nominative or genitive case marker occurred, or they judged only the \is{nominative-genitive-conversion}nominative-marked relative clauses acceptable, but crucially for \textit{both} examples. This means they did not provide varying judgements depending on the two adverbs, and if they showed preference it was a strong preference for nominative in the first place.

A general problem with such examples is that neither context nor matrix clause is provided. Attempting to embed this relative clause in a matrix clause, however, runs the risk that the adverb is understood as referring to the matrix verb, which follows from its status as a CP adverb. In the following sentence in \ref{wallet}, the reference of the adverb is semantically and pragmatically controlled for. A reading where \textit{saiwaini-mo} refers to the matrix clause -- and also to the second subordinate clause -- is not syntactically excluded, but pragmatically marked, since it would be strange to associate `fortunately' with losing something or being in trouble, compared to a fortunate finding of ones wallet. Hence, this adverb should attach at the CP-layer of the embedded relative clause. This is indicated by the brackets. The important part is that, according to the speakers I consulted, both versions, with -\textit{ga} and with -\textit{no}, are equally acceptable.\footnote{I am deeply grateful to Ken Hiraiwa for providing this sentence and to him as well as C. Takeuchi for discussion of these patterns.}

\exg. [[Saiwaini-mo Tarō-ga / Tarō-no mitsuke-te-kure-ta] saifu-wo mata nakushi-te-shimat-ta node,] komat-te-i-ru.\\
fortunately-\textsc{emp} Tarō-\textsc{nom} / Tarō-\textsc{gen} find-\textsc{ger}-do.favor-\textsc{pst} wallet-\textsc{acc} again lose-\textsc{ger}-suffer.from-\textsc{pst} because be.in.trouble-\textsc{ger}-be-\textsc{prs}\\
`[Because I have lost the wallet again [that Taro fortunately has found for me]] I am in trouble.' \label{wallet}

The fact that both \is{nominative-genitive-conversion}case markers are possible seems surprising and is problematic for Miyagawa's analysis. This argument relates directly to another piece of data already mentioned in \citet{harada:1971:ga}, namely the degradation of genitive marked subject relative clauses with intervening material between the verb and the subject as in the following example.

\exg. [kodomo-tachi-ga / *-no minnade ikioi-yoku kake-nobot-ta] kaidan\\
kid-\textsc{pl}-\textsc{nom} / \textsc{gen} everyone vigorous-\textsc{adv} run-climb-\textsc{pst} stairs\\
`the stairs [which the kids vigorously ran up all together]’ \\ \null \hfill \citep[30]{harada:1971:ga} \label{stairs}

The judgments of my consultants with regard to \ref{stairs} again suggest that genitive is rather marked when additional elements intervene, or in other words, when the subject is not adjacent to the noun.\footnote{The intervention patterns are analyzed in \citet[§ 4.2]{miyagawa:2011:genitive} along an economy of movement analysis. The genitive subject (that has received case from D) remains in Spec, vP and does not raise to Spec, TP since T is defective and thus has no EPP feature. Therefore, the genitive marked subject cannot appear high up in the clause because there is no motivation for it to move there. Low (VP) \is{CP adverb}adverbs, however, can be easily accommodated. It remains unclear what a notion such as \textit{economy} can refer to in a language such as Japanese in which scrambling is so productive.} This observation was backed up experimentally in a rating and a self-paced reading study in \citet{nambu:2014:an, nambu:2023:experimental}, who show a significant degradation of genitive-marking with intervening TP or CP (compared to VP) adjuncts, whereas acceptance remained stable for the nominative-marked subject condition. Interestingly, the authors also note in their investigation of reading times a significant degradation of genitive with intervening adjuncts in general. Besides CP and TP adjuncts, they also detected this for locative (more than temporal or manner) PPs, which are arguably in the VP domain \citep[15–23]{nambu:2023:experimental}. It seems as if there is a general intervention effect happening with genitive case, but it is questionable whether this is due to a smaller status of relative clauses containing a genitive-marked subject.\footnote{The authors suggest a processing-based account for locative PPs \citep[25]{nambu:2023:experimental}.}

An alternative piece of data is offered by \citet{ochi:2021:feature} and illustrated in \ref{ochiexample}. Here, a genitive-marked relative clause with three adjuncts appears. This is odd, but not unacceptable, and for some reason less odd than if the subject appeared after \textit{kinō} `yesterday' according to the speakers the author has consulted.\footnote{In his account, all Japanese relative clauses are IPs. The different factor is simply the subject, which can optionally move up to D where it receives genitive case, which it does in \ref{ochiexample}.}

\exg. (?)[kodomo-tachi-no kinō ōzei kyōshitsu-de sawai-da] koto\\
children-\textsc{pl-gen} yesterday many classroom-\textsc{loc} clamor-\textsc{pst} fact\\
`the fact [that many children clamored in the classroom yesterday]’ \\ \null\hfill \citep[326]{ochi:2021:feature} \label{ochiexample}

Generally speaking, there seems to be a tendency for speakers of Modern Japanese to prefer nominative over genitive in relative clauses, as already mentioned in \citet{harada:1971:ga}. The \is{CP adverb}study in \citet{nambu:2014:an, nambu:2023:experimental} can be interpreted along these lines as well, as they do show this tendency independent of the condition \textit{adjacent/non-adjacent}, as also mentioned by the authors. Therefore, all these facts could find a different \is{nominative-genitive-conversion}explanation.

\subsubsection{Focus and genitive} \label{focusrc}

The interaction between \is{focus}\textit{focus} and relative clauses in Japanese gives further important insights. Assuming a projection related to focus is part of the (extended) CP domain, the \is{left periphery}left periphery of the clause \citep{rizzi:1997:the}, predicts that focus is impossible inside genitive case-marked subject relative clauses, if they are not CPs, as such a layer would not be accessible. And this is exactly the analysis argued for in \citet{miyagawa:2017:agreement}.

Recall from Section \ref{externaldp} that focus is marked in Japanese either prosodically, or via focus-sensitive particles – the topic particle -\textit{wa}, which can express a focus function \ref{shanai-de}, and emphatic particles such as -\textit{dake} `only' or -\textit{sae} `even’ \ref{taroganondakusuri} \citep{tomioka:2009:contrastive, tomioka:2016:infor, vermeulen:2007:jap}.

\exg. [Nihongo-wo shanai-de-wa tsuka-u] gaikokujin\\
Japanese-\textsc{acc} office-\textsc{loc}-\textsc{foc} use-\textsc{prs} foreigner\\
`a foreigner [who uses Japanese only in the office]’ \label{shanai} \\ \null \hfill \citep[93]{akaso:2011:on} \label{shanai-de}

\exg. [Tarō-dake-ga non-da] kusuri\\
Taro-only-\textsc{nom} drink-\textsc{pst} medicine\\
`the medicine [(that) only Taro took]’ \null \hfill \citep[96]{akaso:2011:on} \label{taroganondakusuri}

Postponing a precise analysis of focus marking inside relative clauses to Section \ref{focusincp}, the crucial point is that those focus (-sensitive) particles are generally banned from genitive-marked subject relative clauses, as argued in \citet{akaso:2011:on} and \citet{miyagawa:2017:agreement}. Compare \ref{taroganondakusuri} with \ref{tarono}.

\exg.*[Tarō-dake-no non-da] kusuri \label{kusuri} \\
Taro-only-\textsc{gen} drink-\textsc{pst} medicine\\
(`the medicine [(that) only Taro took]’) \null \hfill \citep[96]{akaso:2011:on} \label{tarono}

Unlike many of the contrasts mentioned above, this contrast seems to be more systematic and my consultants had a clearer opinion and shared the judgments of the literature. In other words, the arguments presented in \citet{miyagawa:2017:agreement} seem to hold.

Nevertheless, there are some challenges to this observation. First, note that the genitive particle -\textit{no} and \is{focus}focus-sensitive particles such as -\textit{dake} `only' are not totally incompatible. The following examples show that they can co-occur. Note that the particle, in this case, follows an adverb \ref{kino} and a PP \ref{omote}.

\exg. [kinō-dake Tarō-no non-da] kusuri\\
yesterday-only Taro-\textsc{gen} drink-\textsc{pst} medicine\\
`the medicine [that Taro drank only yesterday]’ \null \hfill \citep[692]{ochi:2017:ga} \label{kino}

\exg. [omote-dake-ni nori-no tsui-te-i-ru] shīto\\
surface-only-\textsc{loc} glue-\textsc{gen} stick-\textsc{ger}-be-\textsc{prs} sheet\\
`a sheet [where glue is only stuck to the surface]’ \label{omote}

Furthermore, as shown in \citet{ochi:2017:ga}, in some Kyūshū dialects, nominative-genitive-conversion \is{nominative-genitive-conversion}has even been retained in main clauses. In these dialects, -\textit{dake} is infelicitous with genitive-marked subjects in general, thus also in main clauses, which are CPs. \is{Kyushu Japanese}

\exg. ??Hanako-dake-no furansugo-no hana-se-ru to yo. (Kyūshū J.)\\
Hanako-only-\textsc{gen} French-\textsc{gen} speak-\textsc{pot-prs} \textsc{comp} \textsc{sfp}\\
`(Listen), only Hanako can speak French!’ \\ \null \hfill \citep[692]{ochi:2017:ga}

These data points show that focus and genitive marking are not mutually exclusive. Nevertheless, the intuition seems to be strong that genitive-marked subject relative clauses resist focus marking. The question is if this is really due to a different size of the relative clause. In the following section, focus will be used as one indication why Japanese relative clauses must be CPs in general.

\subsubsection{Summary}

This subsection discussed arguments, mainly from \citet{miyagawa:2011:genitive, miyagawa:2012:case, miyagawa:2017:agreement} for relative clauses of different sizes based on the phenomenon of nominative-genitive-conversion. This discussion has shown that the data points of Miyagawa are either problematic or can be handled under a different explanation rather than resorting to the systematic distinction between CP and IP status. In other words, the assumption that relative clauses come in different sizes in Japanese does not withstand closer scrutiny.

\subsection{Japanese relative clauses are CPs} \label{kaplan}

Essentially, this leaves us with one option, namely the null hypothesis that all relative clauses in Japanese are CPs. I will provide some arguments here. One very convincing piece of evidence comes from the behavior of focused elements just viewed in Section \ref{focusrc} above. Other arguments come from historical evidence on the presence of complementizers. Finally, I will present the derivation of relative clauses along an operator-movement account, which gives further evidence for CP status.

\subsubsection{Focus and relative clauses in Japanese} \label{focusincp}

As outlined in the examples above, the possible occurrence of \is{focus}focused and topicalized elements inside relative clauses is a rather convincing piece of evidence for CP status. Adopting a \is{left periphery}left-periphery of the clause, a high area within the split CP (\citealt{rizzi:1997:the}; and see \citealt{grohmann:2003:prolific, laenzlinger:2011:elements}), which contains the projections FocP and (several) TopP among others (for instance \citealt{kiss:1998:identificational, kiss:2007:topic, aboh:2014:info}) offers a position for elements in (high) non-canonical positions. This is shown below where such a focused element in \is{Italian} Italian appears in a non-canonical position. Focused elements are displayed in capital letters.\footnote{See Section \ref{externaldp} for the left periphery in the DP; for focus within the DP, see Section \ref{focus,japanese}.}

\exg. Il TUO LIBRO ho letto (, non il suo). (Italian) \\
the your book have read {} not the his\\
`YOUR BOOK, I have read (, not his).’ \null \hfill \citep[286]{rizzi:1997:the}

As \citet{cinque:2020:relcl} points out, \is{reduced relative clause}reduced relative clauses in Italian are incompatible with focus marking, which suggests that they lack a CP alongside a higher FocP layer \citep[218–219]{cinque:2020:relcl}. \is{Italian}

\exg. *Il libro A MARIA dato (non a Carlo) era molto costoso. (Italian)\\
\textsc{det} book to Maria given not to Carlo was very expensive\\
(`The book given TO MARIA (not to Carlo) was very expensive.’) \\ \null \hfill \citep[219]{cinque:2020:relcl}

Above, I already presented examples in which relative clauses in Japanese contain \is{focus}focus-sensitive particles. Those can be attached to the subject, but they can also be used to put emphasis on adverbs \ref{kino} or PPs \ref{omote}. Below, to complement the picture, I give an example of an adjectival relative clause in which the subject, which is the possessee of the relativized head noun, is focused \ref{doresudake}. In this DP, it is \textit{only} the dress of the woman which is highlighted as beautiful, contrasting with other items, or possibly the woman herself.

\exg. doresu-dake-ga utsukushi-k-at-ta josei\\
dress-only-\textsc{nom} beautiful-\textsc{pred-cop-pst} woman\\
`the woman such that only her dress was beautiful' (Lit. `only the dress was beautiful woman’) \label{doresudake}

Such a contrastive focus, involving a \textit{set of alternatives} \citep{zimmermann:2011:focus}, can also be marked via the particle -\textit{wa}, which is typically used as the topic marker but can express focus \citep{vermeulen:2007:jap, vermeulen:2012:info}. In addition to \ref{shanai}, where this marker was added to a postposition, it can also attach to the subject of a relative clause and replace the nominative (or genitive) case marker \ref{ryori-wa}.

\exg. [Tarō-wa tabe-na-k-at-ta] ryōri-wa dore de-s-u ka?\\
Taro-\textsc{top} eat-\textsc{neg-pred-cop-pst} dish-\textsc{top} which \textsc{pred-cop.pol-prs} \textsc{q}\\
`Which one is the dish [that TARO did not eat]?’ \label{ryori-wa}

Moreover, -\textit{wa} inside a relative clause does not necessarily entail contrastive focus, but can also yield presentational focus \ref{wa presentational}.\footnote{This type of focus is less defined by an existing set of alternatives and its contrastive nature, but more by the intention of the speaker to ``express the most important part of the utterance, or what is new in the utterance" \citep[23]{krifka:2007:basic}.}

\exg. Kore-wa [mō waka-i hito-wa yoma-na-i] hon-de-s-u.\\
this-\textsc{top} any.longer young-\textsc{mod} person-\textsc{top} read-\textsc{neg-mod} book-\textsc{pred-cop.pol-prs}\\
`This is a book [such that it is young people who do not read it any longer].’ \label{wa presentational}

\begin{figure}
\fittable{
\begin{tikzpicture}[baseline=-0pt, remember picture]
\tikzset{level distance=25pt, sibling distance=5pt}
\Tree [.DP\textsubscript{internal} {} [.D′\textsubscript{internal} [.FP\textsubscript{indirect} [.TopP {} [.Top′ [.FocP [.DP \edge [roof]; \node (taro1) {Tarō-dake-ga}; ] [.Foc′ [.CP {} [.C′ [.IP \node (taro2) {$<$DP$>$}; [.I′ [.VP {} [.V′ {} [.V non ] ] ] [.I da ] ] ] C ]] Foc ] ] Top ] ] [.F′\textsubscript{indirect} [.... [.NP \edge[roof]; kusuri ] {} ] F ] ] D ] ]
\draw [color=blue][<->] (taro2) |- +(0,-2) -| (taro1);
\end{tikzpicture}
}
\caption{Structure of \ref{taroganondakusuri} with focus movement: \textit{Tarō dake-ga non-da kusuri} `the medicine only Taro took'}
\label{Rc with focus movement}
\end{figure}

\largerpage
The possibility of focus (and topic) marking within relative clauses suggests the presence of a Focus layer (and a Topic layer). This, then, is good evidence of the CP status of relative clauses in Japanese. As with DP-related focus, we do not know whether focused constituents obligatorily move to a higher FocP layer or whether they remain in situ. Recall that movement to a focus position was by no means a requirement for modifiers within the DP (Section \ref{focus,japanese}). Assuming a movement operation, example \ref{taroganondakusuri} would look as in \figref{Rc with focus movement} where the subject DP \textit{Tarō-ga}, having received nominative case in Spec, IP, moves to Spec, FocP. But of course, the focused subject DP could also receive in situ focus through the higher FocP layer.\footnote{I have not indicated an additional CP layer above FocP. In the analysis of \citet[57–65]{cinque:2020:relcl}, it is argued for some languages that this is necessary and why.} \is{focus}

\subsubsection{Complementizers}

\citet{kaplan:1995:the} provide morphological evidence for a CP status of relative clauses in \is{Classical Japanese} Classical Japanese and \is{Korean} Korean which they apply to Modern Japanese. Essentially, their argument is based on the fact that in Classical Japanese and Korean, verbs in attributive structures necessarily co-occur with overt complementizers, which project a CP layer. Their example and analysis of \is{Korean} Korean is depicted in example \ref{korean kaplan} below. \is{complementizer}

\exg. [ecey \textit{pro} manna-ass-te-n] salam (Korean)\\
yesterday \textit{pro} meet-\textsc{pst-rtro-comp} person\\
`the person [\textit{pro} met yesterday]’ \null \hfill \citep[30]{kaplan:1995:the} \label{korean kaplan}

It is the element -\textsc{(nu)n} that the authors interpret as an affixal complementizer. The argument is that it appears after all tensed elements as would be expected from \is{complementizer} complementizers, with C being higher than T \citep[30]{kaplan:1995:the}.

\largerpage[2.5]
Coming to \is{Classical Japanese} Classical Japanese, recall that this language had a much richer inflectional system than Modern Japanese. Importantly, the attributive (\textit{rentaikei}) and sentence-final form (\textit{shūshikei}) for inflecting elements differed morphologically. \citet[31–33]{kaplan:1995:the} therefore analyze the attributive (adnominal) form as projecting a CP layer, analogous to \is{Korean} Korean -\textsc{(nu)n}. This is illustrated in example \ref{keru original}, in which the aspectual verbal element -\textit{keri} is inflected for the attributive form -\textit{keru}, the ending /ru/ being the representative inflectional ending of the attributive form. As expected from its complementizer status, this element follows the tense affix/aspectual auxiliary verb -\textit{tari}, exactly parallel to the Korean example, thus providing another argument for C status.\footnote{Note that the adnominal form in some of the nine attested inflection classes of Classical Japanese is expressed by an additional, distinct, morpheme, which can be straightforwardly analyzed as a complementizer, but in some classes, sentence-final form and adnominal form exhibit different inflectional morphemes. This is the case for elements following the \textit{ri}-inflectional class, to which -\textit{keri} also belongs, and also for the historical inflectional paradigm of \textit{i}-adjectives which alternated between -\textit{shi} (sentence-final form) and -\textit{ki} (adnominal form) as shown below.

\par

\ex. \ag. taka-shi\\
high-\textsc{shk}\\
`is high' \null \hfill \textit{sentence-final}
\bg. taka-ki mono\\
high-\textsc{rtk} thing\\
`a high thing' \null \hfill \textit{attributive}

\par

See the discussion in Section \ref{complexdirect} on why -\textit{ki}, although it can be analyzed as a complementizer, does not necessarily only entail indirect modification.} I give the corresponding example in sentence-final form for clarity in \ref{keru}.\footnote{I gloss /ru/ as \textsc{comp} here, but it can be glossed as \textsc{rtk} (\textit{rentaikei}) with more emphasis on the inflection form.} \is{Classical Japanese} \is{complementizer}

\exg. ki-tar-i-ke-\textbf{ru} kariginu\\
wear-\textsc{perf-ryk-pst-comp} hunting.clothes\\
`the hunting clothes \textit{pro} was wearing' \\ \null \hfill Classical Japanese \textit{attributive form} \citep[32]{kaplan:1995:the} \label{keru original}

\exg. Kariginu-wo ki-tar-i-ke-\textbf{ri}.\\
hunting.clothes-\textsc{acc} wear-\textsc{perf-ryk-pst-shk}\\
`\textit{pro} was wearing hunting clothes.’ \\ \null \hfill Classical Japanese \textit{sentence-final form} \label{keru}

%\largerpage[2.5]
Another convincing piece of evidence that the attributive form fulfilled the role of a complementizer in \is{Classical Japanese} Classical Japanese is the phenomenon of \is{focus concord}\textit{focus concord} (Japanese: \textit{Kakari-Musubi}, Lit. `dependent-concluding’), as presented in \citet{kinsui:1995:iwayuru} and advocated for in \citet[§ 3.5]{hiraiwa:2005:dimensions}.\footnote{For descriptions of focus concord in Classical Japanese, see among others \citet{morishige:1971:nihon,funagi:1987:kakari}. See \citet{narrog:2019:origin} for a recent contribution.} \textit{Kakari-Musubi} refers to a phenomenon in which certain elements in the clause affect the inflectional form of the matrix verb. For example, clause-internal \textit{wh}-elements and focus elements trigger inflection of the verb in sentence-final position to attributive, not sentence-final, form. This is shown in example \ref{manyoshu}, in which the \textit{wh}-element \textit{izure-no} `when’ triggers an embedded question highlighted via the question particle -\textit{ka} and the matrix verb occurs in attributive form (\citealt[125–127]{hiraiwa:2005:dimensions} and see \citealt{nomura:1995:ka}).

\exg. Miyuki fu-ru koshi-no Ohoyama yuki-sugi-te izure-no hi-ni-\textbf{ka} wa-ga sato-wo mi-\textbf{mu} (Classical Japanese).\\
snow fall-\textsc{prs.rtk} mount-\textsc{mod} Ōyama go-pass-\textsc{ger} when-\textsc{mod} day-\textsc{loc-q} I-\textsc{gen} home-\textsc{acc} see-\textsc{comp}\\
`Crossing the snowy Mt. Ōyama, on which day can I see my home?' \\ \null \hfill (\citealt[125]{hiraiwa:2005:dimensions} from Manyōshū 3153) \label{manyoshu}

When verbs appear as complements of matrix verbs, they are also inflected for adnominal form \citep[10–11]{kinsui:1995:iwayuru}. In other words, the adnominal form performs the complementizing function in the absence of a complementizer. %\is{Classical Japanese} %\is{complementizer}

\exg. [Tomo-no empō-yori ki-tar-\textbf{u}]-wo yoroko-bu.\\
friend-\textsc{mod} away-\textsc{abl} come-\textsc{pst-comp}-\textsc{acc} delighted.at-\textsc{shk}\\
`\textit{pro} is delighted at (the fact) that a friend came all the long way.’ \\ \null \hfill \citep[10]{kinsui:1995:iwayuru}

These pieces of data give strong evidence that in \is{Classical Japanese} Classical Japanese the adnominal form served the purpose of a complementizer.

\citet{kaplan:1995:the} argue that the complementizer layer has been preserved in Modern Japanese, although it is not filled overtly, but via a null complementizer \citep[33–35]{kaplan:1995:the}. Indeed, it does seem convenient to not have to assume a categorial change of the category of relative clauses in the course of the language from CP to IP. The authors flesh out their analysis and aruge that  \textsc{na}, the non-past attributive form of \textit{na}-adjectives, is a copula that originates in V and raises through I to C ``where it supports the null complementizer" \citep[34]{kaplan:1995:the}. They give the following structure for a relative clause with an adjective where a nominal subject is present that relates to possession, similar to the examples mentioned in Section \ref{indirect mod}. \is{complementizer}

\exg. [\textsc{np} [\textsc{cp} [\textsc{ip} kami-ga kirei t\textsc{v} t\textsc{i}] [\textsc{c} [\textsc{i} [\textsc{v} na]]]] hito]\\
{} {} {} hair-\textsc{nom} pretty {} {} {} {} {} \textsc{cop} person\\
`a person whose hair is pretty' \\ \null \hfill Modern Japanese \citep[34]{kaplan:1995:the}

While I agree that these arguments in favor of a CP layer in Classical and in Modern Japanese (in certain instances) are rather convincing, the analysis of \textsc{na} always projecting such a layer cannot be maintained for the same reasons that Nishiyama’s analysis cannot be universally maintained, as explained in Section \ref{why it fails}. Besides the fact that it is not clear what triggers this movement to C other than the necessity to support the null complementizer, this analysis denies the possibility of direct modification, just like Nishiyama’s. We saw examples in Section \ref{complexdirect} where an additional CP layer, at least the possibility, did not immediately exclude direct modification. But this does not mean all modifiers always project a CP. %\is{complementizer}

\subsection{Derivation of indirect modifiers} \label{operator}

Further evidence for CP status comes from the derivation of Japanese relative clauses adopted here. Recall that the intriguing thing about the category of relative clauses in general is that the head noun needs to be active, somehow represented, in the matrix clause as well as inside the relative clause (the \is{connectivity problem}\textit{Connectivity Problem}, \citealt{salzmann:2017:reconstruction}). Recall furthermore that there are many technical implementations to solve this problem. In Japanese, the standard analysis is to assume a (silent) relative \is{operator}operator which ensures correct interpretability of the head noun inside the relative clause (see \citealt{miyamoto:2017:RC} a.o.), and this will be presented in the next subsection (for general accounts of operator movement within relative clauses see, among others, \citealt{chomsky:1977:on, jackendoff:1977:x} and \citealt[§ 2]{salzmann:2017:reconstruction} for an overview). 

\subsubsection{Operator movement}

Schematically, I assume that a head-external relative clause with a verbal nucleus is derived along the lines in \figref{derivation kareop} based on example \ref{kareop}, repeated here including the \is{operator}operator.
 
\exg. kare-ga \textit{Op}\textsubscript{1} kat-ta hon\textsubscript{1}\\
he-\textsc{nom} \textit{Op}\textsubscript{1} buy-\textsc{pst} book\textsubscript{1}\\
`the book (which) he bought' \label{kareop}

\begin{figure}
\centering
\begin{tikzpicture}[baseline=0pt, remember picture]
\tikzset{level distance=25pt, sibling distance=5pt}
 \Tree [.DP\textsubscript{internal} {} [.D′\textsubscript{internal} [.FP\textsubscript{indirect} [.CP \node (op1) {Op\textsubscript{1}}; [.C′ [.IP [.DP \edge[roof]; kare-ga ] [.I′ [.VP {} [.V′ \node (op2) {$<$Op$>$}; [.V kat ] ] ] [.I ta ] ] ] C ] ] [.F′\textsubscript{indirect} [.... [.NP \edge[roof]; hon\textsubscript{1} ] {} ] F ] ] D ] ]
\draw [color=blue][->] (op2) |- +(0,-2) -| (op1);
 \end{tikzpicture}
\caption{Derivation of relative clauses with operator movement (= \ref{kareop})}
\label{derivation kareop}
\end{figure}

Besides the assumptions that relative clauses are clear indirect modifiers and underlyingly single-headed (for reasons of simplicity, but see \citealt{cinque:2020:relcl} and Section \ref{dP}), I adopt an \is{operator}operator account \citep{ishii:1991:operators, miyamoto:2017:RC}, according to which it is a relative clause operator which ensures the correct interpretability of the noun within the relative clause and mediates the relationship between noun and relative clause, solving the \is{connectivity problem}connectivity problem \citep{salzmann:2017:reconstruction}. This operator \textit{Op} starts from within the relative clause, in this case, from the internal argument position of the verb since it is an object relative clause. It then raises to Spec, CP, which means such a position must be available as a landing site. It is exactly this operation, the raising of the operator, for which we can find evidence to be reviewed below.

At first glance, similar results can be achieved by assuming a \is{raising}\textit{raising} analysis within a \is{double-headed relative clause}double-headed framework, with the difference that the noun itself is the internal head of the relative clause, in addition to an identical counterpart (the external head in the matrix clause), and then moves to Spec, CP. The external head is then deleted following this operation. This derivation likewise predicts movement effects in an identical fashion (see \citealt[22–36]{cinque:2020:relcl} for arguments and \citealt[65–71]{cinque:2020:relcl} for a discussion on prenominal RC-languages).

Complications arise when trying to apply this derivation to other kinds of relative clauses in Japanese, for instance head-internal relative clauses \citep{hiraiwa:2012:mechanism, hiraiwa:2017:a}. Furthermore, some sort of remnant movement of the relative clause would need to be stipulated following the raising of the noun to ensure the correct word order. These aspects would complicate the analysis, which is why I will stick with the simpler operator derivation.

\subsubsection{Evidence}

The \is{operator}operator-derivation can be motivated by identifying various \is{reconstruction effects}reconstruction effects, which, in turn, provide clear evidence for movement \citep{schachter:1973:focus, vries:2002:syntax}, most of which have been reported for Japanese in \citet{ishii:1991:operators} and are discussed in \citet{miyamoto:2017:RC} and \citet{whitman:2013:prehead}. A prominent reconstruction effect involves pronominal anaphor binding with the heavy pronouns \textit{karejishin/kanojojishin} `himself/herself' or \textit{jibunjishin} `oneself' (\citealt[615]{miyamoto:2017:RC} quoting \citealt[29]{ishii:1991:operators}), similar to English \textit{the picture of himself}\textsubscript{1} [\textit{which John}\textsubscript{1} \textit{likes e}]]), \citep{salzmann:2017:reconstruction}. Observe \ref{anaphorbinding}.

\exg. Marī-wa [Jon\textsubscript{1}-ga \textit{e}\textsubscript{2} taipu-shi-ta [karejishin\textsubscript{1}-no ronbun\textsubscript{2}-wo]] mot-te-ki-ta.\\
Mary-\textsc{top} John-\textsc{nom} \textit{e} type-do-\textsc{pst} himself-\textsc{mod} paper-\textsc{acc} bring-\textsc{ger}-come-\textsc{pst}\\
`Mary brought himself's paper that John typed.’ \\ \null \hfill \citep[615]{miyamoto:2017:RC} \label{anaphorbinding}

In order for the reflexive \textit{karejishin} `himself’ to find its antecedent \textit{Jon}, the two elements, or some instances of them (the \is{operator}operator), must have been in a c-command relation at some point in the derivation. The underlining matrix clause in \ref{anaphormatrix} confirms that movement must have happened \citep[615–616]{miyamoto:2017:RC}.\footnote{Allegedly the same holds with raised quantifiers, according to \citet{whitman:2013:prehead} even with the canonical reflexive pronoun \textit{jibun} `self’.} \is{reconstruction effects}

\exg. Jon\textsubscript{1}-ga karejishin\textsubscript{1}-no ronbun\textsubscript{2}-wo taipu-shi-ta.\\
John-\textsc{nom} himself-\textsc{mod} paper-\textsc{acc} type-do-\textsc{pst}\\
`John typed his paper.’ \label{anaphormatrix}

Furthermore, even though Japanese has been argued, for example in \citet{kuno:1973:structure}, to lack island effects in relative clauses in general, it does show them clearly with object relative clauses (\citealt[623–626]{miyamoto:2017:RC}, see also \citealt[369–373]{whitman:2013:prehead}).\footnote{See \citet{takahashi:2025:relative} for a recent contribution and discussion. The authors show, based on experimental data, that extraction out of relative clauses is marked in Japanese in general, but even more so with object relative clauses.}

\exg. *[Bill-ga, [[\textit{e}\textsubscript{1} \textit{e}\textsubscript{2} kai-ta] hon\textsubscript{2}-wo] yakushi-te-i-ru] gakusha\textsubscript{1}\\
Bill-\textsc{nom} \textit{e}\textsubscript{1} \textit{e}\textsubscript{2} write-\textsc{pst} book\textsubscript{2}-\textsc{acc} translate-\textsc{ger}-be-\textsc{prs} scholar\textsubscript{1}\\
(`the scholar\textsubscript{1} that Bill is translating the book\textsubscript{2} that t\textsubscript{1} wrote’) \\ \null \hfill \citep[178]{inoue:1976:togo} \label{islandjapanese}

The corresponding root clauses of \ref{islandjapanese} are repeated below for convenience.

\exg. Gakusha\textsubscript{1}-ga hon\textsubscript{2}-wo kai-ta. Bill-ga sono hon\textsubscript{2}-wo yakushi-te-i-ru.\\
scholar\textsubscript{1}-\textsc{nom} book\textsubscript{2}-\textsc{acc} write-\textsc{pst} Bill-\textsc{nom} that book\textsubscript{2}-\textsc{acc} translate-\textsc{ger}-be-\textsc{prs}\\
`A scholar\textsubscript{1} wrote a book\textsubscript{2}. Bill is translating that book\textsubscript{2}.’

\largerpage[-4]
Finally, Japanese allows relativization of idioms \citep{inoue:1976:togo, ishii:1991:operators, miyamoto:2017:RC, whitman:2013:prehead}, which is another argument for movement within relative clauses often given (\citealt{schachter:1973:focus, carlson:1977:amount}, see also \citealt{bhatt:2002:raising, bhatt:2015:rc, sauerland:2003:a, cinque:2020:relcl}). The argument goes that if a noun is relativized and an \is{reconstruction effects}idiomatic meaning with the modifying verb is retained, these elements must have formed constituents at one point of the derivation. Observe the example \ref{idiom rc} below.\footnote{Idioms are a very difficult case, even in English, and recent discussion has cast doubt upon the validity of the argument that they must be involved in movement derivations \citep{heycock:2019:relative, webelhuth:2019:idioms}. Furthermore, as mentioned in \citet[48]{bhatt:2002:raising}, this argument does not rule out other analyses. See also the appendix in \citet{cinque:2020:relcl}. In \citet{hoshi:2004:para}, it was argued that the relativization of heads taking part in idiomatic constructions is ruled out in Japanese, but \citet{inoue:1976:togo, whitman:2013:prehead} and \citet{miyamoto:2017:RC} have argued otherwise. Ken Hiraiwa (p.c.) and other speakers have told me that they do not find the idiomatic expression reported in the text entirely natural.} \is{English}

\exg. [Karera-ga magarinarinimo \textit{e}\textsubscript{1} tsuke-ta] ketchaku\textsubscript{1}-wa, amari yorokoba-re-na-k-at-ta.\\
they-\textsc{nom} somehow.or.other \textit{e} reach-\textsc{pst} conclusion-\textsc{top} much be.happy-\textsc{pass-neg-pred-cop-pst}\\
`The end [they somehow or other brought it to] was not very well received (by others).’ (\textit{ketchaku-wo tsukeru} = to reach a conclusion = to bring to an end) \\ \null \hfill \citep[214]{inoue:1976:togo} \label{idiom rc}

These pieces of data suggest that movement takes place in Japanese relative clauses, in turn motivating the \is{operator}operator derivation. And this derivation only leaves room for an analysis of relative clauses as CPs since the operator needs a landing site.

\subsubsection{Relative clauses and major subjects} \label{majorsubject}

There is a class of relative clauses that have been argued to be derived from so-called \textit{major subject constructions} (\citealt[§ 9]{kuroda:1992:japanese}; cf. \citealt{kuno:1973:structure, heycock:1993:syntactic, whitman:2013:prehead, miyamoto:2017:RC}). Examples were briefly mentioned in Section \ref{indirect mod} (and also Section \ref{reason and manner}) and look as follows. Recall that the examples below are cases of relative clauses where the non-verbal modifiers are followed by the monomoraic realizations of \is{MOD@\textsc{mod}}\textsc{mod}, but that forms of the copula can follow the respective modifiers instead. \is{major subject construction}

\exg. [Tarō-ga shujinkō-no] monogatari\\
Taro-\textsc{nom} main.protagonist-\textsc{mod} story\\
`story [in which Taro is the main protagonist]’ \null \hfill \citep[66]{watanabe:2010:notes} \label{shujinko1}

\exg. [shūshoku-ga taihen-na] butsurigaku\\
job-\textsc{nom} hard-\textsc{mod} physics\\
`physics, [where finding a job is difficult]’ \null \hfill \citep[255]{kuno:1973:structure} \label{butsurigaku}

In \ref{shujinkodouble} and \ref{butsurigakudouble}, the underlying constructions of \ref{shujinko1} and \ref{butsurigaku} involving a major subject are given. The examples are adapted from the authors mentioned above. \ref{butsurigakudouble} is also given in \citet[179]{heycock:1993:syntactic}.\footnote{A reviewer wonders whether \ref{shujinko1} might not be derived from a prepositional phrase in which (\textit{sono}) \textit{monogatari} is a locative expression. However, note that this expression could serve as a major subject in a major subject construction as well, since, as pointed out in Section \ref{reason and manner}, this is generally possible for locative (and temporal) expressions (see example \ref{sono hi}). In other words, even if we take the relativized noun to have been part of a prepositional phrase, the relative clause is arguably derived from a major subject construction.} \is{major subject construction} \is{operator}

\exg. Sono monogatari-ga Tarō-ga shujinkō-d-a.\\
this story-\textsc{nom} Taro-\textsc{nom} main.protagonist-\textsc{pred-cop.prs}\\
`(As for) this story, Taro is the main protagonist.’ \label{shujinkodouble}

\exg. Butsurigaku-ga shūshoku-ga taihen-d-a.\\
physics-\textsc{nom} employment-\textsc{nom} difficult-\textsc{pred-cop.prs}\\
`Physics is such that finding a job is difficult.’ \label{butsurigakudouble}

In major subject constructions, two nouns are bestowed with the nominative particle. This is thus a case where a subject-marked constituent does not bear a theta role \citep{heycock:1993:syntactic}. Rather, the higher \textit{ga}-marked noun seems to have a \is{focus}focus (or potentially topic) interpretation (cf. \citealt{vermeulen:2005:poss}).

Major subject constructions have been used as an underlying mechanism to account for seemingly gapless relative clauses such as \ref{butsurigaku} since \citet{kuroda:1992:japanese}. This suggests that the relative clause in \ref{shujinko1} has the complex structure depicted in \figref{Structure major subject}. See also \citet{heycock:1993:syntactic} and especially \citet{miyamoto:2017:RC}, on whom I base the derivation here, for more arguments and discussion, also of example \ref{butsurigaku}.

\begin{figure}
\centering
\begin{tikzpicture}[baseline=0pt, remember picture]
\tikzset{level distance=25pt, sibling distance=2pt}
 \Tree [.CP Op [.C′ [.IP [.DP\textsubscript{2} {} [.D′ [.FP {} [.F′ [.... [.NP \edge[roof]; {Tarō-ga} ] {} ] F ] ] D ] ] [.I′ [.DP\textsubscript{1} {} [.D′ [.FP $<$Op$>$ [.F′ [.... [.NP {} [.N′ {} [.N shujinkō ] ] ] {} ] F ] ] D ] ] I ] ] C ] ]
\end{tikzpicture}
\caption{Structure of relative clause in \ref{shujinko1}}
\label{Structure major subject}
\end{figure}

\largerpage
The modifying clause contains two nouns, \textit{Tarō} and \textit{shujinkō} `main protagonist’, which, in fact, must both project a DP. The DP \textit{Tarō}, the subject within the relative clause, is situated in Spec, IP of the CP in which \textit{shujinkō} is embedded (although it is questionable whether it originates from somewhere within that DP\textsubscript{1}). Furthermore, the relative clause contains an \is{operator}operator corresponding to the modified noun \textit{monogatari} `story’. The \is{major subject construction}major subject construction given above shows that the operator must originate from somewhere within the relative clause, ensuring the correct interpretation. Accounting for \is{operator}its role as a modifier, under the theory adopted here, I assume it is situated in the extended projection of DP\textsubscript{1}.\footnote{\citet{miyamoto:2017:RC}, among others, assumes a subject position, but likely nothing hinges on this. The point is that the operator originates from somewhere within the relative clause.}

The derivation is slightly different, and in fact clearer, for relative clauses of the following kind with an adjectival nucleus co-occurring with a subject, which in turn expresses a possession-relation with the modified noun.

\exg. [doresu-ga utsukushi-k-at-ta] josei\\
dress-\textsc{nom} beautiful-\textsc{pred-cop-pst} woman\\
`the woman [whose dress was beautiful]’ \label{utsukushikatta op}

In this sentence, the dress belongs to the woman and is described as beautiful. Such sentences, arguably, also have a corresponding \is{major subject construction}major subject construction, namely \ref{doresumajor} (and see Section \ref{indirect mod}). Here, \textit{sono josei} `that woman' fulfills several characteristics of a major subject: It has a \is{focus}focus interpretation, it does not bear a theta role, and it seems to be followed by a \is{operator}prosodic break \citep[1346–1348]{vermeulen:2005:poss}.

\exg. Sono josei-ga (sono josei-no) doresu-ga utsukushi-k-at-ta.\\
that woman-\textsc{nom} that woman-\textsc{mod} dress-\textsc{nom} beautiful-\textsc{pred-cop-pst}\\
`As for that woman, her dress was beautiful.’ \label{doresumajor}

The relevant difference from \ref{shujinko1} is that the \is{operator}operator, representing the noun \textit{josei} `woman’, is closely related to the subject DP within the relative clause \textit{doresu} `dress’, not to the modifier, which is not a noun, but an adjective. I therefore assume that the operator starts out from the extended projection of that DP (here DP\textsubscript{2}). This explains the structural differences from the derivation of the relative clause in \ref{shujinko1}, depicted in \figref{Structure major subject}, and the derivation of \ref{utsukushikatta op} in \figref{Structure major subject utsukushikatta}. \is{major subject construction}

% \begin{figure}
% \tiny
% \centering
% \begin{tikzpicture}[baseline=0pt, remember picture]
% \tikzset{level distance=25pt, sibling distance=5pt}
% \Tree [.DP\textsubscript{internal} {} [.D′\textsubscript{external} [.FP\textsubscript{indirect} [.CP Op\textsubscript{1} [.C′ [.IP [.DP\textsubscript{2} {} [.D′ [.FP\textsubscript{Poss} $<$Op\textsubscript{1}$>$ [.F′\textsubscript{Poss} [.... [.NP \edge[roof]; doresu-ga ] {} ] F ] ] D ] ] [.I′ [.CopP {} [.Cop′ [.PredP {} [.Pred′ [.AP {} [.A′ $<$DP\textsubscript{2}$>$ [.A utsukushi ] ] ] [.Pred k(u) ] ] ] [.Cop at ] ] ] [.I ta ] ] ] {} ] ] [.F′\textsubscript{indirect} [.NP \edge[roof]; josei\textsubscript{1} ] F ] ] D ] ]
% \end{tikzpicture}
% \end{figure}


\begin{figure}
\begin{forest}
[DP\textsubscript{internal}
    [~]
    [D′\textsubscript{external}
        [FP\textsubscript{indirect}
            [CP
                [Op\textsubscript{1}]
                [C′
                    [IP
                        [DP\textsubscript{2}
                            [~]
                            [D′
                                [FP\textsubscript{Poss}
                                    [$<$Op\textsubscript{1}$>$ ]
                                    [F′\textsubscript{Poss}
                                        […
                                            [.NP
                                                [doresu-ga, roof]
                                            ]
                                            []
                                        ]
                                        [F]
                                    ]
                                ]
                                [D]
                            ]
                        ]
                        [I′
                            [CopP
                                [~]
                                [Cop′
                                    [PredP
                                        [~]
                                        [Pred′
                                            [AP
                                                [~]
                                                [A′
                                                    [$<$DP\textsubscript{2}$>$]
                                                    [A
                                                        [utsukushi]
                                                    ]
                                                ]
                                            ]
                                            [Pred
                                                [k(u)]
                                            ]
                                        ]
                                    ]
                                    [Cop
                                        [at]
                                    ]
                                    ]
                            ]
                            [I
                                [ta]
                            ]
                    ]
                ]
                [~]
        ]
        ]
        [F′\textsubscript{indirect}
            [NP
                [josei\textsubscript{1},roof]
            ]
            [F]
    ]
]
    [D]
]
]
    \end{forest}
\caption{Structure of \ref{utsukushikatta op}}
\label{Structure major subject utsukushikatta}
\end{figure}

Crucially, the DP \textit{doresu} `dress’, in which the operator, co-indexed with the head noun, is contained, starts out as the complement of the adjective, ensuring the meaning that it is the dress that is beautiful not the woman, and subsequently raises to Spec, IP of the relative clause. The operator, in turn, raises from the extended projection of DP\textsubscript{2} to Spec, CP. That this must be the case, that some sort of movement takes place, was shown through the availability of reconstruction effects. Here, it is also possible to argue for a precise position of this operator within the extended projection, as it is most likely a position related to possession, ensuring the possessive relationship between `dress’ and `woman’. I tentatively give a functional projection FP\textsubscript{Poss} (although this might not be relevant, given the broad nature of functional projections in Japanese). Overall, the structure is slightly different than for \ref{shujinko1} (and probably even more complex for \ref{butsurigaku}). \is{major subject construction}

To summarize, the derivation of relative clauses assumed in this book makes use of operator movement and this gives good evidence for the CP status of Japanese relative clauses. A derivation along these lines is assumed to be universal, thus holds also in seemingly simpler DPs such as \ref{kareop} mentioned at the beginning of the chapter. In each case, an operator, which is co-indexed with the head noun, raises from within the relative clause to Spec, CP.

\subsection{Summary}

This section showed why the size of relative clauses in Japanese is hard to determine, and eventually presented a unifying analysis as CPs. It was shown that neither the arguments that relative clauses cannot be bigger than IPs, nor that they are of different sizes were convincing. The crucial part to take away is that relative clauses in Japanese only come in one size, which makes the dichotomy \textit{reduced} vs. \textit{full} as well as \textit{CP} vs. \textit{IP} superfluous, and with it the necessity of dedicated projections for these categories.

\section{On the position of relative clauses} \label{complexity}

Assuming there are several functional projections in the indirect domain, above as well as below numerals, and that relative clauses are all CPs suggests that placement is entirely arbitrary. The question is if this is really the case or if there are factors influencing the position of relative clauses. In this section, it will be argued that the higher positions are reached via independent movement and that a factor influencing this movement is \is{complexity}complexity.

\subsection{Introduction}

Recall that cross-linguistically the following structural hierarchy is at issue: \textit{demonstrative} $>$ \textit{full relative clause} $>$ \textit{numeral} $>$ \textit{adjective/reduced relative clause}. Translated into terms of head proximity, full relative clauses are further away from the noun than \is{numeral}numerals, which in turn are further away than reduced relative clauses (cf. \citealt[§ 3.5]{cinque:2020:relcl}; \citealt{cinque:2017:micro, greenberg:1963:some, hawkins:1983:word, rijkhoff:2002:noun}). This is especially easy to visualize in languages in which all modifiers strictly appear on one side of the noun.

While this applies to Japanese, being strictly head-final, this language does not differentiate between full and reduced relative clauses. This predicts that relative clauses can productively appear to the left or to the right of numerals, which are also situated in the indirect domain. And in fact, speakers generally do not have any preference regarding the order of these elements.

\ex. \label{pati} \ag. [pātī-de at-ta [futari-no josei]]\\
party-\textsc{loc} meet-\textsc{pst} two-\textsc{mod} woman \label{pati3}\\
\bg. [futari-no [pātī-de at-ta josei]]\\
two-\textsc{mod} party-\textsc{loc} meet-\textsc{pst} woman\\
`the two women who I met at a party' \label{pati4}

In \ref{pati3}, the relative clause in linear order precedes the numeral, while in \ref{pati4} the relative clause linearly follows the numeral. Taking linear order as an indication of structural position suggests that in \ref{pati3} the relative clause is situated in a functional projection above FP\textsubscript{Num}, whereas the relative clause in \ref{pati4} is in a functional projection below it. Whether the latter case is the result of movement, or the relative clause was base-generated in a higher position can at first glance not be deduced. Crucially, the two positions are not the result of the relative clauses being of a different size in \ref{pati3} and \ref{pati4}.

The question now is whether the two options, and ordering of these constituents in general, are simply random. Importantly, speakers find both orders in \ref{pati} perfectly fine and do not have a preference for one order over the other.\footnote{Recall that this is not the case for direct modifiers, which are significantly degraded to the left of numerals.} To answer this question, one will first have to investigate if there is a way to make one of the options more or less natural, or more or less acceptable, in contrast to the alternative.

These considerations naturally also apply to adjectival relative clauses, those containing an adjectival nucleus. Recall that \textit{prima facie} adjectives denoting non-past are ambiguous between a direct and an indirect status, which is why specific tests need to be applied to identify direct modifiers (Section \ref{testsdirect}). Unambiguous indirect adjectival modifiers can come in several types, but a clear indication is inflection in past tense as well as co-occurrence with subjects.\footnote{Recall that some modifiers analyzed as direct do allow subjects, as discussed in Section \ref{complexdirect}.} Starting with inflected adjectives, speakers again find both orders in the ordering pairs presented below grammatical and, more importantly, do not have a preference for one option: \textit{adjectival relative clause} $\succ$ \textit{numeral} as well as \textit{numeral} $\succ$ \textit{adjectival relative clause}. The following examples are all taken and adapted from findings in the corpus study reported further below in Section \ref{section corpus study}.\footnote{Although it is a fine line between grammaticality on the one hand and preference on the other, this difference is crucial to be aware of. When discussing such minimal pairs of examples, I checked always for both parameters, although in most cases only a difference in preference emerged. I mark the non-preferred version in the following via the symbol \textit{\#}.} \is{numeral}

\ex. \label{kanashikatta} \ag. [kanashi-k-at-ta [mittsu-no dorama]]\\
sad-\textsc{pred-cop-pst} three-\textsc{mod} soap.opera\\
\bg. [mittsu-no [kanashi-k-at-ta dorama]]\\
three-\textsc{mod} sad-\textsc{pred-cop-pst} soap.opera\\
`three soap operas that were sad' \label{soapoperajapanese}

\ex. \label{kanashikattafutari} \ag. [utsukushi-k-at-ta [futari-no josei]]\\
beautiful-\textsc{pred-cop-pst} two-\textsc{mod} woman\\
\bg. [futari-no [utsukushi-k-at-ta josei]]\\
two-\textsc{mod} beautiful-\textsc{pred-cop-pst} woman\\
`the two women who were beautiful'

Corresponding examples in other languages do show ordering restrictions. For instance, in \is{German} German, when a participial -- that is, \is{reduced relative clause}reduced -- relative clause and a numeral co-occur, placing the numeral to the left of the reduced relative clause is strongly preferred. In \ref{soapoperagerman}, I translated the examples of \ref{soapoperajapanese} into German, and, although of course there might be subtle differences in information structure depending on the context, according to my judgment, the version in which the reduced relative clause follows the numeral is not only much more natural than the reversed order, but I find the reversed order quite ungrammatical (see also the pairs in example \ref{numeral rrc german} in Section \ref{dP}).\footnote{Japanese also contrasts with Mandarin, another strict head-final language. \citet{hsu:2017:alternatives} reports on many more occurrences of relative clauses in a lower position than in a higher position above numerals and demonstratives in a corpus study he performed and judges the latter case as marked and focused.}

\ex. \ag. die drei traurig gewesenen Seifenopern (German)\\
the three sad be.\textsc{pst.ptcp.pl} soap.opera\\
\bg. ??die traurig gewesenen drei Seifenopern \\
the sad be.\textsc{pst.ptcp.pl} three soap.opera\\
`the three soap operas that were sad' \label{soapoperagerman}

In fact, there are ways to achieve a similar effect in Japanese, a preference for one order over the other. Significantly, this order is the reverse of the \is{German} German case. One way to achieve this is to add degree modifiers such as \textit{totemo} `very’ to the relative clause. Almost all speakers I consulted expressed a higher preference for the version with the relative clause in first position, contrary to the minimal pair without the degree adverb, where both options are fine. Given below is the relevant altered version of example \ref{kanashikatta} with the degree adverb. \is{complexity}

\ex. \label{dorama} \ag. [totemo kanashi-k-at-ta [mittsu-no dorama]]\\
very sad-\textsc{pred-cop-pst} three-\textsc{mod} soap.opera\\
\bg. \#[mittsu-no [totemo kanashi-k-at-ta dorama]]\\
three-\textsc{mod} very sad-\textsc{pred-cop-pst} soap.opera\\
`three soap operas that were very sad'

The same effect is achieved when an overt subject appears. Again, speakers showed a stronger preference for the relative clause in first position, now tested for the corresponding example to \ref{kanashikattafutari}.

\ex. \label{dorama with subject} \ag. [doresu-ga utsukushi-k-at-ta [futari-no josei]]\\
dress-\textsc{nom} beautiful-\textsc{pred-cop-pst} two-\textsc{mod} woman\\
\bg. \#[futari-no [doresu-ga utsukushi-k-at-ta josei]]\\
two-\textsc{mod} dress-\textsc{nom} beautiful-\textsc{pred-cop-pst} woman\\
`the two women whose dresses were beautiful'

Importantly, if a certain order is preferred, it is always the one in which the relative clause precedes the \is{numeral}numeral, never the other. But why do speakers prefer the position of the relative clause to the left of the numeral, once a degree modifier or a subject are present?

These relative clauses have been made more \textit{complex} or \textit{heavier}. Although \textit{complexity} and \textit{heaviness} are rather evasive terms and hard to define, let alone measure, I will attempt a descriptive definition below. \is{complexity}

For \textit{syntactic complexity}, I adopt the following definition of \citet{hawkins:2007:efficiency}, who relies on an idea going back to \citet{miller:1963:finitary}.

\ex. Complexity is a function of the amount of structure that is associated with the terminal elements, or words, of a sentence. More structure means, in effect, that more linguistic properties have to be processed in addition to recognizing or producing the words themselves. \\ \null \hfill \citep[8]{hawkins:2007:efficiency}

This means we can speak of higher \is{complexity}complexity when a phrase contains additional constituents, especially functional ones. This is definitely the case when an overt subject appears or when an adverbial phrase is merged in the CP domain of the relative clause.

\textit{Prosodic complexity} or \textit{heaviness} on the other hand can be described as equal to length. In other words, the more morae a relative clause consists of, the longer, and in turn the more complex it is.\footnote{Some speakers have a higher tolerance than others in this respect. For example Shōta Momma (p.c.) has pointed out to me that for him, the more complex a relative clause to the right of a \is{numeral}numeral is made, the less natural he finds it. However, he would still not willingly shift the two units around.}

In the following, I will present more empirical evidence for the observation that higher complexity correlates with higher preference of relative clauses to the left of numerals, based on a corpus study and subsequent discussion with native speakers on this topic.

\subsection{Complex relative clauses in higher positions: Corpus study} \label{section corpus study}

In order to get a clearer picture of the phenomenon, I have performed a search query for all occurrences of all adjectival relative clauses containing an adjective inflected for past and a numeral within the same DP, thus modifying the same noun, in the Japanese web corpus \is{jaTenTen}jaTenTen11 \citep{jakubi:2013:ten}.\footnote{This corpus is accessible via Sketch Engine, \url{https://auth.sketchengine.eu} \citep{kilgarriff:2014:sketch}.} The goal was to see which order of the two constituents is generally more attested and in what way additional constituents inside the relative clause have influence on its position. \is{complexity}

There are several reasons why I have chosen to investigate (solely) adjectival relative clauses. From a practical perspective, increasing complexity of this type of relative clause is a straightforward task since degree adverbs are usually freely compatible with adjectives -- which is not the case for most verbs and by extension verbal relative clauses, where such combinations are often semantically impossible.\footnote{A related point is that adjectives were considerably easier to test in the corpus at hand, mainly because I could focus on inflections in past tense. Not only were those clear signs of indirect modification, there were no other tense or aspect forms to test in attributive position. For verbs, I would have needed to include search operations for non-past and other tense and aspect forms, such as progressive, since these all mark instances of indirect modification. Furthermore, syntactic dependencies are not annotated in the interface at hand, which means that getting results for verbal relative clauses as nominal modifiers is quite cumbersome. In fact, one would need to look for all occurrences of a verb and a noun at an arbitrary distance within the same sentence and then to annotate manually whether there are dependencies and what constituents intervene, if any. This aspect is therefore left for future research and potentially a more fitting manner of elicitation.} From a theoretical perspective, I expected this study to complete the picture of adjectives in Japanese, and to give a nice contrast to direct modifiers, out of which adjectives are argued to be the prototypical representatives. On the other hand, one must be aware of the fact that adjectives inflected for past and \is{numeral}numerals modifying the same noun might not be entirely natural. \is{complexity} \is{jaTenTen}

Concretely, I performed searches for numerals,\footnote{Numerals in Japanese consist of a cardinal number morpheme and a numeral classifier reflecting the physical appearance (among other characteristics) of the head noun. There is also an underspecified option in which no classifier appears. In each case, \is{numeral}numerals are within the DP obligatorily licensed via \textsc{mod} spelled out as /no/.} either followed or preceded by an \textit{i}- or a \textit{na}-adjective inflected for past, then followed by a noun. This means in total four different search scenarios were tested. Results were looked through and false positives removed in the process. The final results after cleaning the data look as follows.\footnote{Needless to say, since extraction of false positives was performed manually and this process is rather error-prone, the results must be interpreted with caution.} \is{jaTenTen}

\ex. \emph{Order of numerals and adjectives in past in Japanese attested in jaTenTen11 (false positives removed, 2022/09/26)}
\a. \textit{i}-adjective in past $\succ$ numeral (iA\_Num): 149 results
\b. Numeral $\succ$ \textit{i}-adjective in past (Num\_iA): 26 results
\c. \textit{na}-adjective in past $\succ$ numeral (naA\_Num): 74 results
\d. Numeral $\succ$ \textit{na}-adjective in past (Num\_naA): 7 results

It is striking, especially for a web corpus, that so few results were obtained. This suggests that the combination of inflected adjectives and \is{numeral}numerals is rare in naturally occurring speech. In practically all cases, an appropriate context is needed to justify an inflection of the adjective in past over non-past, as also mentioned by the consulted native speakers (see Section \ref{section closer look}).

The first obvious observation is that a higher number of occurrences for the order \textit{adjectival relative clause} $\succ$ \textit{numeral} is found than for the reversed order. Importantly, the type of adjective, whether the adjective is an \textit{i}- or a \textit{na}-adjective, is not an influencing factor.

In the next step, I annotated every relative clause manually according to whether additional constituents appear or not. \figref{Co-occurrences of adjectival relative clauses and numerals in Sketch Engine} displays all relevant numbers. It indicates how many relative clauses contained additional constituents for each ordering pair and both adjective groups.

\begin{figure}

%\vspace{-3cm}
\includegraphics[scale=1.2, width=12cm]{rcconstituents.png}
 \caption{Co-occurrences of adjectival relative clauses and numerals in Sketch Engine}
 \label{Co-occurrences of adjectival relative clauses and numerals in Sketch Engine}
\end{figure}

Taking additional constituents as an indicator of higher \is{complexity}complexity, these results confirm the initial observation: More complex adjectival relative clauses are favored to the left of \is{numeral}numerals. In fact, 81\% (121 out of 149) of relative clauses containing \textit{i}-adjectives and 64\% (48 out of 74) of those containing \textit{na}-adjectives which appear to the left of the numeral contain at least one additional constituent. For the reversed position, no such trend is detectable: In the position below the numeral, only 20\% (5 out of 21) of relative clauses with \textit{i}- and 42\% (3 out of 7) of relative clauses with \textit{na}-adjectives contain an additional constituent. \is{jaTenTen}

In the final step, I have analyzed which constituents exactly appear in those relative clauses and annotated each relative clause according to the type of constituent it co-occurs with. This allowed a categorization of the relative clauses with constituents into different categories. There are five types of constituents, and, of course, a relative clause can contain more than one additional constituent, although this was uncommon. The most prominent constituents are \textit{degree adverbs} and \textit{subjects}, as exemplified above (examples \ref{dorama} and \ref{dorama with subject}). Other types of constituents are \textit{adverbial adjuncts}, \textit{verbal clauses} and \textit{compounds}. The relevant numbers are given for the orders \textit{adjectival relative clause} $\succ$ \textit{numeral} for both adjective groups in tables \ref{orderingtable1} and \ref{orderingtable2}.

\begin{table}[p]
\begin{tabular}{lr}
\lsptoprule
Type & Amount \\
\midrule
compound construction  & 1  \\
subject & 59  \\
indirect object & 0 \\
degree adverb & 18 \\
adverbial adjunct & 8 \\
verbal clause & 18 \\
adjectival clause & 0 \\
compound + degree adverb & 1 \\
subject + degree adverb & 26 \\
verbal clause + degree adverb & 1 \\
\lspbottomrule
\end{tabular}
\caption{types of constituents inside adjectival relative clauses in the order \textit{i}-adjective $\succ$ numeral}
\label{orderingtable1}
\end{table}


\begin{table}[p]
\begin{tabular}{lr}
\lsptoprule
Type & Amount \\
\midrule
compound construction  & 4  \\
subject & 8  \\
indirect object & 1  \\
degree adverb & 4 \\
adverbial adjunct & 18 \\
verbal clause & 9 \\
adjectival clause & 1 \\
compound + degree adverb & 0 \\
subject + degree adverb & 3 \\
verbal clause + degree adverb & 0 \\
\lspbottomrule
\end{tabular}
\caption{types of constituents inside adjectival relative clauses in the order \textit{na}-adjective $\succ$ numeral}
\label{orderingtable2}
\end{table}

\largerpage
To illustrate what is meant by these \is{complexity}labels, representative examples are presented in the following examples. For reasons of space and convenience, I give only the DPs and for reasons of readibility, the sources are given in footnotes.

\exg. [kaiseki-kanō-de-at-ta [go-ko-no hai]]\\
analysis-possible-\textsc{pred-cop-pst} five-\textsc{cl-mod} embryo\\
`five embryos which were possible to analyze'\\ \null \hfill \textit{na}-adjective $\succ$ numeral; compound\footnote{\url{castle104.com}, 2022/09/26.}

\exg. [inshō-ga yo-k-at-ta [hitotsu-no riyū]] \\
impression-\textsc{nom} good-\textsc{pred-cop-pst} one-\textsc{mod} reason\\
`one reason why the impression was good' \\ \null \hfill \textit{i}-adjective $\succ$ numeral; subject\footnote{\url{http://cocolog-nifty.com}, 2022/09/26.}

\exg. [kono mecha kata-k-at-ta [ni-hon-no neji]]\\
this very hard-\textsc{pred-cop-pst} two-\textsc{cl-mod} screw\\
`those two screws which were very hard'\\ \null \hfill \textit{i}-adjective $\succ$ numeral; degree adverb\footnote{\url{https://monote.net}, 2022/09/26.}

\exg. [kodomo-no koro-ni suki-d-at-ta [ni-hiki-no neko]]\\ %\label{neko}
child-\textsc{mod} time-\textsc{loc} fond-\textsc{pred-cop-pst} two-\textsc{cl-mod} cat\\
`the two cats I was fond of as a child' \\ \null \hfill \textit{na}-adjective $\succ$ numeral; temporal adverbial adjunct\footnote{\url{http://honeypeace.blog11.fc2.com/blog-date-20070808.html}, 2022/09/26.}

\exg. [kono kōen-de shit-te yo-k-at-ta [futatsu-no kotoba]] \\
this lecture-\textsc{loc} know-\textsc{ger} good-\textsc{pred-cop-pst} two-\textsc{mod} word\\
`those two words for which it was good that I learned them in that lecture' \\ \null \hfill \textit{i}-adjective $\succ$ numeral; verbal clause\footnote{\url{fc2.com}, 2022/09/26.}

\exg. [higai-ga mottomo ōki-k-at-ta [itsutsu-no mura]] \\ %\label{higai}
damage-\textsc{nom} most big-\textsc{pred-cop-pst} five-\textsc{mod} village\\
`the five villages where the damage was worst' \\ \null \hfill \textit{i}-adjective $\succ$ numeral; subject + degree adverb\footnote{\url{http://www.ekenzo.jp/shopdetail/006002000016/}, 2022/09/26.}

The emerging picture from these findings further confirms the initial hypothesis: Indirect modifiers of higher \is{complexity}complexity are preferred to the left of \is{numeral}numerals. From a structural perspective, this means they are situated in a functional projection higher than FP\textsubscript{Num}.

\subsection{A closer look at the findings} \label{section closer look}
\largerpage[2]
Since the problem remains that only a few relevant examples were attested and because corpora, although being representative samples of a language, are only representative to a certain extent, I have followed up this corpus study with interviews and discussions with native speakers. I had three aims in mind: First, to find out whether representative examples are actually acceptable for speakers. Second, to examine which order of adjectival relative clause and \is{numeral}numeral they generally prefer and what role additional constituents play. Third, to find out whether native speakers have intuitions on factors that might influence their preference regarding the order of relative clauses and numerals.\footnote{These interviews were conducted over Zoom with three speakers in March 2023. I am incredibly grateful to my consultants.} \is{complexity}

The first thing to note is that the speakers shared the initial impression I had from the corpus analysis. The fact that only a few results are attested is because combinations of inflected adjectives and numerals within the same DP are not entirely natural. In terms of the analysis, as mentioned above, pairs of simple relative clauses containing an inflected adjective but no additional constituents in combination with a numeral were by all speakers judged equally fine in both orders, both in terms of grammaticality and acceptability. No speaker had any particular preference for one order over the other. If the adjectival relative clause, however, was presented with an additional constituent, for example a degree modifier such as \textit{totemo} `very’, they expressed a clearly higher preference for the position of the relative clause to the left of the numeral. The relevant example is repeated from above (=\ref{dorama}).

\ex. \ag. [totemo kanashi-k-at-ta [mittsu-no dorama]]\\
very sad-\textsc{pred-cop-pst} three-\textsc{mod} soap.opera\\
\bg. \#[mittsu-no [totemo kanashi-k-at-ta dorama]]\\
three-\textsc{mod} very sad-\textsc{pred-cop-pst} soap.opera\\
`three soap operas that were very sad'

When asked for potential reasons for their choice, all consulted speakers noted that for them the sentence becomes easier to parse and that the length of the relative clause influenced their choice to place it DP-initially. One speaker specified that the order in which the numeral is to the left of the relative clause is not only odd, but stylistically awkward with longer relative clauses because the numeral seems too far detached from the noun.

Especially in relative clauses containing subjects, the order \textit{adjectival relative clause} $\succ$ \textit{numeral} was clearly preferred. Some speakers even noted that for them, the reversed order is decreased not only in preference but in grammaticality. An illustrative example is \ref{questions,many}.

\ex. \label{questions,many} \ag. [toiawase-ga ō-k-at-ta [yattsu-no shitsumon]]\\
inquiry-\textsc{nom} many-\textsc{pred-cop-pst} eight-\textsc{mod} question\\
\bg. \#[yattsu-no [toiawase-ga ō-k-at-ta shitsumon]]\\
eight-\textsc{mod} inquiry-\textsc{nom} many-\textsc{pred-cop-pst} question\\
`the eight questions with many inquiries’ (Lit. `the eight questions where the questions were many’)

A reason for this might be the desire to avoid \is{attachment ambiguity}attachment ambiguities. Such ambiguities can arise when the \is{numeral}numeral can logically refer to not only the noun to be modified but also the subject contained in the relative clause. Syntactically, the numeral forms a constituent with the nominal subject within the relative clause, not with the modified head noun. This is indicated below with brackets and in the English translation. \is{complexity}

\exg. [[yattsu-no toiawase-ga] ō-k-at-ta shitsumon]\\
eight-\textsc{mod} inquiry-\textsc{nom} many-\textsc{pred-cop-pst} question\\
`question(s) in which eight inquiries were much’ (= showed up a lot)

\largerpage
In other words, speakers might interpret DPs in which a numeral precedes a relative clause with a nominal subject in such a way that the numeral is part of the relative clause itself and modifies the subject, rather than the head noun of the DP. \textit{Yattsu-no} `eight’ can be understood as a modifier of the noun \textit{toiawase} `inquiry’, or of the noun \textit{shitsumon} `question’. To avoid such ambiguities, speakers might prefer the reversed order. Another mechanism to avoid attachment ambiguities is to insert a comma, or a pause in spoken language as an intonational break.\footnote{I thank Jiro Inaba (p.c.) for this observation.}

Importantly, however, relative clauses whose subject cannot be referred to by the relevant numeral,\footnote{For example \textit{futari-no} `two’ can only refer to human entities, not to dresses etc.} hence where no attachment ambiguities are expected, were nevertheless preferred to the left of the numeral by all speakers.

This discussion shows that \is{complexity}syntax-external factors have to be considered when it comes to the order of these elements. This is similar to the effect of \is{scope}scope on the order of direct modifiers discussed in Section \ref{optionality} (and see \citealt{ramchand:2014:deriving}). However, the results are not as clear, as also evidenced by the subtle judgments. We cannot determine a certain threshold of \textit{complexity} which bans a relative clause from a lower position. Nevertheless, the observation holds that more complex relative clauses are definitely preferred to the left of numerals. Possible syntactic correlations of this fact will be reviewed in the following.

\subsection{Discussion} \label{RC discussion}

The fact that (restrictive) relative clauses are possible to the left and to the right of \is{numeral}numerals falls out from the assumption that there are several different merging positions above and below FP\textsubscript{Num}. These are not reserved for relative clauses of different sizes; as was argued, Japanese relative clauses are all CPs. This is thus crucially different from languages such as English where this is the main motivation for different projections above and below FP\textsubscript{Num} for the respective relative clauses.\footnote{For \is{non-restrictive relative clause}\textit{non-restrictive} relative clauses, which must merge DP-externally, see Section \ref{RCs above dem}.}

It is tempting to reduce this phenomenon, again, to the lack of the relevant categories \is{Units of Language}(\textit{Units of Language}, \citealt{wiltschko:2014:universal}) in the linguistic inventory of Japanese in order to overtly realize different projections for \textit{full} and \textit{reduced} relative clauses. However, since Japanese does not overtly manifest the different types of relative clauses in the first place, it is hard to justify this conclusion.

In any case, these findings are highly relevant from a cross-linguistic point of view, especially for research on other prenominal-RC languages, as well as such languages which also do not mark a difference between the two categories of relative clause morphologically or syntactically.

To summarize, the two relative clauses in \ref{pati3} and \ref{pati4} occupy two different positions within the indirect domain and this is shown in \figref{Relative clause vs. numeral}.\footnote{Recall that numerals are selected via \textsc{mod} as well, but this is simplified here.}

% \begin{figure}
% \begin{subfigure}{\textwidth}
% \centering
% \begin{tikzpicture}[baseline=0pt, remember picture]
% \tikzset{level distance=25pt, sibling distance=5pt}
% \Tree [.DP\textsubscript{internal} {} [.D′\textsubscript{internal} [.FP\textsubscript{indirect} [.CP \edge[roof]; {pātī-de at-ta} ] [.F′\textsubscript{indirect} [.FP\textsubscript{Num} [.NumP \edge[roof]; {futari-no} ] [.F′\textsubscript{Num} [.FP\textsubscript{indirect} {} [.F′\textsubscript{indirect} [.... [.NP \edge[roof]; {josei} ] {} ] F ] ] F ] ] F ] ] D ] ]
% \end{tikzpicture}
% \caption{Structure: Relative clause $\succ$ numeral (= \ref{pati3})}
% \label{Structure RC Num}
% \end{subfigure}
%
% \begin{subfigure}{\textwidth}
% \centering
% \begin{tikzpicture}[baseline=0pt, remember picture]
% \tikzset{level distance=25pt, sibling distance=5pt}
% \Tree [.DP\textsubscript{internal} {} [.D′\textsubscript{internal} [.FP\textsubscript{indirect} {} [.F′\textsubscript{indirect} [.FP\textsubscript{Num} [.NumP \edge[roof]; {futari-no} ] [.F′\textsubscript{Num} [.FP\textsubscript{indirect} [.CP \edge[roof]; {pātī-de at-ta} ] [.F′\textsubscript{indirect} [.... [.NP \edge[roof]; {josei} ] {} ] F ] ] F ] ] F ] ] D ] ]
% \end{tikzpicture}
% \caption{Structure: Numeral $\succ$ relative clause (= \ref{pati4})}
% \label{Structure Num Rc}
% \end{subfigure}
%
% \caption{Two possible orders of relative clause and numeral}
% \label{Relative clause vs. numeral}
%
% \end{figure}

\begin{figure}
\begin{subfigure}{.45\textwidth}
\begin{forest}
for tree={
  l sep=25pt,
  s sep=5pt,
}
[DP\textsubscript{internal}
  [{}]
  [D′\textsubscript{internal}
    [FP\textsubscript{indirect}
      [CP
        [pātī-de at-ta, roof]
      ]
      [F′\textsubscript{indirect}
        [FP\textsubscript{Num}
          [NumP
            [futari-no, roof]
          ]
          [F′\textsubscript{Num}
            [FP\textsubscript{indirect}
              [{}]
              [F′\textsubscript{indirect}
                [....
                  [NP
                    [josei, roof]
                  ]
                  [{}]
                ]
                [F]
              ]
            ]
            [F]
          ]
        ]
        [F]
      ]
    ]
    [D]
  ]
]
\end{forest}
\caption{Structure: Relative clause $\succ$ numeral\newline (= \ref{pati3})}
\label{Structure RC Num}
\end{subfigure}
\hfill
\begin{subfigure}{.45\textwidth}
\begin{forest}
for tree={
  l sep=25pt,
  s sep=5pt,
}
[DP\textsubscript{internal}
  [{}]
  [D′\textsubscript{internal}
    [FP\textsubscript{indirect}
      [{}]
      [F′\textsubscript{indirect}
        [FP\textsubscript{Num}
          [NumP
            [futari-no, roof]
          ]
          [F′\textsubscript{Num}
            [FP\textsubscript{indirect}
              [CP
                [pātī-de at-ta, roof]
              ]
              [F′\textsubscript{indirect}
                [....
                  [NP
                    [josei, roof]
                  ]
                  [{}]
                ]
                [F]
              ]
            ]
            [F]
          ]
        ]
        [F]
      ]
    ]
    [D]
  ]
]
\end{forest}
\caption{Structure: Numeral $\succ$ relative clause\newline (= \ref{pati4})}
\label{Structure Num Rc}
\end{subfigure}

\caption{Two possible orders of relative clause and numeral}
\label{Relative clause vs. numeral}

\end{figure}

The question is, what the underlying reason for the two ordering possibilities? Can relative clauses simply be base-generated in different positions? Or is one position a derived position? Based on the fact that \is{complexity}complexity is a driving factor that leads speakers to prefer the higher position in a systematic fashion, I opt for the second option: Concretely, relative clauses are base-generated below FP\textsubscript{Num} but can optionally move across this projection individually (cf. \citealt{laenzlinger:2011:elements, kim:2019:syntax}). Furthermore, this movement option is available to all relative clauses, but it correlates with higher complexity. In other words, the more complex a relative clause is, the more prone it is to move up. This means that the derivation of the relative clause in \ref{pati3} is as in Figure \ref{movement Rc above num}.

\begin{figure}
\centering
 \begin{tikzpicture}[baseline=0pt, remember picture]
\tikzset{level distance=25pt, sibling distance=5pt}
\Tree [.DP\textsubscript{internal} {} [.D′\textsubscript{internal} [.FP\textsubscript{indirect} [. CP \edge[roof]; \node (cp1) {pātī-de at-ta}; ] [.F′\textsubscript{indirect} [.FP\textsubscript{Num} [.NumP \edge[roof]; {futari-no} ] [.F′\textsubscript{Num} [.FP\textsubscript{indirect} \node (cp2) { $<$CP$>$}; [.F′\textsubscript{indirect} [.... [.NP \edge[roof]; {josei} ] {} ] F ] ] F ] ] F ] ] D ] ]
\draw [color=blue][->] (cp2) |- +(0,-2) -| (cp1);
\end{tikzpicture}
\caption{Movement of relative clause above NumP}
\label{movement Rc above num}
\end{figure}

\largerpage[-1]
The seemingly free order of relative clauses and \is{numeral}numerals with potentially only complexity as a driving factor for movement is not straightforwardly explainable in a purely cartographic account, or any purely syntactic account. This holds especially if the only possible movement operations involve (phrasal) movement of the lexical head -- either alone or in one of the different \is{pied-piping}pied-piping fashions \citep{cinque:2017:micro, cinque:2021:carto, cinque:2023:linearization, cinque:2025:superiority}. However, it becomes clear that the lexical head is not involved in these derivations, hence relative clauses must move individually. Furthermore, adopting the standard assumption that movement is only possible to the left, the lower position must be where relative clauses are base-generated. As a result of left bound movement, relative clauses reach a position above FP\textsubscript{Num} and they do this for syntax-external reasons, i.e. \is{complexity}complexity. This accounts for the systematic preference for more complex relative clauses in a higher position.\footnote{A reviewer correctly points out that complexity should not be visible in syntax. However, there is, to my knowledge, no account to better capture the behavior explained here. Note that assuming different areas for base-generation is not an appealing option either. Even if more than one position were accessible for base-generation, the syntax would face the same problem, namely to decide when to merge relative clauses and to decide to merge more complex ones later. It is, at this point, unclear to me how to model this theoretically. \label{footnote complexity}}

This proposal matches earlier proposals for movement of relative clauses (cf. \citealt{kim:2019:syntax} on individual movement of \is{Korean} Korean relative clauses depending on their involvement in the discourse) but also of other modifiers (focus movement, \citealt{giusti:1996:is, dimitrova:1998:fragments, ihsane:2001:specific, aboh:2004:topic, aboh:2006:when, svenonius:2008:position}).

Having established that movement is triggered by syntax-external factors, the movement process observable with \is{complexity}complex modifiers seems to be some sort of rescue operation. Looking outside of Cartography, and generally syntax-based accounts, the fact that longer elements are preferred before shorter elements is in conformity with the so-called \textit{long-before-short principle} that predicts for SOV languages that longer elements are best to the left of shorter elements. This was, for instance, tested empirically for \is{Persian} Persian, another SOV language, in \citet{faghiri:2020:word}. This principle is the direct opposite of the well-established \textit{Gesetz der wachsenden Glieder} `Law of increasing constituents' of \citet{behaghel:1909:glieder}, which is not designed for (strictly head-final) SOV languages. \is{long-before-short principle}

Another interesting parallel concerns compound formation in Japanese in relation to prosodic \is{compound accentuation}accentuation.\footnote{I thank an anonymous reviewer for bringing this parallel to my attention.} Compounds consisting of three elements can be left- and right-branching. In left-branching compounds, the compound only bears one accent and the three elements are phrased as one unit. In right-branching compounds, on the other hand, there is an additional accent after the first element, causing two accentuation phrases (\citealt[12–14]{kubozono:2005:rendaku}; cf. \citealt{nishiyama:2017:phrasal}). This is shown in examples \ref{example left-branching} and \ref{example right-branching}, adapted from \citet[12–14]{kubozono:2005:rendaku}. Accents mark accented morae, square brackets indicate morphological units, and curly brackets indicate accentual phrases.

\ex. Left-branching compounds \\\label{example left-branching}
$\lbrack$ Do’itsu + bu'ngaku$\rbrack$ + kyōkai $\rightarrow$ \{doitsu-bungaku-kyo'okai\}  \\
`Germany + literature + association' $\rightarrow$ `German Association of literature'

\newpage
\ex. Right-branching compounds \\\label{example right-branching}
Do’itsu + $\lbrack$ bu'ngaku + kyōkai$\rbrack$ $\rightarrow$ \{do'itsu\} \{bungaku-kyo'okai\}  \\
`Germany + literature + association' $\rightarrow$ `Association of German literature'

This means that right-branching compounds are the marked case since they violate one of the accentuation rules of Japanese \citep[153]{nishiyama:2017:phrasal}. This is translatable into the long-before-short principle in a straightforward manner. In left-branching compounds, the first two elements form a morphological compound first to which the third is added, meaning a longer sequence (two elements) precedes a shorter sequence (one element). In right-branching compounds, the second and third element have formed a morphological compound to which the first element is added. Therefore, a shorter sequence (one element) precedes a longer sequence (two elements). In accordance with the \is{long-before-short principle}long-before-short principle, the first option should be favored, and this is what happens. Right-branching \is{compound accentuation}compounds are more difficult to parse and hence involve marked accentuation. 

To summarize, relative clauses in Japanese originate in a lower area within the indirect domain and can optionally move upwards to a position above \is{numeral}numerals. One driving factor for this movement operation is \is{complexity}complexity, a syntax-external factor. These observations fall out nicely under the general long-before-short principle observable in SOV langauges. It is desirable to find more instantiations of this principle in Japanese and other SOV languages, especially pertaining to the DP, and/or the factor complexity.

\subsection{Relative clauses above demonstratives} \label{RCs above dem}

This remaining section addresses the distribution of \is{non-restrictive relative clause}\textit{non-restrictive} relative clauses, which are confined to a position above demonstratives, one of the factors that point towards a DP-external domain. Recall the following example from \citet{ishizuka:2008:RC} which shows that those types of relative clause cannot occur below demonstratives but must occur above them.

\exg. *Ito-san-ni-wa musuko-ga hitori i-ru. [Sono [sakunen isha-ni nat-ta musuko-ga]] kekkon-shi-ta.\\
Ito-Mr.-\textsc{loc-top} son-\textsc{nom} one be-\textsc{prs} that last.year doctor-\textsc{loc} become-\textsc{pst} son-\textsc{nom} marriage-do-\textsc{pst}\\
(`Mr. Ito has one son. That son, who happened to become a doctor last year, got married.’) \\ \null \hfill *demonstrative $\succ$ non-restrictive relative clause \citep[5]{ishizuka:2008:RC}

\citet{ishizuka:2008:RC} only considers verbal relative clauses, but the same facts hold for adjectival relative clauses. If they appear higher than demonstratives they can be interpreted restrictively or non-restrictively. A sentence such as \ref{imamade}, where a \textit{na}-adjective precedes the demonstrative and is inflected for past, could be uttered in a context in which Mr./Mrs. Itō has one son and offers additional information about him (non-restrictive), as well as in a context in which s/he has multiple sons and wants to identify that son who got married (namely the shy one).

\exg. [Imamade totemo uchiki-d-at-ta [sono musuko-ga]] kekkon-shi-ta.\\
up.until.now very shy-\textsc{pred-cop-pst} that son-\textsc{nom} marriage-do-\textsc{pst}\\
1. $\approx$ `That son who was very shy up until now got married.’ \\
2. $\approx$ `That son, who happened to be very shy up until now, got married.’ \label{imamade}

On the other hand, it is also impossible for \is{non-restrictive relative clause}non-restrictive adjectival relative clauses to appear below demonstratives, shown in the adaptation of Ishizuka's example in \ref{ishizukaadj} where a non-restrictive interpretation is forced.\footnote{Again, speakers' judgment is subtle and some might find such examples acceptable, especially in colloquial speech. However, it seems that when choosing between two options in such a context, speakers might always go for the option in which the non-restrictive relative clause precedes the demonstrative. Future research will shed light on this phenomenon.}

\exg. *Ito-san-ni-wa musuko-ga hitori i-ru. [Sono [totemo uchiki-d-at-ta musuko-ga]] kekkon-shi-ta.\\
Ito-Mr.-\textsc{loc-top} son-\textsc{nom} one be-\textsc{prs} that very shy-\textsc{pred-cop-pst} son-\textsc{nom} marriage-do-\textsc{pst}\\
(`Mr. Ito has one son. That son, who happened to be very shy, got married.’) \label{ishizukaadj}

From these facts, I draw the conclusion that non-restrictive relative clauses merge in the DP-external domain -- the \is{left periphery}left periphery -- and hence are outside the scope of D, predicting their non-restrictive interpretation. This is shown in \figref{non-restrictive rc - dem}.

\begin{figure}
\centering
\begin{tikzpicture}[baseline=0pt, remember picture]
\tikzset{level distance=25pt, sibling distance=5pt}
\Tree [.DP\textsubscript{external} {} [.D′\textsubscript{external} [.FP [.CP\textsubscript{non-res} \edge[roof]; {isha-ni natta} ] [.F′ [.DP\textsubscript{internal} sono [.D′\textsubscript{internal} [.... [.NP \edge[roof]; musuko ] {} ] D ] ] F ] ] D ] ]
\end{tikzpicture}
\caption{Base-generated position: Non-restrictive relative clause $>$ demonstrative}
\label{non-restrictive rc - dem}
\end{figure}

\largerpage[-2.5]
On the other hand, it was shown that restrictive relative clauses merge in the indirect domain below FP\textsubscript{Num}. They hence have not only the option to move to a functional projection above \is{numeral}numerals, but also to the left periphery, accounting for restrictive relative clauses to the left of demonstratives. Again, they do this independently as depicted in \figref{restrictive rc - dem}.\footnote{\citet{ishizuka:2008:RC} proposes an analysis in which two D-projections are assumed in Japanese, a lower one, DemP, where the first part of the demonstrative, namely \textit{ko}- `speaker-proximate', \textit{so}- `hearer-proximate' and \textit{a}- `distant', is generated, and a higher one, DP, headed by \textsc{no}. Having two projections available gives two possible landing sites for relative clauses, which originate very low in her system \citep{kayne:1994:anti}. In her derivation, movement of the deictic element to Spec, DP yields the correct lexical item. Then the NP moves out of the relative clause to Spec, CP, followed by remnant movement of the IP to Spec, DemP. In this position, it is in the scope of D, therefore can only receive a restrictive interpretation. Additionally, in a further step, the remnant can optionally raise from Spec, DemP to Spec, DP leaving the c-command domain of D and gaining access to a non-restrictive interpretation (see \citealt[8–9]{ishizuka:2008:RC}). One problem is that the two positions, DP and DemP, lack motivation beyond the derivation proposed here. \label{footnote ishizuka}}

\begin{figure}
\centering
\begin{tikzpicture}[baseline=0pt, remember picture]
\tikzset{level distance=25pt, sibling distance=5pt}
\Tree [.DP\textsubscript{external} {} [.D′\textsubscript{external} [.FP [.CP\textsubscript{res} \edge[roof]; \node (isha1) {isha-ni natta}; ] [.F′ [.DP\textsubscript{internal} sono [.D′\textsubscript{internal} [.FP\textsubscript{indirect} \node (isha2) {$<$CP$>$}; [.F′\textsubscript{indirect} [.... [.NP \edge[roof]; musuko ] {} ] F ] ] D ] ] F ] ] D ] ]
\draw [color=blue][->] (isha2) |- +(0,-2) -| (isha1);
\end{tikzpicture}
\caption{Movement: Restrictive relative clause $\succ$ demonstrative}
\label{restrictive rc - dem}
\end{figure}

Importantly, the two types of relative clauses, restrictive and \is{non-restrictive relative clause}non-restrictive, occupy different positions, as non-restrictive relative clauses originate DP-externally but restrictive ones move from within the DP-internal domain. The analysis presented here is in line with the analysis of \citet[223]{cinque:2020:relcl} and \citet{cinque:2008:two}, who reduces the differences in interpretation to two different merging positions as well: one outside (= above) demonstratives for non-restrictive and one inside (= below) demonstratives for restrictive relative clauses. As in our system, the additional pattern of Japanese is explained via optional raising of restrictive relative clauses across DP\textsubscript{internal} to a position below non-restrictive relative clauses \citep[116]{cinque:2008:two}. Also recall that \citet{kim:2019:syntax} proposes the projection SpplP (\textit{SupplementaryP}) for non-restrictive, but importantly supplementary, relative clauses in the \is{left periphery}left periphery of Korean, and this is essentially the counterpart to what I simply give as FP\textsubscript{indirect}.\footnote{In this respect, note that Japanese offers a parametric difference from \is{Mandarin}Mandarin, where relative clauses in the lower position have access to a restrictive as well as a non-restrictive interpretation, but in the high position are usually interpreted restrictively, although open to a non-restrictive interpretation (\citealt[235–236]{cinque:2020:relcl}; \citealt{hsu:2017:alternatives}). Assuming the position in the left periphery is equally reached via independent movement of the relative clause means that the relation of movement and interpretation is different from Japanese and should be investigated. Also recall the German data given in Footnote \ref{footnote German RC} in Section \ref{RC LP}.}

%Finally, Japanese provides evidence that demonstratives can move from the DP-internal to the \is{left periphery}DP-external domain. This gives rise to the patterns in which \is{non-restrictive relative clause}non-restrictive relative clauses appear to the right of demonstratives.

\newpage

A remaining question to be answered is whether there is a driving force for the movement of restrictive relative clauses to the \is{left periphery}left periphery.\footnote{Another issue to investigate is the nature of non-restrictive relative clauses. \citet{cinque:2008:two} proposes a difference between \textit{integrated} and \textit{non-integrated} non-restrictive relative clauses and points out that languages such as Italian and other Romance varieties possess both types, but that there might be languages in which only one type exists, among them presumably Japanese \citep[121–124]{cinque:2008:two}. This also concerns the question as to whether non-restrictive relative clauses are base-generated higher than \is{quantificational modifier}quantificational modifiers in the left periphery or not. Finally, I have assumed that the projections for non-restrictive and restrictive relative clauses in the left periphery are different. This must be the case as stacking of the two types of relative clauses is generally allowed \citep{cinque:2008:two, ishizuka:2008:RC}. However, there are different opinions on the ordering principles among non-restrictive and restrictive relative clauses and while \citet{cinque:2008:two} proposes a hierarchy \textit{non-restrictive relative clause} $>$ \textit{restrictive relative clause}, his judgments are not shared by \citet{ishizuka:2008:RC}.} I investigated whether a similar effect in terms of \is{complexity}complexity can be observed; in other words, whether complexity can also be responsible for an increased preference for relative clauses above demonstratives. As it turns out, this is not the case, exemplified in \ref{adjective order lp} with an adjectival relative clause containing a subject. While this relative clause would be strongly preferred above \is{numeral}numerals within the indirect domain, there does not seem to be any preference for a higher position above a demonstrative. This suggests that, although complexity encourages movement above FP\textsubscript{Num}, it does not necessarily do so above the internal D-head. This means then that the driving factor behind the movement of restrictive relative clauses to a position above demonstratives remains undetermined.

\ex. \label{adjective order lp} Context: Two friends went to a restaurant yesterday and had several dishes there. Speaker A asks Speaker B about one specific dish.
 \ag. Kinō it-ta resutoran-no, [ano [sūpu-ga sugoi kara-k-at-ta ryōri-wa]] nan d-a kke?\\
yesterday go-\textsc{pst} restaurant-\textsc{mod} that soup-\textsc{nom} very spicy-\textsc{pred-cop-pst} dish-\textsc{top} what \textsc{pred-cop.pst} \textsc{q} \\
\bg. Kinō it-ta resutoran-no, [sūpu-ga sugoi kara-k-at-ta [ano ryōri-wa]] nan d-a kke?\\
yesterday go-\textsc{pst} restaurant-\textsc{mod} soup-\textsc{nom} very spicy-\textsc{pred-cop-pst} that dish-\textsc{top} what \textsc{pred-cop.pst} \textsc{q}\\
`What was that dish in which the soup was very spicy at the restaurant we went to yesterday?’

To summarize, \is{non-restrictive relative clause}non-restrictive relative clauses are base-generated DP-externally. Restrictive relative clauses merge in the indirect domain and can optionally move to the left periphery.

\section{Chapter summary and closing remarks}

Although indirect modification has been argued to be the only type of modification available in Japanese, many questions have been left open and were addressed in the preceding pages. The following conclusions can be drawn.

\ex. \a. The complex elements following Japanese modifiers are combinations of the predicative copula, the dummy copula and a tense marker.
\b. Relative clauses in Japanese must be at least of size IP, but we cannot determine whether they are bigger than that or not. I have assumed a CP status for several reasons, among them the possibility of focused elements and the better compatibility with the operator-movement approach adopted in this book.
\c. The difference between full and reduced relative clauses does not exist.
\d. The indirect domain is comprised of several functional projections reserved for indirect modifiers and a projection reserved for numerals. The projections reserved for indirect modifiers are accordingly not reserved for modifiers of a specific type (full or reduced) or specific syntactic size (IP or CP).
\d. Indirect modifiers are base-generated below FP\textsubscript{Num} but can independently move up to a higher position. One driving factor influencing this movement operation is complexity.
\e. Non-restrictive relative clauses merge DP-externally.
\e. Restrictive relative clauses can optionally move to the left periphery.

These findings are similar to the findings in the direct domain and compatible with previous accounts in the cartographic research field. On the other hand, they emphasize that a certain degree of tolerance is needed when it comes to movement and ordering. Abandoning syntactic size and category as criteria that influence the position of indirect modifiers in Japanese leads to the conclusion that the functional projections inside the indirect domain only correlate to these parameters when they are overtly realized in a given language. In, for example, English or German, relative clauses of different sizes can be distinguished in the surface structure and therefore dedicated functional projections are justified. However, this is not the case in Japanese. From a perspective of language uniformity, I suggest that the structure of the indirect domain presented by Japanese, and similar languages, is the underlying structure. Languages then show parametric, language-specific differences as to whether they realize different functional projections dedicated to different types of indirect modifiers or not.

